\chapter{Circuits on a finite set of parties}
\label{ch:party}

This chapter states the calculus of circuits run by a finite set of parties,
following \cite[Sections~5--7]{OpenAI2026PolynomialPEPS}. Every Hilbert space
is finite dimensional, and at every time the memory is the tensor product of
the current spaces of the parties. A circuit on parties is built from local
operations, which act at one party, and gates, which act on the parties of a
cell; it is read as a circuit with one gate or one local family per layer on
the geometry of all finite sets of parties, and its operator is the Kraus
operator of its single branch. Pair sources prepare a normalized vector on
fresh registers at two parties, and pair effects contract such registers. The
chapter first sets up this syntax and the operations on it, then proves the
elimination of pair effects of \cite[Lemma~5.1]{OpenAI2026PolynomialPEPS}, and
closes with ownership sheets, their births, deaths and exchanges
\cite[Lemmas~6.5 and~6.6]{OpenAI2026PolynomialPEPS}, and stacks of sheets as
used in \cite[Proposition~7.1]{OpenAI2026PolynomialPEPS}.

\section{The syntax of circuits on parties}
\label{sec:party_syntax}

Throughout this chapter $P$ is a finite set of parties, a system
$X=(X_p)_{p\in P}$ assigns a finite basis set to every party,
$\mathcal H_X=\bigotimes_p\C^{X_p}$, and $X_P$ and $X_D$ are the configurations of
all parties and of a set $D$ of parties, as in Chapter~\ref{ch:sites}. This
section defines the circuits run by the parties
\cite[Section~5.1]{OpenAI2026PolynomialPEPS}, reads them as circuits on the
geometry of all finite sets of parties, and counts their gates. It then
introduces the gates of the compression argument: pair sources and pair effects,
their monomials and weighted sums, and words of pair sources.

\subsection{Party circuits and their operators}
\label{sec:party_syntax_circuits}

\begin{definition}[Circuits on a finite set of parties]
    \label{def:party_circuit}
    \lean{QuantumCircuit.Party.PartyCircuit, QuantumCircuit.Party.PartyCircuit.toCircuit,
        QuantumCircuit.Party.PartyCircuit.op, QuantumCircuit.Party.PartyCircuit.IsAtom}
    \leanok
    \uses{def:circuits_gate_family, def:circuits_subsets_geometry, def:circuits_gate_at,
        def:circuits_of_cells, thm:circuits_composition, def:circuits_single,
        thm:circuits_single_op}
    Let $\mathcal G$ be a gate family on $P$. A party circuit from a system $X$ to a
    system $Y$ with gates in $\mathcal G$ is built by the following rules
    \cite[Section~5.1]{OpenAI2026PolynomialPEPS}:
    \begin{enumerate}
        \item the empty circuit from $X$ to $X$;
        \item one gate $g$ of $\mathcal G$ on a set $D\subseteq P$, its cell, from $X$
              to a system $Y$ with $Y_s=X_s$ for every $s\notin D$; the other parties
              are spectators, and the cell may be empty;
        \item one local operation at a party $p$: a matrix $A$ from $\C^{X_p}$ to
              $\C^{Y_p}$, where $Y_s=X_s$ for every $s\neq p$;
        \item local operations at every party at once: matrices $A_p$ from
              $\C^{X_p}$ to $\C^{Y_p}$ for all $p\in P$;
        \item the composition $c'\circ c$ of a party circuit $c$ from $X$ to $Y$ and a
              party circuit $c'$ from $Y$ to $Z$, first $c$ and then $c'$.
    \end{enumerate}
    The party circuits given by the first four rules are the party circuits of one
    step. The circuit $\widehat c$ of a party circuit $c$ is the circuit on the
    geometry of all finite sets of parties defined rule by rule: the empty circuit;
    the circuit of the layer of $g$ on $D$; the circuit of the layer of $A$ at $p$;
    the circuit of the layer with no cell and the one-member family $\{A_p\}$ at every
    party $p$; and the composition of $\widehat c$ and $\widehat{c'}$. When every gate
    of $\mathcal G$ has one outcome, $\widehat c$ has one branch
    (Theorem~\ref{thm:party_circuit_op}), and for an interpretation of $\mathcal G$ the
    operator $T_c$ of $c$, a matrix from $\mathcal H_X$ to $\mathcal H_Y$, is the
    Kraus operator of that branch.
\end{definition}

\begin{theorem}[The circuit and the operator of a party circuit]
    \label{thm:party_circuit_op}
    \lean{QuantumCircuit.Party.PartyCircuit.isUnbranched_toCircuit,
        QuantumCircuit.Party.PartyCircuit.isSequential_toCircuit,
        QuantumCircuit.Party.PartyCircuit.isSingle_toCircuit,
        QuantumCircuit.Party.PartyCircuit.op_done, QuantumCircuit.Party.PartyCircuit.op_comp,
        QuantumCircuit.Party.PartyCircuit.op_gate, QuantumCircuit.Party.PartyCircuit.op_localOp,
        QuantumCircuit.Party.PartyCircuit.op_localOps,
        QuantumCircuit.Party.PartyCircuit.op_induction}
    \leanok
    \uses{def:party_circuit, def:circuits_local, def:circuits_gate_at, def:circuits_single,
        def:sites_place_between, def:elementary_finite_tensor_product}
    Let $c$ be a party circuit from $X$ to $Y$. Its circuit $\widehat c$ is unbranched
    and sequential: every layer has at most one cell, and a layer with a cell is the
    identity at the parties outside it. If every gate has one outcome, $\widehat c$ is
    single. For an interpretation of a gate family with one outcome per gate, the
    operators are, rule by rule,
    \begin{align}
        T_{\mathrm{empty}} & =1,                           &
        T_{c'\circ c}      & =T_{c'}T_c,            \notag   \\
        T_g                & =M_g\otimes1,                 &
        T_A                & =A\otimes1,                   &
        T_{(A_p)_p}        & =\bigotimes_pA_p,    \notag
    \end{align}
    where $M_g$ is the matrix of $g$ at its outcome, and $M_g\otimes1$ and $A\otimes1$
    are the placements between $X$ and $Y$ on $D$ and on $\{p\}$, with $A$ read on the
    configurations of $\{p\}$. Consequently, let a property of matrices from
    $\mathcal H_X$ to $\mathcal H_Y$, for all systems $Y$, hold for the identity of
    $\mathcal H_X$ and pass from $M$ to $T_sM$ for every party circuit $s$ of one step.
    Then it holds for $T_c$ for every party circuit $c$ from $X$.
\end{theorem}
\begin{proof}\leanok
    \uses{lem:circuits_duality, lem:circuits_local_comp, def:circuits_gate_at,
        def:circuits_single, thm:circuits_single_op, lem:circuits_gate_at,
        thm:circuits_of_cells}
    By induction on $c$. A gate, a local operation at one party and local operations
    at every party give the circuit of one layer, which is unbranched
    (Lemma~\ref{lem:circuits_duality}). Each of these layers is one step, and the
    circuit of a layer of one step is sequential
    (Definition~\ref{def:circuits_gate_at}); the layer of a gate with one outcome, and
    a layer without cells whose onsite families have one member, give single circuits
    (Definition~\ref{def:circuits_single}). The three properties pass to compositions
    (Lemma~\ref{lem:circuits_local_comp} and Definitions~\ref{def:circuits_gate_at}
    and~\ref{def:circuits_single}). The operator of a single circuit of one layer is a
    Kraus operator of the layer (Theorem~\ref{thm:circuits_single_op}); for the layer
    of a gate it is $M_g\otimes1$ at the unique outcome, for the layer of $A$ at $p$ it
    is $A\otimes1$ (Lemma~\ref{lem:circuits_gate_at}), and for the layer without cells
    it is $\bigotimes_pA_p$ (Theorem~\ref{thm:circuits_of_cells}). The operator of a
    composition of single circuits is $T_{c'}T_c$ by
    Theorem~\ref{thm:circuits_single_op}. For the last statement, show by induction on
    a circuit $c$ from $Y$ to $Z$ that the property passes from $M$ to $T_cM$: a circuit
    of one step is the hypothesis, and $T_{c'\circ c}M=T_{c'}(T_cM)$. Apply this to
    $M=1$.
\end{proof}

\begin{definition}[Allowed and locally contractive party circuits]
    \label{def:party_allowed}
    \lean{QuantumCircuit.Party.PartyCircuit.IsAllowed,
        QuantumCircuit.Party.PartyCircuit.IsLocalContractive}
    \leanok
    \uses{def:party_circuit}
    Let $\mathcal G$ be a gate family with one outcome per gate, with an
    interpretation, and let $\mathcal A$ be a condition on the gates of $\mathcal G$. A
    party circuit is allowed for $\mathcal A$ when its local operations are
    contractions, $\lVert A\rVert\le1$ and $\lVert A_p\rVert\le1$, and every gate $g$
    satisfies $\mathcal A$ and has $\lVert M_g\rVert\le1$
    \cite[Section~5.1]{OpenAI2026PolynomialPEPS}. A party circuit is locally
    contractive when its local operations are contractions; its gates are not
    constrained.
\end{definition}

\begin{theorem}[An allowed party circuit is a contraction]
    \label{thm:party_op_contraction}
    \lean{QuantumCircuit.Party.PartyCircuit.norm_op_le_one,
        QuantumCircuit.Party.PartyCircuit.IsAllowed.forall_toCircuit,
        QuantumCircuit.Party.PartyCircuit.IsAllowed.isLocalContractive}
    \leanok
    \uses{def:party_allowed, def:circuits_local, def:kraus_family}
    Let $c$ be an allowed party circuit from $X$ to $Y$. Every gate and every onsite
    family of its circuit $\widehat c$ is trace-nonincreasing, $c$ is locally
    contractive, and $\lVert T_c\rVert\le1$
    \cite[Section~5.1]{OpenAI2026PolynomialPEPS}.
\end{theorem}
\begin{proof}\leanok
    \uses{thm:circuits_single_op, lem:circuits_gate_at, lem:circuits_local_gate_unitary,
        lem:kraus_family_channel, lem:circuits_local_comp, thm:party_circuit_op}
    By induction on $c$. A family with one member $V$ is trace-nonincreasing if and
    only if $\lVert V\rVert\le1$ (Theorem~\ref{thm:circuits_single_op}); in
    particular the family $\{1\}$ is, being unitary and hence trace-preserving
    (Lemmas~\ref{lem:circuits_local_gate_unitary} and~\ref{lem:kraus_family_channel}).
    Hence the layer of an allowed gate and the layer of a contraction at one party are
    trace-nonincreasing at every gate and onsite family
    (Lemma~\ref{lem:circuits_gate_at}), and so is the layer of contractions at every
    party. The property passes to compositions (Lemma~\ref{lem:circuits_local_comp}).
    Local contractivity is part of allowedness. Since $\widehat c$ is single
    (Theorem~\ref{thm:party_circuit_op}), Theorem~\ref{thm:circuits_single_op} gives
    $\lVert T_c\rVert\le1$.
\end{proof}

\subsection{Gates along a party circuit and their counts}
\label{sec:party_syntax_counts}

A gate whose cell consists of exactly one party is private. In the geometry of all
finite sets of parties a private cell costs $0$ and every other cell, the empty one
included, costs $1$ (Definition~\ref{def:circuits_subsets_geometry}).

\begin{definition}[Gate occurrences, gate counts and participation]
    \label{def:party_gate_counts}
    \lean{QuantumCircuit.GateOccurrence, QuantumCircuit.Party.PartyCircuit.gateLocations,
        QuantumCircuit.Party.PartyCircuit.participants, QuantumCircuit.Party.PartyCircuit.gates,
        QuantumCircuit.Party.PartyCircuit.gateCells, QuantumCircuit.Party.PartyCircuit.numGates,
        QuantumCircuit.Party.PartyCircuit.gateCount,
        QuantumCircuit.Party.PartyCircuit.participation,
        QuantumCircuit.Party.PartyCircuit.participationCount,
        QuantumCircuit.Party.PartyCircuit.ForallGate}
    \leanok
    \uses{def:party_circuit, def:circuits_subsets_geometry}
    Let $c$ be a party circuit. Its gate occurrences form a finite set, one for each
    gate, kept distinct even when two gates coincide: there is none for the empty
    circuit and for local operations, one for a gate, and the disjoint union of those
    of $c$ and $c'$ for $c'\circ c$. The participants of an occurrence are the parties
    of its cell \cite[Section~5.1]{OpenAI2026PolynomialPEPS}. The gates of $c$, each
    with its cell and its two systems, and their cells form lists in chronological
    order. Read along $c$, with the constructors of a composition adding:
    \begin{enumerate}
        \item the number of gates $n(c)$ counts every gate, private ones and those on
              the empty cell included;
        \item the gate count $\#(c)$ is the sum of the costs of the cells of the gates,
              the number of gates that are not private;
        \item the participation $\pi_p(c)$ of a party $p$ is the number of gates that
              are not private and whose cell contains $p$;
        \item the participation count $\nu_p(c)$ is the number of gates whose cell
              contains $p$, private ones included.
    \end{enumerate}
    A condition on gates holds at every gate of $c$ when every gate of $c$, with its
    cell and systems, satisfies it.
\end{definition}

\begin{theorem}[Relations among the counts]
    \label{thm:party_gate_count_relations}
    \lean{QuantumCircuit.Party.PartyCircuit.gateCells_eq_map_gates,
        QuantumCircuit.Party.PartyCircuit.length_gates,
        QuantumCircuit.Party.PartyCircuit.card_gateLocations,
        QuantumCircuit.Party.PartyCircuit.gateCount_le_numGates,
        QuantumCircuit.Party.PartyCircuit.gateCount_eq_numGates,
        QuantumCircuit.Party.PartyCircuit.participation_eq_participationCount,
        QuantumCircuit.Party.PartyCircuit.forall_participants_of_forallGate,
        QuantumCircuit.Party.PartyCircuit.natCard_participants_mem,
        QuantumCircuit.Party.PartyCircuit.gateCount_gate_empty}
    \leanok
    \uses{def:party_gate_counts}
    For every party circuit $c$, the list of cells is the list of the cells of the
    gates, and the length of the list of gates and the number of gate occurrences are
    both $n(c)$. Always $\#(c)\le n(c)$, and when no gate of $c$ is private,
    $\#(c)=n(c)$ and $\pi_p(c)=\nu_p(c)$ for every party $p$. A condition on cells that
    holds at every gate holds at the participants of every occurrence. The number of
    occurrences whose participants contain $p$ is $\nu_p(c)$
    \cite[Section~5.1]{OpenAI2026PolynomialPEPS}. One gate on the empty cell has gate
    count $1$.
\end{theorem}
\begin{proof}\leanok
    \uses{def:circuits_subsets_geometry}
    By induction on $c$, the constructors of a composition adding on both sides. For a
    gate on $D$ the cost of $D$ is at most $1$, it is $1$ when $|D|\neq1$, and it is
    positive exactly when $|D|\neq1$ (Definition~\ref{def:circuits_subsets_geometry}).
    The occurrences of a composition containing $p$ are the disjoint union of those of
    the two parts.
\end{proof}

\begin{theorem}[The counts of the circuit of a party circuit]
    \label{thm:party_counts_circuit}
    \lean{QuantumCircuit.Party.PartyCircuit.gateCount_toCircuit,
        QuantumCircuit.Party.PartyCircuit.participation_toCircuit}
    \leanok
    \uses{def:party_gate_counts, def:circuits_costs}
    If every gate has one outcome, then for every party circuit $c$ the gate count of
    its circuit $\widehat c$ is $\#(c)$, and the number of gates of $\widehat c$
    meeting a party $p$ is $\pi_p(c)$.
\end{theorem}
\begin{proof}\leanok
    \uses{thm:party_circuit_op, lem:circuits_gate_at, def:circuits_costs, lem:circuits_seq}
    By induction on $c$. The cost of the circuit of one layer is the cost of the
    layer. The layer of a gate on $D$ has gate count the cost of $D$ and meets $p$ once
    exactly when $p\in D$ and $D$ has positive cost, and the layer of a local operation
    costs nothing (Lemma~\ref{lem:circuits_gate_at}); a layer without cells costs
    nothing (Definition~\ref{def:circuits_costs}). For a composition, $\widehat c$ is
    unbranched with a branch (Theorem~\ref{thm:party_circuit_op}), so the costs of
    $\widehat c$ and $\widehat{c'}$ add (Lemma~\ref{lem:circuits_seq}).
\end{proof}

\subsection{Relabelling bases and replacing gates}
\label{sec:party_syntax_relabel}

\begin{definition}[Relabelling the bases and replacing the gates]
    \label{def:party_relabel}
    \lean{QuantumCircuit.Party.PartyCircuit.relabel,
        QuantumCircuit.Party.PartyCircuit.gates_relabel,
        QuantumCircuit.Party.PartyCircuit.gateCells_relabel,
        QuantumCircuit.Party.PartyCircuit.numGates_relabel,
        QuantumCircuit.Party.PartyCircuit.gateCount_relabel,
        QuantumCircuit.Party.PartyCircuit.mapGates,
        QuantumCircuit.Party.PartyCircuit.mapGatesLocationsEquiv}
    \leanok
    \uses{def:party_circuit, def:party_gate_counts, def:circuits_gate_morphism}
    For bijections $e_s\colon Y_s\to X_s$, the relabelling $R_e$ is the party circuit
    from $X$ to $Y$ of local operations at every party, at $s$ the matrix from
    $\C^{X_s}$ to $\C^{Y_s}$ whose entry at $(y,x)$ is $1$ if $e_s(y)=x$ and $0$
    otherwise. It has no gate and no cell, so $n(R_e)=\#(R_e)=0$. For a morphism
    $\varphi$ of gate families, the party circuit $\varphi_*c$ replaces every gate $g$
    of $c$ by $\varphi(g)$ and keeps every other rule; the gate occurrences of $c$ and
    of $\varphi_*c$ are identified rule by rule.
\end{definition}

\begin{theorem}[The operator of a relabelling]
    \label{thm:party_relabel_op}
    \lean{QuantumCircuit.Party.PartyCircuit.op_relabel,
        QuantumCircuit.Party.PartyCircuit.op_relabel_mul,
        QuantumCircuit.Party.PartyCircuit.op_relabel_symm,
        QuantumCircuit.Party.PartyCircuit.isAllowed_relabel}
    \leanok
    \uses{def:party_relabel, def:party_allowed}
    For a gate family with one outcome per gate and any interpretation, the operator
    of $R_e$ is the matrix of the bijection $y\mapsto(e_s(y_s))_s$ of configurations:
    $(T_{R_e})_{yx}=1$ if $e_s(y_s)=x_s$ for every $s$ and $0$ otherwise. Hence, for
    every matrix $M$ with rows indexed by $X_P$, the product $T_{R_e}M$ is $M$ with its
    rows reindexed, $(T_{R_e}M)_{y,w}=M_{(e_s(y_s))_s,w}$; no inverse appears. For
    bijections $e_s\colon X_s\to Y_s$, the operator of the relabelling along their
    inverses is the matrix of the inverse of the bijection of configurations. A
    relabelling is allowed for every condition on gates.
\end{theorem}
\begin{proof}\leanok
    \uses{thm:party_circuit_op, thm:sites_product_span, lem:matrix_l2_contractions}
    By Theorem~\ref{thm:party_circuit_op}, $T_{R_e}$ is the Kronecker product of the
    matrices of the $e_s$, whose entry at $(y,x)$ is $\prod_s[e_s(y_s)=x_s]$, the
    entry of the matrix of the bijection of configurations
    (Theorem~\ref{thm:sites_product_span}). Multiplying by the matrix of a bijection
    reindexes the rows. The matrix of a bijection is a contraction
    (Lemma~\ref{lem:matrix_l2_contractions}).
\end{proof}

\begin{theorem}[Replacing the gates]
    \label{thm:party_map_gates}
    \lean{QuantumCircuit.Party.PartyCircuit.op_mapGates,
        QuantumCircuit.Party.PartyCircuit.gates_mapGates,
        QuantumCircuit.Party.PartyCircuit.gateCells_mapGates,
        QuantumCircuit.Party.PartyCircuit.numGates_mapGates,
        QuantumCircuit.Party.PartyCircuit.gateCount_mapGates,
        QuantumCircuit.Party.PartyCircuit.participation_mapGates,
        QuantumCircuit.Party.PartyCircuit.participants_mapGatesLocationsEquiv,
        QuantumCircuit.Party.PartyCircuit.IsAllowed.mapGates}
    \leanok
    \uses{def:party_relabel, def:party_allowed}
    Let $\varphi$ be a morphism of gate families. The list of gates of $\varphi_*c$ is
    the list of gates of $c$ with every gate replaced by its image, cells and systems
    kept. Hence $\varphi_*c$ has the cells, the number of gates, the gate count and the
    participations of $c$, and corresponding occurrences have the same participants.
    Let both families have one outcome per gate and let $\varphi$ be compatible with
    interpretations of the two families. Then $T_{\varphi_*c}=T_c$. If moreover $c$ is
    allowed for a condition $\mathcal A$, and the image of every gate satisfying
    $\mathcal A$ satisfies a condition $\mathcal A_2$, then $\varphi_*c$ is allowed for
    $\mathcal A_2$.
\end{theorem}
\begin{proof}\leanok
    \uses{thm:party_circuit_op, def:circuits_single, def:circuits_gate_morphism}
    By induction on $c$. For a gate $g$ on $D$, compatibility and uniqueness of the
    outcome give $M_{\varphi(g)}=M_g$, so the placements on $D$ agree
    (Theorem~\ref{thm:party_circuit_op}) and the norms agree; local operations and
    compositions are unchanged.
\end{proof}

\subsection{Fresh pair registers}
\label{sec:party_syntax_add_pair}

\begin{definition}[Adding fresh factors at two parties]
    \label{def:party_add_pair}
    \lean{QuantumCircuit.Party.addPair, QuantumCircuit.Party.addPairLeft,
        QuantumCircuit.Party.addPairRight, QuantumCircuit.Party.addPair_of_ne,
        QuantumCircuit.Party.addPair_eq_of_notMem, QuantumCircuit.Party.addPairConfigEquiv,
        QuantumCircuit.Party.addPairConfigEquiv_apply_fst,
        QuantumCircuit.Party.addPairConfigEquiv_apply_snd_left,
        QuantumCircuit.Party.addPairConfigEquiv_apply_snd_right,
        QuantumCircuit.Party.addPairConfigEquiv_apply_snd_of_ne,
        QuantumCircuit.Party.addPairConfigEquiv_symm_apply_left,
        QuantumCircuit.Party.addPairConfigEquiv_symm_apply_right,
        QuantumCircuit.Party.addPairConfigEquiv_symm_apply_of_ne,
        QuantumCircuit.Party.pairFactor, QuantumCircuit.Party.pairFactor.unique,
        QuantumCircuit.Party.addPairSiteEquiv, QuantumCircuit.Party.card_pairFactor,
        QuantumCircuit.Party.card_addPair, QuantumCircuit.Party.pairFactorConfigEquiv,
        QuantumCircuit.Party.addPairMatrix}
    \leanok
    \uses{def:sites_config}
    Let $p\neq q$ be parties and $\alpha,\beta$ finite sets. The system
    $X^{pq}_{\alpha\beta}$ has the basis $\alpha\times X_p$ at $p$,
    $\beta\times X_q$ at $q$, and $X_s$ at every other party; it agrees with $X$
    outside $\{p,q\}$. Its configurations are pairs of a configuration of the fresh
    factors and a configuration of $X$: the bijection
    $X^{pq}_{\alpha\beta,P}\to(\alpha\times\beta)\times X_P$ sends $z$ to $((a,b),x)$,
    where $z_p=(a,x_p)$, $z_q=(b,x_q)$ and $z_s=x_s$ for $s\notin\{p,q\}$
    \cite[Section~5.1]{OpenAI2026PolynomialPEPS}. The fresh factors alone form the
    system $F^{pq}_{\alpha\beta}$ with the basis $\alpha$ at $p$, $\beta$ at $q$, and a
    single basis vector at every other party; its bases are single points when $\alpha$
    and $\beta$ are, its configurations are the pairs $\alpha\times\beta$, and
    \begin{align}
        |F^{pq}_{\alpha\beta,s}|
         & =\bigl(|\alpha|\text{ if }s=p\text{, else }1\bigr)
        \bigl(|\beta|\text{ if }s=q\text{, else }1\bigr). \notag
    \end{align}
    At every party $s$ the basis of $X^{pq}_{\alpha\beta}$ is
    $F^{pq}_{\alpha\beta,s}\times X_s$, so
    $|X^{pq}_{\alpha\beta,s}|=|F^{pq}_{\alpha\beta,s}|\,|X_s|$. For matrices $C_s$
    from $\C^{Y_s}$ to $\C^{Y'_s}$, the operation $C_s$ beside the fresh factors is the
    matrix $1\otimes C_s$ from $\C^{Y^{pq}_{\alpha\beta,s}}$ to
    $\C^{Y'^{pq}_{\alpha\beta,s}}$, read through these bases.
\end{definition}

\begin{theorem}[Party operations beside the fresh factors]
    \label{thm:party_add_pair_matrix}
    \lean{QuantumCircuit.Party.addPairConfigEquiv_eq_siteEquiv,
        QuantumCircuit.Party.norm_addPairMatrix_le,
        QuantumCircuit.Party.piKronecker_addPairMatrix}
    \leanok
    \uses{def:party_add_pair, def:elementary_finite_tensor_product}
    The bijection of configurations of $X^{pq}_{\alpha\beta}$ is read party by party:
    it is the composite of the bijections
    $X^{pq}_{\alpha\beta,s}\to F^{pq}_{\alpha\beta,s}\times X_s$ at every party, of the
    splitting of a family of pairs into a configuration of $F^{pq}_{\alpha\beta}$ and
    one of $X$, and of the identification of the configurations of
    $F^{pq}_{\alpha\beta}$ with $\alpha\times\beta$. For matrices $C_s$ as in Definition~\ref{def:party_add_pair},
    the operation beside the fresh factors has norm at most $\lVert C_s\rVert$, and
    their Kronecker product over all parties is
    \begin{align}
        \bigotimes_s(1\otimes C_s) & =1_{\alpha\times\beta}\otimes\bigotimes_sC_s, \notag
    \end{align}
    the right side read along the bijections of configurations of
    $Y^{pq}_{\alpha\beta}$ and $Y'^{pq}_{\alpha\beta}$.
\end{theorem}
\begin{proof}\leanok
    \uses{thm:detailed_balance_tensor_norm, thm:sites_product_span}
    The first statement is checked on the fresh parts and at every party separately.
    Reindexing does not increase the norm, and
    $\lVert1\otimes C\rVert\le\lVert C\rVert$
    (Theorem~\ref{thm:detailed_balance_tensor_norm}). For the Kronecker product,
    compute an entry with the first statement: the product over the parties of the
    entries of the identity on the fresh factors is the entry of the identity on the
    configurations of $F^{pq}_{\alpha\beta}$, that is of $1_{\alpha\times\beta}$, and
    the product of the entries of the $C_s$ is the entry of $\bigotimes_sC_s$
    (Theorem~\ref{thm:sites_product_span}).
\end{proof}

\subsection{Pair sources, pair effects and monomials}
\label{sec:party_syntax_pair_gates}

For a bijection $e\colon m\to\alpha\times n$ and $v\in\C^\alpha$, the preparation
$J_e(v)$ is the matrix from $\C^n$ to $\C^m$ with entry $v_a$ at $(y,x)$ when
$e(y)=(a,x)$ and $0$ otherwise; it maps $\ket x$ to $v\otimes\ket x$ read through $e$.
The effect $\bra v\otimes1$ from $\C^{\alpha\times n}$ to $\C^n$ is its adjoint for
$e$ the identity, and $J_e(v)^\dagger$ is this effect with its columns read through $e$ (Definitions~\ref{def:construction_idle_gates}
and~\ref{def:party_effect_matrices}).

\begin{definition}[Pair sources, pair effects and monomials]
    \label{def:party_pair_gates}
    \lean{QuantumCircuit.Party.PairSource, QuantumCircuit.Party.SourceInventory,
        QuantumCircuit.Party.SourceInventory.IsNormalized,
        QuantumCircuit.Party.PairSource.partyPair,
        QuantumCircuit.Party.SourceGate, QuantumCircuit.Party.SourceGate.left_mem,
        QuantumCircuit.Party.SourceGate.right_mem, QuantumCircuit.Party.SourceGate.ofAddPair,
        QuantumCircuit.Party.SourceGate.regionMatrix,
        QuantumCircuit.Party.SourceGate.outputEquiv, QuantumCircuit.Party.SourceGate.ev,
        QuantumCircuit.Party.SourceGate.IsAllowed,
        QuantumCircuit.Party.SourceGate.ev_ofAddPair_apply,
        QuantumCircuit.Party.SourceGate.reindexOutput,
        QuantumCircuit.Party.SourceGate.ev_reindexOutput,
        QuantumCircuit.Party.SourceGate.ne_left_of_notMem,
        QuantumCircuit.Party.SourceGate.ne_right_of_notMem,
        QuantumCircuit.Party.SourceGate.toAddPair,
        QuantumCircuit.Party.SourceGate.addPairLeft_toAddPair,
        QuantumCircuit.Party.SourceGate.addPairRight_toAddPair,
        QuantumCircuit.Party.SourceGate.toAddPair_of_notMem,
        QuantumCircuit.Party.SourceGate.configEquiv,
        QuantumCircuit.Party.SourceGate.configEquiv_apply_fst,
        QuantumCircuit.Party.SourceGate.configEquiv_apply_snd_left,
        QuantumCircuit.Party.SourceGate.configEquiv_apply_snd_right,
        QuantumCircuit.Party.SourceGate.configEquiv_apply_snd_of_notMem,
        QuantumCircuit.Party.PairEffectGate, QuantumCircuit.Party.PairEffectGate.left_mem,
        QuantumCircuit.Party.PairEffectGate.right_mem,
        QuantumCircuit.Party.PairEffectGate.ofAddPair,
        QuantumCircuit.Party.PairEffectGate.toSourceGate,
        QuantumCircuit.Party.PairEffectGate.regionMatrix,
        QuantumCircuit.Party.PairEffectGate.reindexInput,
        QuantumCircuit.Party.PairEffectGate.ev, QuantumCircuit.Party.PairEffectGate.IsAllowed,
        QuantumCircuit.Party.PairEffectGate.ev_ofAddPair_apply,
        QuantumCircuit.Party.PairEffectGate.ev_reindexInput,
        QuantumCircuit.Party.PairEffectGate.toAddPair,
        QuantumCircuit.Party.PairEffectGate.configEquiv,
        QuantumCircuit.Party.ChainGate, QuantumCircuit.Party.chainEv,
        QuantumCircuit.Party.ChainGate.IsAllowed, QuantumCircuit.Party.Monomial}
    \leanok
    \uses{def:party_add_pair, def:party_circuit, def:circuits_gate_family,
        def:circuits_gate_sum, def:sites_pair_factor_config}
    A pair source datum consists of two distinct parties $l\neq r$, finite sets $\alpha$
    and $\beta$, and a vector $\eta\in\C^{\alpha\times\beta}$; its unordered pair of
    endpoints is $\{l,r\}$. A source inventory is a finite list of such data, repeated
    data counting as distinct occurrences, and it is normalized when every vector has
    norm $1$ \cite[Lemma~5.1]{OpenAI2026PolynomialPEPS}.

    A pair source on a cell $D$ from $X$ to $Y$ is a pair source datum with
    $D=\{l,r\}$, bijections $Y_l\to\alpha\times X_l$ and $Y_r\to\beta\times X_r$, and
    $Y_s=X_s$ for every $s\notin D$; thus $l,r\in D$, and a party outside $D$ is neither
    endpoint \cite[Section~5.1]{OpenAI2026PolynomialPEPS}. Its matrix on the cell, from
    $\C^{X_D}$ to $\C^{Y_D}$, prepares $\eta$ on the fresh factors and is the identity
    on the old ones: its entry at $(y,x)$ is $\eta_{ab}$ when $y_l\mapsto(a,x_l)$ and
    $y_r\mapsto(b,x_r)$, and $0$ otherwise. Correspondingly, a configuration of $D$ in
    $Y$ is a pair of a configuration of the fresh factors and a configuration of $D$ in
    $X$. A pair source has one outcome, its interpretation is its matrix, and it is
    allowed when $\lVert\eta\rVert=1$. The canonical source of $\eta$ on $\{p,q\}$ goes
    from $X$ to $X^{pq}_{\alpha\beta}$ with the identifications of
    Definition~\ref{def:party_add_pair}. Relabelling the output along bijections
    $e_s\colon Y'_s\to Y_s$, with $Y'_s=Y_s$ outside $D$, composes the two bijections of
    the source with $e_l$ and $e_r$ and reindexes the rows of the matrix by
    $y'\mapsto(e_s(y'_s))_s$. The output of a source is read in canonical form by the
    bijections $Y_s\to X^{lr}_{\alpha\beta,s}$ given at $l$ and $r$ by the bijections of
    the source and elsewhere by the identification $Y_s=X_s$; through them a
    configuration $y$ of $Y$ is the pair $((a,b),x)$ with
    $y_l\mapsto(a,x_l)$, $y_r\mapsto(b,x_r)$ and $x_s=y_s$ outside $D$.

    A pair effect on $D$ from $X$ to $Y$ has the same data with bijections
    $X_l\to\alpha\times Y_l$ and $X_r\to\beta\times Y_r$, and $Y_s=X_s$ outside $D$; it
    removes the factors $\alpha$ and $\beta$. Its matrix on the cell has the entry
    $\overline{\eta_{ab}}$ at $(y,x)$ when $x_l\mapsto(a,y_l)$ and $x_r\mapsto(b,y_r)$,
    and $0$ otherwise. The source with the same data goes from $Y$ to $X$. The
    canonical effect of $\eta$ on $\{p,q\}$ goes from $Y^{pq}_{\alpha\beta}$ to $Y$, a
    relabelling of the input reindexes the columns, and the input is read in canonical
    form through the source with the same data. A pair effect has one outcome, its
    interpretation is its matrix, and it is allowed when $\lVert\eta\rVert=1$
    \cite[Section~5.1]{OpenAI2026PolynomialPEPS}.

    The gates of monomials are the pair sources and the pair effects, with the
    interpretation of the sum of the two gate families; each has one outcome, and each
    is allowed when its vector is normalized. A monomial on a finite set $D$ of parties,
    from $X$ to $Y$ with $X,Y$ systems on $D$, is a party circuit on the parties of $D$
    whose gates are pair sources and pair effects.
\end{definition}

\begin{theorem}[Matrices of pair sources and pair effects]
    \label{thm:party_pair_gate_matrices}
    \lean{QuantumCircuit.Party.SourceGate.regionMatrix_eq_prepMatrix,
        QuantumCircuit.Party.PairEffectGate.regionMatrix_eq_conjTranspose,
        QuantumCircuit.Party.SourceGate.norm_ev_le, QuantumCircuit.Party.SourceGate.norm_ev,
        QuantumCircuit.Party.PairEffectGate.norm_ev_le,
        QuantumCircuit.Party.PairEffectGate.norm_ev,
        QuantumCircuit.Party.SourceGate.eq_reindexOutput,
        QuantumCircuit.Party.SourceGate.eq_reindexOutput_ofAddPair}
    \leanok
    \uses{def:party_pair_gates, def:party_effect_matrices}
    Let $S$ be the matrix of a pair source on $D$ with vector $\eta$. Then $S$ is the
    preparation $J_e(\eta)$ along the bijection $e\colon Y_D\to(\alpha\times\beta)\times X_D$,
    and the matrix of a pair effect is the adjoint of the matrix of the source with the
    same data. Both have norm at most $\lVert\eta\rVert$, with equality when $D$ has a
    configuration in the input of the source, respectively in the output of the effect;
    a normalized source or effect is a contraction. Two pair sources on $D$ from $X$
    with the same pair source datum differ by a relabelling of the output, and every
    pair source is, with $D=\{l,r\}$, the canonical source of its datum followed by a
    relabelling of the output.
\end{theorem}
\begin{proof}\leanok
    \uses{def:sites_pair_factor_config, def:party_effect_matrices}
    The entries of $S$ and of the preparation agree once the configurations of
    $D=\{l,r\}$ are read as pairs (Definition~\ref{def:sites_pair_factor_config}), and
    the entries of the effect are the conjugates of those of the source, transposed.
    The norms are those of a preparation (Definition~\ref{def:party_effect_matrices}),
    and $\lVert M^\dagger\rVert=\lVert M\rVert$. For two sources with bijections
    $u_l,u_r$ and $u'_l,u'_r$, the relabelling is $u_l^{-1}u'_l$ at $l$,
    $u_r^{-1}u'_r$ at $r$, and the identification of equal bases elsewhere; the
    canonical source is one such source.
\end{proof}

\begin{theorem}[The operator of one pair source or pair effect]
    \label{thm:party_pair_gate_op}
    \lean{QuantumCircuit.Party.SourceGate.configEquiv_ofAddPair,
        QuantumCircuit.Party.SourceGate.op_gate, QuantumCircuit.Party.PairEffectGate.op_gate,
        QuantumCircuit.Party.PairEffectGate.op_gate_eq_effectMatrix,
        QuantumCircuit.Party.ChainGate.op_gate_inl, QuantumCircuit.Party.ChainGate.op_gate_inr}
    \leanok
    \uses{def:party_pair_gates, def:party_effect_matrices, def:party_circuit}
    Let $g$ be a pair source on $D$ from $X$ to $Y$ with vector $\eta$, and let
    $\kappa\colon Y_P\to(\alpha\times\beta)\times X_P$ be the reading of its output in
    canonical form; for the canonical source $\kappa$ is the bijection of
    configurations of $X^{pq}_{\alpha\beta}$. The party circuit of the one gate $g$ has
    the operator $J_\kappa(\eta)$, which maps $\ket x$ to $\eta\otimes\ket x$.
    For a pair effect from $X$ to $Y$ with input read through
    $\kappa\colon X_P\to(\alpha\times\beta)\times Y_P$, the operator of its one gate is
    $J_\kappa(\eta)^\dagger$, the effect $\bra\eta\otimes1$ with its columns read
    through $\kappa$. Among the gates of monomials, one pair source and one pair effect
    have the same operators.
\end{theorem}
\begin{proof}\leanok
    \uses{thm:party_circuit_op, def:sites_place_between, lem:sites_place_between_algebra,
        thm:party_pair_gate_matrices, def:party_effect_matrices, def:circuits_gate_sum}
    The canonical reading is checked coordinate by coordinate. By
    Theorem~\ref{thm:party_circuit_op} the operator of the source is the placement
    $S\otimes1$ between $X$ and $Y$ on $D$, whose entry at $(y,x)$ is the entry of $S$
    when $y$ and $x$ agree outside $D$ (Definition~\ref{def:sites_place_between}). The
    two conditions, $y_l\mapsto(a,x_l)$, $y_r\mapsto(b,x_r)$, and agreement outside
    $D$, together say $\kappa(y)=((a,b),x)$, and the entry is then $\eta_{ab}$. For the
    effect, $(E\otimes1)=(S'^\dagger\otimes1)=(S'\otimes1)^\dagger$ for the source $S'$
    with the same data (Theorem~\ref{thm:party_pair_gate_matrices} and
    Lemma~\ref{lem:sites_place_between_algebra}), and the adjoint of a preparation is
    the effect with its columns reindexed (Definition~\ref{def:party_effect_matrices}).
    The interpretation of the sum of the two gate families restricts to the
    interpretations of each.
\end{proof}

\subsection{Weighted gates and original circuits}
\label{sec:party_syntax_original}

\begin{definition}[Weighted gates and effect circuits]
    \label{def:party_weighted_gate}
    \lean{QuantumCircuit.Party.WeightedGate, QuantumCircuit.Party.WeightedGate.ev,
        QuantumCircuit.Party.WeightedGate.IsAllowed, QuantumCircuit.Party.EffectCircuit}
    \leanok
    \uses{def:party_pair_gates, def:party_allowed}
    A weighted gate on a cell $D$ from $X$ to $Y$, its gate expansion, is a finite list
    of pairs $(\lambda_k,M_k)$ of a coefficient $\lambda_k\in\C$ and a monomial $M_k$ on
    the parties of $D$ from $X|_D$ to $Y|_D$
    \cite[Section~5.1]{OpenAI2026PolynomialPEPS}. It has one outcome, and its matrix,
    from $\C^{X_D}$ to $\C^{Y_D}$, is $\sum_k\lambda_kT_{M_k}$. It is allowed when every
    monomial $M_k$ is allowed: its local operations are contractions and its pair
    sources and pair effects are normalized. An effect circuit is a party circuit whose
    gates are weighted gates, with no further condition.
\end{definition}

\begin{definition}[Original gates and original circuits]
    \label{def:party_original_circuit}
    \lean{QuantumCircuit.Party.OriginalGate, QuantumCircuit.Party.OriginalGate.ev,
        QuantumCircuit.Party.OriginalGate.toWeightedGateMorphism,
        QuantumCircuit.Party.OriginalGate.isCompatible_toWeightedGateMorphism,
        QuantumCircuit.Party.OriginalGate.ofWeightedGate,
        QuantumCircuit.Party.OriginalGate.toWeightedGate_ofWeightedGate,
        QuantumCircuit.Party.OriginalCircuit, QuantumCircuit.Party.OriginalCircuit.done,
        QuantumCircuit.Party.OriginalCircuit.comp, QuantumCircuit.Party.OriginalCircuit.gate,
        QuantumCircuit.Party.OriginalCircuit.localOp,
        QuantumCircuit.Party.OriginalCircuit.localOps,
        QuantumCircuit.Party.OriginalCircuit.toPartyCircuit_done,
        QuantumCircuit.Party.OriginalCircuit.toPartyCircuit_comp,
        QuantumCircuit.Party.OriginalCircuit.toPartyCircuit_gate,
        QuantumCircuit.Party.OriginalCircuit.toPartyCircuit_localOp,
        QuantumCircuit.Party.OriginalCircuit.toPartyCircuit_localOps,
        QuantumCircuit.Party.OriginalCircuit.op, QuantumCircuit.Party.OriginalCircuit.op_done,
        QuantumCircuit.Party.OriginalCircuit.toEffectCircuit}
    \leanok
    \uses{def:party_weighted_gate, def:party_relabel, def:circuits_gate_morphism}
    An original gate on a cell $D$ from $X$ to $Y$ is a weighted gate on $D$ that is not
    private, $|D|\neq1$, whose monomials are allowed and whose matrix is a contraction;
    the empty cell is allowed \cite[Section~5.1]{OpenAI2026PolynomialPEPS}. It has one
    outcome and its matrix is that of its expansion. Forgetting the conditions is a
    morphism from original gates to weighted gates compatible with the
    interpretations, and a weighted gate satisfying the three conditions is the
    expansion of an original gate. An original circuit from $X$ to $Y$ is a party
    circuit whose gates are original gates and whose local operations are
    contractions. Original circuits are built by the empty circuit, one original gate,
    one contraction at one party, contractions at every party, and composition, and
    these are the corresponding party circuits. The operator $T_w$ of an original
    circuit $w$ is that of its party circuit, the empty circuit having the operator
    $1$, and its effect circuit is the image of its party circuit under the forgetful
    morphism.
\end{definition}

\begin{theorem}[Operators of original circuits]
    \label{thm:party_original_op}
    \lean{QuantumCircuit.Party.OriginalCircuit.induction,
        QuantumCircuit.Party.OriginalCircuit.op_comp,
        QuantumCircuit.Party.OriginalCircuit.op_gate,
        QuantumCircuit.Party.OriginalCircuit.op_localOp,
        QuantumCircuit.Party.OriginalCircuit.op_localOps,
        QuantumCircuit.Party.PartyCircuit.isAllowed_of_isLocalContractive,
        QuantumCircuit.Party.OriginalCircuit.isAllowed,
        QuantumCircuit.Party.OriginalCircuit.norm_op_le_one}
    \leanok
    \uses{def:party_original_circuit, def:party_allowed, def:sites_place_between,
        def:elementary_finite_tensor_product}
    A property of original circuits holds for all of them when it holds for the empty
    circuits, for one original gate, for one contraction at one party and for
    contractions at every party, and passes to compositions. The operator of a
    composition is $T_{w'\circ w}=T_{w'}T_w$, that of one original gate is the
    placement between $X$ and $Y$ on its cell of the matrix of its expansion, that of
    one contraction at a party $p$ is its placement on $\{p\}$, and that of
    contractions at every party is their Kronecker product. A locally contractive
    party circuit of original gates is allowed for the condition that the monomials of
    every gate are allowed; hence every original circuit is allowed, and
    $\lVert T_w\rVert\le1$ \cite[Section~5.1]{OpenAI2026PolynomialPEPS}.
\end{theorem}
\begin{proof}\leanok
    \uses{thm:party_circuit_op, thm:party_op_contraction}
    The induction is the induction on the underlying party circuit, the contractivity
    of the local operations of a composition splitting into that of the two parts. The
    operators are those of Theorem~\ref{thm:party_circuit_op}. Allowedness is checked
    by induction on the party circuit: every original gate is a contraction with
    allowed monomials. Theorem~\ref{thm:party_op_contraction} gives the norm bound.
\end{proof}

\begin{theorem}[The effect circuit of an original circuit]
    \label{thm:party_original_effect}
    \lean{QuantumCircuit.Party.PartyCircuit.forallGate_card_ne_one,
        QuantumCircuit.Party.PartyCircuit.forallGate_card_ne_one_mapGates,
        QuantumCircuit.Party.OriginalCircuit.forallGate_card_ne_one,
        QuantumCircuit.Party.OriginalCircuit.op_toEffectCircuit,
        QuantumCircuit.Party.OriginalCircuit.isAllowed_toEffectCircuit,
        QuantumCircuit.Party.OriginalCircuit.nonprivate}
    \leanok
    \uses{def:party_original_circuit, def:party_gate_counts, def:party_allowed}
    No gate of a party circuit of original gates is private, and neither is any gate of
    its image under the forgetful morphism; in particular no gate of an original
    circuit or of its effect circuit is private
    \cite[Section~5.1]{OpenAI2026PolynomialPEPS}. The effect circuit of an original
    circuit $w$ has the operator $T_w$ and is allowed for the condition that the
    monomials of every weighted gate are allowed.
\end{theorem}
\begin{proof}\leanok
    \uses{thm:party_map_gates, thm:party_original_op}
    The first statements follow by induction on the circuit from $|D|\neq1$ for every
    original gate, the forgetful morphism keeping the cells. The forgetful morphism is
    compatible with the interpretations and keeps the condition on the monomials, so
    Theorem~\ref{thm:party_map_gates} transfers the operator and allowedness of
    Theorem~\ref{thm:party_original_op}.
\end{proof}

\subsection{Words of pair sources}
\label{sec:party_syntax_words}

\begin{definition}[Words of pair sources]
    \label{def:party_word}
    \lean{QuantumCircuit.Party.Word, QuantumCircuit.Party.Word.IsAllowed,
        QuantumCircuit.Party.Word.op, QuantumCircuit.Party.IsSourceOnly,
        QuantumCircuit.Party.Word.source, QuantumCircuit.Party.Word.sourceAt,
        QuantumCircuit.Party.Word.localPrepare, QuantumCircuit.Party.Word.localDiscard,
        QuantumCircuit.Party.Word.prepareProduct, QuantumCircuit.Party.Word.sources,
        QuantumCircuit.Party.Word.sourceCount, QuantumCircuit.Party.Word.sources_comp,
        QuantumCircuit.Party.Word.sources_relabel, QuantumCircuit.Party.Word.sources_source,
        QuantumCircuit.Party.Word.sources_prepareProduct, QuantumCircuit.Party.Word.UsesOnly,
        QuantumCircuit.Party.Word.toChain}
    \leanok
    \uses{def:party_pair_gates, def:party_allowed, def:party_relabel,
        def:circuits_gate_sum}
    A word from $X$ to $Y$ is a party circuit whose gates are pair sources. It is
    allowed when its local operations are contractions and its pair sources are
    normalized, and an operator is source-only when it is the operator $T_w$ of an
    allowed word $w$ \cite[Lemma~5.1]{OpenAI2026PolynomialPEPS}. Particular words are:
    \begin{enumerate}
        \item the word of one pair source, and the canonical pair source word of
              $\eta\in\C^{\alpha\times\beta}$ on $p\neq q$, from $X$ to
              $X^{pq}_{\alpha\beta}$;
        \item the local preparation of $v\in\C^\alpha$ at a party $p$, the local
              operation $J_e(v)$ at $p$ for a bijection
              $e\colon Y_p\to\alpha\times X_p$, where $Y_s=X_s$ for $s\neq p$; with
              $\alpha=\alpha_1\times\alpha_2$ it prepares a pair vector at one party;
        \item the local bra of $v\in\C^\alpha$ at $p$, the local operation
              $J_e(v)^\dagger$ for a bijection $e\colon X_p\to\alpha\times Y_p$,
              where $Y_s=X_s$ for $s\neq p$;
        \item the preparation of the product vector $\bigotimes_p\phi_p$, for vectors
              $\phi_p\in\C^{X_p}$, from the system with a single basis vector at every
              party, by the local operation $\ket{\phi_p}$ at every party $p$.
    \end{enumerate}
    The source list of a word is the list of the data of its pair sources, those of the
    later part first: it is empty for the empty word and for local operations, it is
    $(g)$ for one source $g$, and it is the list of $w'$ followed by that of $w$ for
    $w'\circ w$; thus relabellings and product preparations have no source. The source
    count is the length of the source list. A word uses only the parties of a set $R$
    when its sources have cells inside $R$, its local operations at one party act at
    parties of $R$, and at every party $s\notin R$ its local operations at every party
    have $Y_s=X_s$ and are the identity. Every word is a monomial with no pair effect,
    each pair source read as a gate of monomials.
\end{definition}

\begin{theorem}[Operators of words]
    \label{thm:party_word_op}
    \lean{QuantumCircuit.Party.Word.op_comp, QuantumCircuit.Party.Word.op_source,
        QuantumCircuit.Party.Word.op_sourceAt,
        QuantumCircuit.Party.Word.op_source_eq_op_relabel_mul,
        QuantumCircuit.Party.Word.op_localPrepare,
        QuantumCircuit.Party.Word.op_localPrepare_eq_prepMatrix,
        QuantumCircuit.Party.Word.op_localDiscard,
        QuantumCircuit.Party.Word.op_prepareProduct_mulVec,
        QuantumCircuit.Party.Word.op_toChain}
    \leanok
    \uses{def:party_word, def:sites_pair_factor_config, def:sites_place_between,
        def:siteconfig_product_vector, def:party_effect_matrices}
    The operator of a composition of words is $T_{w'\circ w}=T_{w'}T_w$. The operator of
    one pair source is the preparation of its vector with the output read in canonical
    form (Theorem~\ref{thm:party_pair_gate_op}); for the canonical source on $p\neq q$ it
    is $J_e(\eta)$ along the bijection $e$ of configurations of
    $X^{pq}_{\alpha\beta}$, which maps $\ket x$ to $\eta\otimes\ket x$, and every pair
    source is the canonical source of its datum followed by the relabelling of its
    output into canonical form:
    \begin{align}
        T_g & =T_{R}\,T_{\mathrm{can}}, \notag
    \end{align}
    where $T_{\mathrm{can}}$ is the operator of the canonical source of the datum of $g$
    from $X$ to $X^{lr}_{\alpha\beta}$ and $R$ is the relabelling along the reading of
    the output of $g$ in canonical form.
    The local preparation of $v$ at $p$ has the operator $J_e(v)\otimes1$, the
    placement on $\{p\}$, which is the preparation $J_{\hat e}(v)$ along the
    bijection $\hat e\colon Y_P\to\alpha\times X_P$ that splits off the factor $\alpha$
    at $p$ (Definition~\ref{def:sites_pair_factor_config}). The local bra of $v$ at $p$
    has the operator $\bra v\otimes1$ with its columns read along the corresponding
    bijection $X_P\to\alpha\times Y_P$. The product preparation maps the unique basis
    vector to $\bigotimes_p\phi_p$. A word read as a monomial has the operator of the
    word.
\end{theorem}
\begin{proof}\leanok
    \uses{thm:party_circuit_op, thm:party_pair_gate_op, thm:party_relabel_op,
        def:sites_pair_factor_config, def:sites_place_between, thm:party_map_gates,
        def:circuits_gate_sum}
    The composition and the local operations are Theorem~\ref{thm:party_circuit_op},
    and one source is Theorem~\ref{thm:party_pair_gate_op}. The relabelling $R$ after
    the canonical source reindexes the rows of $J_e(\eta)$ along the reading in
    canonical form (Theorem~\ref{thm:party_relabel_op}), which is the operator of the
    source. For the local preparation compare entries: the entry of the placement at
    $(y,z)$ is $v_a$ when $e(y_p)=(a,z_p)$ and $y_s=z_s$ for $s\neq p$, which is the
    condition that $\hat e(y)=(a,z)$ (Definition~\ref{def:sites_pair_factor_config});
    the local bra is the same computation for the adjoint. For the product
    preparation, the Kronecker product of the columns $\ket{\phi_p}$ applied to the
    basis vector has the coefficient $\prod_p(\phi_p)_{x_p}$ at $x$. The inclusion of
    pair sources among the gates of monomials is compatible with the interpretations
    (Definition~\ref{def:circuits_gate_sum}), so Theorem~\ref{thm:party_map_gates}
    applies.
\end{proof}

\begin{theorem}[A fresh source commutes with local operations on the old factors]
    \label{thm:party_source_commutes}
    \lean{QuantumCircuit.Party.Word.op_sourceAt_mul_piKronecker}
    \leanok
    \uses{def:party_word, thm:party_add_pair_matrix}
    Let $p\neq q$, let $\eta\in\C^{\alpha\times\beta}$, and let $A_s$ be matrices from
    $\C^{Y_s}$ to $\C^{Y'_s}$. Local operations followed by the canonical source equal
    the canonical source followed by the same operations beside the fresh factors:
    \begin{align}
        T_{\mathrm{can},Y'}\bigotimes_sA_s & =\bigotimes_s(1\otimes A_s)\,T_{\mathrm{can},Y},
        \notag
    \end{align}
    where $T_{\mathrm{can},Y}$ is the operator of the canonical source of $\eta$ on
    $\{p,q\}$ from $Y$ to $Y^{pq}_{\alpha\beta}$
    \cite[Theorem~5.2]{OpenAI2026PolynomialPEPS}.
\end{theorem}
\begin{proof}\leanok
    \uses{thm:party_word_op, thm:party_add_pair_matrix, def:party_effect_matrices}
    By Theorem~\ref{thm:party_word_op} both canonical sources are preparations of
    $\eta$ along the bijections of configurations, and by
    Theorem~\ref{thm:party_add_pair_matrix} the right-hand Kronecker product is
    $1_{\alpha\times\beta}\otimes\bigotimes_sA_s$ along the same bijections. At the
    configuration $((a,b),y')$ and the input $x$ both sides have the entry
    $\eta_{ab}(\bigotimes_sA_s)_{y'x}$.
\end{proof}

\begin{theorem}[Allowed words]
    \label{thm:party_word_allowed}
    \lean{QuantumCircuit.Party.Word.isAllowed_done, QuantumCircuit.Party.Word.IsAllowed.comp,
        QuantumCircuit.Party.Word.isAllowed_localOp,
        QuantumCircuit.Party.Word.isAllowed_localOps,
        QuantumCircuit.Party.Word.isAllowed_source,
        QuantumCircuit.Party.Word.isAllowed_localPrepare,
        QuantumCircuit.Party.Word.isAllowed_localDiscard,
        QuantumCircuit.Party.Word.isAllowed_relabel,
        QuantumCircuit.Party.Word.isAllowed_sourceAt,
        QuantumCircuit.Party.Word.isAllowed_prepareProduct,
        QuantumCircuit.Party.Word.IsAllowed.toChain}
    \leanok
    \uses{def:party_word}
    The empty word, compositions of allowed words, contractions at one party and at
    every party, and relabellings are allowed words. One pair source with a normalized
    vector, in particular the canonical source of a unit vector, is allowed. The local
    preparation and the local bra of a unit vector are allowed, and so is the
    preparation of a product of unit vectors. An allowed word read as a monomial is an
    allowed monomial.
\end{theorem}
\begin{proof}\leanok
    \uses{thm:party_pair_gate_matrices, def:party_effect_matrices,
        lem:matrix_l2_contractions, thm:party_map_gates}
    A normalized source is a contraction
    (Theorem~\ref{thm:party_pair_gate_matrices}). A preparation of a unit vector is a
    contraction (Definition~\ref{def:party_effect_matrices}), and so is its adjoint;
    the column $\ket{\phi_p}$ is the preparation of $\phi_p$ on a single basis vector.
    The matrix of a bijection is a contraction (Lemma~\ref{lem:matrix_l2_contractions}).
    Reading a word as a monomial keeps allowedness by Theorem~\ref{thm:party_map_gates}.
\end{proof}

\begin{theorem}[Source lists of words]
    \label{thm:party_word_sources}
    \lean{QuantumCircuit.Party.Word.numGates_eq_sourceCount,
        QuantumCircuit.Party.Word.forallGate_card_ne_one,
        QuantumCircuit.Party.Word.gateCount_eq_sourceCount,
        QuantumCircuit.Party.Word.isNormalized_sources,
        QuantumCircuit.Party.Word.usesOnly_relabel,
        QuantumCircuit.Party.Word.UsesOnly.mem_sources,
        QuantumCircuit.Party.Word.gates_toChain}
    \leanok
    \uses{def:party_word, def:party_gate_counts}
    For every word $w$, the number of gates and the gate count are both the source
    count, since no pair source is private; the gates of $w$ read as a monomial are its
    pair sources. The source list of an allowed word is normalized. A relabelling
    that is the identification of equal bases outside $R$ uses only the parties of $R$,
    and every source of a word using only the parties of $R$ joins two parties of $R$.
\end{theorem}
\begin{proof}\leanok
    \uses{thm:party_gate_count_relations, thm:party_map_gates}
    By induction on $w$. A source has the cell $\{l,r\}$ with $l\neq r$, of two
    parties, so Theorem~\ref{thm:party_gate_count_relations} gives $\#(w)=n(w)$. The
    gates of the monomial are given by Theorem~\ref{thm:party_map_gates}. The other
    statements are read off the rules.
\end{proof}

\begin{definition}[Sources and effects of monomials]
    \label{def:party_monomial_effects}
    \lean{QuantumCircuit.Party.Monomial.sources, QuantumCircuit.Party.Monomial.effects,
        QuantumCircuit.Party.Monomial.effectParties,
        QuantumCircuit.Party.Monomial.effectCount,
        QuantumCircuit.Party.Monomial.sourceCount}
    \leanok
    \uses{def:party_pair_gates, def:party_gate_counts}
    For a monomial $M$, its sources are the data of its pair sources, the latest first,
    and its effects are the data of its pair effects in chronological order, read off
    the list of its gates. The endpoints of its effects form the list of the pairs
    $(l,r)$ of the effects. The effect count and the source count are the lengths of
    the two lists \cite[Section~5.1]{OpenAI2026PolynomialPEPS}.
\end{definition}

\begin{theorem}[Counting the sources and effects of a monomial]
    \label{thm:party_monomial_effects}
    \lean{QuantumCircuit.Party.Monomial.effectCount_add_sourceCount,
        QuantumCircuit.Party.Monomial.effectCount_comp,
        QuantumCircuit.Party.Monomial.sourceCount_comp,
        QuantumCircuit.Party.Monomial.effects_comp,
        QuantumCircuit.Party.Monomial.norm_vector_of_mem_effects,
        QuantumCircuit.Party.Monomial.effects_toChain}
    \leanok
    \uses{def:party_monomial_effects, def:party_allowed, def:party_word}
    For a monomial $M$, the effect count plus the source count is $n(M)$. Under
    composition the effect counts and the source counts add, and the effects of
    $M'\circ M$ are those of $M$ followed by those of $M'$. The effects of an allowed
    monomial are normalized, and a word read as a monomial has no effect.
\end{theorem}
\begin{proof}\leanok
    \uses{thm:party_gate_count_relations, thm:party_word_sources}
    Every gate of a monomial is a source or an effect, and the number of gates is the
    length of the list of gates (Theorem~\ref{thm:party_gate_count_relations}); count
    by induction on that list. The lists of gates of a composition are concatenated.
    Normalization follows by induction on $M$. The gates of a word read as a monomial
    are pair sources (Theorem~\ref{thm:party_word_sources}).
\end{proof}

\section{Operations on circuits on parties}
\label{sec:party_syntax_operations}

This section runs a circuit on parties beside untouched spectator factors, reads
a circuit without gates party by party, moves pair sources along a map of
parties, and reads a word on blocks of parties, on a subset of parties, and on
the parties of one cell. In each case the operator of the new circuit is the
operator of the old one tensored with an identity or relabelled along a
bijection of configurations, as in \cite[Section~5.1]{OpenAI2026PolynomialPEPS},
where an operation acts on some registers and leaves the others untouched.

Throughout, $P$ is a set of parties, a system is a family $X=(X_p)_{p\in P}$ of
finite basis sets, and, when $P$ is finite, $\mathcal H_X=\bigotimes_p\C^{X_p}$
has the configurations $X_P=\prod_pX_p$ as its basis. For a party circuit $c$
from $X$ to $Y$ whose gates have one outcome each and an interpretation
$\mathrm{ev}$ of its gates, $[c]$ denotes its operator from $\mathcal H_X$ to
$\mathcal H_Y$; for a word, the interpretation is that of pair sources. For a
region $D\subseteq P$ and a matrix $K$ from $\C^{X_D}$ to $\C^{Y_D}$, where $Y$
agrees with $X$ outside $D$, $K\otimes1$ is its placement between $X$ and $Y$
(Definition~\ref{def:sites_place_between}).

\subsection{Circuits beside spectator factors}
\label{sec:party_frame}

Let $W=(W_p)_{p\in P}$ be a system of finite basis sets, the spectators. The
system $W\times X$ has the basis $W_p\times X_p$ at $p$, and its configurations
on a region $D$ are paired with $W_D\times X_D$.

\begin{definition}[Circuits beside spectator factors]
    \label{def:party_gate_frame}
    \lean{QuantumCircuit.Party.GateFrame, QuantumCircuit.Party.GateFrame.IsCompatible,
        QuantumCircuit.Party.SourceGate.frame,
        QuantumCircuit.Party.SourceGate.toPairSource_frame,
        QuantumCircuit.Party.PairEffectGate.frame,
        QuantumCircuit.Party.PairEffectGate.toPairSource_frame,
        QuantumCircuit.Party.PartyCircuit.frame}
    \leanok
    \uses{def:party_circuit, def:party_pair_gates, def:party_weighted_gate,
        def:circuits_gate_family}
    A family of gates runs beside spectator factors when every gate $g$ on a cell
    $D$ from $X$ to $Y$ and every spectator system $W$ give a gate $g^W$ on $D$ from
    $W\times X$ to $W\times Y$, with the outcomes of $g^W$ in bijection with those
    of $g$. An interpretation $\mathrm{ev}$ is compatible with framing when, for
    every outcome $o$ of $g^W$ with corresponding outcome $o'$ of $g$,
    \begin{align}
        \mathrm{ev}(g^W,o)=1_{W_D}\otimes\mathrm{ev}(g,o'),
        \label{eq:party_frame_compatible}
    \end{align}
    read along the pairings $(W\times X)_D\cong W_D\times X_D$ and
    $(W\times Y)_D\cong W_D\times Y_D$. A pair source or a pair effect beside $W$
    keeps its endpoints, its fresh bases and its vector, the fresh factor placed in
    front of the spectator factor at each endpoint. The gates of monomials are
    framed as pair sources or pair effects, and a weighted gate on $D$ is framed by
    framing each of its monomials beside the spectators $(W_p)_{p\in D}$, with the
    same coefficients. The circuit $c^W$ beside $W$ is defined by induction on
    the circuit $c$: the empty circuit stays empty, a gate $g$ becomes $g^W$ on the
    same cell, a local operation $A$ at $p$ becomes $1_{W_p}\otimes A$ at $p$, local
    operations $A_p$ at every party become $1_{W_p}\otimes A_p$, and a composition
    becomes the composition of the framed circuits.
\end{definition}

\begin{theorem}[Framing keeps the syntax]
    \label{thm:party_frame_syntax}
    \lean{QuantumCircuit.Party.PartyCircuit.UsesOnly.frame,
        QuantumCircuit.Party.PartyCircuit.gateCells_frame,
        QuantumCircuit.Party.PartyCircuit.numGates_frame,
        QuantumCircuit.Party.PartyCircuit.gateCount_frame,
        QuantumCircuit.Party.PartyCircuit.participation_frame,
        QuantumCircuit.Party.PartyCircuit.participationCount_frame,
        QuantumCircuit.Party.PartyCircuit.frameLocationsEquiv,
        QuantumCircuit.Party.PartyCircuit.participants_frameLocationsEquiv,
        QuantumCircuit.Party.PartyCircuit.mapGates_frame,
        QuantumCircuit.Party.Word.sources_frame,
        QuantumCircuit.Party.Monomial.sources_frame,
        QuantumCircuit.Party.Monomial.effects_frame,
        QuantumCircuit.Party.Monomial.effectCount_frame,
        QuantumCircuit.Party.Monomial.sourceCount_frame,
        QuantumCircuit.Party.PartyCircuit.forallGate_isRight_frame}
    \leanok
    \uses{def:party_gate_frame, def:party_gate_counts, def:party_word,
        def:party_pair_gates, def:party_relabel, def:circuits_gate_morphism}
    Let $c$ be a party circuit and $W$ a spectator system. If $c$ uses only the
    parties of a set $R$, so does $c^W$. The circuit $c^W$ has the same list of
    gate cells, the same number of gates, the same gate count, and, at every party
    $p$, the same participation and the same number of gates containing $p$. Its
    gate occurrences are in bijection with those of $c$, corresponding occurrences
    having the same cell. If a morphism $\varphi$ of gate families satisfies
    $\varphi(g^W)=\varphi(g)^W$ for every gate $g$, then mapping $c^W$ along
    $\varphi$ gives the circuit beside $W$ of the image of $c$. A word beside $W$
    has the sources of the word. A monomial beside $W$ has the pair sources and
    the pair effects of the monomial, hence the same numbers of each, and a chain
    all of whose gates are pair effects stays such a chain.
\end{theorem}
\begin{proof}\leanok
    By induction on the circuit. Framing does not change the cell of a gate or the
    pair data of a source or an effect, and it adds no gate. For the parties
    outside $R$ at which the local operations $A_p$ are the identifications of
    equal bases $Y_p=X_p$, the matrices $1_{W_p}\otimes A_p$ are the
    identifications of $W_p\times Y_p=W_p\times X_p$. The list of gates of a
    composition is the concatenation of the two lists, so the lists of pair sources
    and pair effects of a monomial are those of the framed monomial.
\end{proof}

\begin{theorem}[Pair sources and pair effects beside spectators]
    \label{thm:party_frame_pair_compatible}
    \lean{QuantumCircuit.Party.SourceGate.regionMatrix_frame,
        QuantumCircuit.Party.PairEffectGate.regionMatrix_frame,
        QuantumCircuit.Party.SourceGate.isCompatible_frame,
        QuantumCircuit.Party.PairEffectGate.isCompatible_frame,
        QuantumCircuit.Party.chainEv_isCompatible_frame}
    \leanok
    \uses{def:party_gate_frame, def:party_pair_gates}
    The matrix of a pair source beside $W$ on its cell $D$ is $1_{W_D}\otimes M$,
    where $M$ is the matrix of the source on $D$, read along the pairings of
    configurations, and likewise for a pair effect. Hence the interpretations of
    pair sources, of pair effects and of the gates of monomials are compatible with
    framing.
\end{theorem}
\begin{proof}\leanok
    The cell $D$ consists of the two endpoints. The element of the matrix of the
    framed source at a pair of configurations vanishes unless the spectator parts
    agree at both endpoints, that is, on all of $D$, and it is then the element of
    $M$ at the remaining parts; this is the element of $1_{W_D}\otimes M$. The
    matrix of a pair effect is the adjoint of the matrix of the source with the same
    pair data, and the adjoint of $1_{W_D}\otimes M$, relabelled, is
    $1_{W_D}\otimes M^\dagger$, relabelled. A gate of a monomial is one or the
    other.
\end{proof}

\begin{theorem}[The operator of a circuit beside spectators]
    \label{thm:party_op_frame}
    \lean{QuantumCircuit.Party.PartyCircuit.op_frame,
        QuantumCircuit.Party.PartyCircuit.op_frame_unit,
        QuantumCircuit.Party.Word.op_frame}
    \leanok
    \uses{def:party_gate_frame, def:party_circuit, def:party_word}
    Let $P$ be finite, let the gates have one outcome each, and let $\mathrm{ev}$ be
    compatible with framing. For every party circuit $c$ from $X$ to $Y$ and every
    spectator system $W$,
    \begin{align}
        [c^W]=1_{W_P}\otimes[c],
        \label{eq:party_op_frame}
    \end{align}
    read along the pairings $(W\times X)_P\cong W_P\times X_P$ and
    $(W\times Y)_P\cong W_P\times Y_P$. In particular this holds for every word.
    If every $W_p$ has one element, then the element of $[c^W]$ at $(y,x)$ is the
    element of $[c]$ at the second components of $y$ and $x$.
\end{theorem}
\begin{proof}\leanok
    \uses{thm:party_frame_pair_compatible, lem:sites_place_between_products,
        thm:sites_product_span, def:circuits_single}
    By induction on the circuit. The empty circuit gives $1=1\otimes1$. For a gate
    $g$ on $D$, the operator of $g^W$ is the placement of
    $\mathrm{ev}(g^W,o)=1_{W_D}\otimes\mathrm{ev}(g,o')$
    by~(\ref{eq:party_frame_compatible}), the outcomes being unique, and placing
    $1_{W_D}\otimes K$ between $W\times X$ and $W\times Y$ gives
    $1_{W_P}\otimes(K\otimes1)$ by Lemma~\ref{lem:sites_place_between_products}. A
    local operation $1_{W_p}\otimes A$, read on the one-party region $\{p\}$, is
    $1_{W_{\{p\}}}\otimes A$ read there, and the same lemma applies. For local
    operations at every party, $\bigotimes_p(1_{W_p}\otimes A_p)
        =1_{W_P}\otimes\bigotimes_pA_p$ after regrouping the factors. For a
    composition, $(1\otimes[c'])(1\otimes[c])=1\otimes[c'][c]$. If every $W_p$ has
    one element, $1_{W_P}$ is the $1\times1$ identity.
\end{proof}

\begin{theorem}[Weighted gates beside spectators]
    \label{thm:party_frame_weighted}
    \lean{QuantumCircuit.Party.WeightedGate.isCompatible_frame}
    \leanok
    \uses{def:party_gate_frame, def:party_weighted_gate}
    The interpretation of weighted gates is compatible with framing.
\end{theorem}
\begin{proof}\leanok
    \uses{thm:party_op_frame, thm:party_frame_pair_compatible}
    The matrix of a weighted gate is $\sum_tc_t[M_t]$ over its terms, and that of the
    framed gate is $\sum_tc_t[M_t^W]$. By Theorem~\ref{thm:party_op_frame}, applied
    to the monomials with the compatible interpretation of their gates,
    $[M_t^W]=1\otimes[M_t]$, read along the pairings. By induction on the list of
    terms, the Kronecker product with $1$ and the relabelling are linear, which
    gives $1\otimes\sum_tc_t[M_t]$.
\end{proof}

\begin{theorem}[Framing keeps circuits allowed]
    \label{thm:party_frame_allowed}
    \lean{QuantumCircuit.Party.PartyCircuit.isAllowed_frame,
        QuantumCircuit.Party.Word.isAllowed_frame,
        QuantumCircuit.Party.PartyCircuit.IsLocalContractive.frame,
        QuantumCircuit.Party.ChainGate.isAllowed_frame,
        QuantumCircuit.Party.WeightedGate.isAllowed_frame}
    \leanok
    \uses{def:party_gate_frame, def:party_allowed, def:party_word,
        def:party_pair_gates, def:party_weighted_gate}
    Let $\mathrm{ev}$ be compatible with framing and let a condition on gates pass
    from every gate $g$ to $g^W$. If a party circuit is allowed for this condition,
    so is the circuit beside $W$. A circuit beside $W$ is locally contractive when
    the circuit is, and a word, an allowed gate of a monomial, and an allowed
    weighted gate stay allowed beside $W$.
\end{theorem}
\begin{proof}\leanok
    \uses{thm:detailed_balance_tensor_norm, lem:matrix_l2_contractions,
        thm:party_frame_pair_compatible}
    By induction on the circuit. A gate $g^W$ has by~(\ref{eq:party_frame_compatible})
    the matrix $1\otimes\mathrm{ev}(g,o')$, relabelled, whose norm is at most
    $\lVert1\otimes\mathrm{ev}(g,o')\rVert\le\lVert\mathrm{ev}(g,o')\rVert\le1$, and
    it satisfies the condition on gates by hypothesis. A local operation
    $1\otimes A$ has $\lVert1\otimes A\rVert\le\lVert A\rVert\le1$. For words and
    monomials the interpretations of pair sources and of the gates of monomials are
    compatible by Theorem~\ref{thm:party_frame_pair_compatible}, and the condition on
    their gates concerns the vectors, which framing keeps, so the first statement
    applies; every framed monomial of an allowed weighted gate is then allowed.
\end{proof}

\subsection{Circuits without gates}
\label{sec:party_local}

\begin{definition}[Circuits without gates and their party operations]
    \label{def:party_local_circuit}
    \lean{QuantumCircuit.Party.PartyCircuit.localOpFamily,
        QuantumCircuit.Party.PartyCircuit.localOpFamily_self,
        QuantumCircuit.Party.PartyCircuit.localOpFamily_of_ne,
        QuantumCircuit.Party.PartyCircuit.IsLocal,
        QuantumCircuit.Party.PartyCircuit.isLocal_iff_gates_eq_nil,
        QuantumCircuit.Party.Word.isLocal_iff_sources_eq_nil,
        QuantumCircuit.Party.PartyCircuit.partyOp,
        QuantumCircuit.Party.PartyCircuit.regionOp,
        QuantumCircuit.Party.PartyCircuit.parties}
    \leanok
    \uses{def:party_circuit, def:party_gate_counts, def:party_word,
        thm:sites_product_span}
    The party family of a local operation $A$ at $p$ from $X$ to $Y$, where
    $Y_s=X_s$ for $s\neq p$, is $A$ at $p$ and the matrix of the identification
    $X_s=Y_s$ at every $s\neq p$. A party circuit is local when it has no gate,
    that is, when its list of gates is empty; a word is local exactly when it has
    no pair source. The party operation $c_p$ of $c$ at $p$ is defined by induction
    on $c$: it is $1$ for the empty circuit, $0$ for a gate, the party family at $p$
    for a local operation, $A_p$ for local operations $(A_s)_s$ at every party, and
    $c'_pc_p$ for the composition of $c$ followed by $c'$. For a finite region $R$,
    the region operation is $c_R=\bigotimes_{p\in R}c_p$. When $P$ is finite, the
    parties of $c$ are those named by its operations: the cells of its gates, the
    party of each local operation, whatever its matrix, and, for local operations
    $(A_s)_s$ at every party, the parties $p$ at which $A_p$ is not the matrix of an
    identification of equal bases $Y_p=X_p$.
\end{definition}

\begin{theorem}[A circuit without gates acts party by party]
    \label{thm:party_local_op}
    \lean{QuantumCircuit.Party.PartyCircuit.op_localOp_eq_piKronecker,
        QuantumCircuit.Party.PartyCircuit.op_eq_piKronecker_partyOp,
        QuantumCircuit.Party.PartyCircuit.op_eq_reindex_kronecker_regionOp}
    \leanok
    \uses{def:party_local_circuit, def:sites_config}
    Let $P$ be finite and let the gates have one outcome each. The operator of one
    local operation is the product operator $\bigotimes_sm_s$ of its party family
    $(m_s)_s$. The operator of a local circuit $c$ is the product of its party
    operations,
    \begin{align}
        [c]=\bigotimes_{p\in P}c_p,
        \label{eq:party_local_op}
    \end{align}
    and, for every region $R$, in the coordinates of the cuts along $R$,
    $[c]=c_R\otimes\bigotimes_{p\notin R}c_p$. This is the collection of the local
    maps of a word without pair sources by party in the proof of
    \cite[Theorem~5.2]{OpenAI2026PolynomialPEPS}.
\end{theorem}
\begin{proof}\leanok
    \uses{def:party_circuit, def:sites_place_between, thm:sites_product_span,
        lem:sites_product_intertwine}
    The element of the placement of $A$ at $(y,x)$ is $A_{y_px_p}$ if $y$ and $x$
    agree outside $p$ through the identifications, and $0$ otherwise; the element
    of $\bigotimes_sm_s$ is $A_{y_px_p}\prod_{s\neq p}(m_s)_{y_sx_s}$, and the product
    over $s\neq p$ is $1$ exactly when $y$ and $x$ agree outside $p$. The identity
    (\ref{eq:party_local_op}) follows by induction on the circuit: the empty circuit
    gives $1=\bigotimes_p1$, local operations at every party have the product
    operator as their operator, and for a composition
    $(\bigotimes_pc'_p)(\bigotimes_pc_p)=\bigotimes_pc'_pc_p$. The last statement is
    (\ref{eq:party_local_op}) with the product split along the cut, by
    Lemma~\ref{lem:sites_product_intertwine}.
\end{proof}

\begin{theorem}[The party operations of an allowed circuit are contractions]
    \label{thm:party_local_contraction}
    \lean{QuantumCircuit.Party.PartyCircuit.norm_localOpFamily_le_one,
        QuantumCircuit.Party.PartyCircuit.norm_partyOp_le_one,
        QuantumCircuit.Party.PartyCircuit.norm_regionOp_le_one,
        QuantumCircuit.Party.PartyCircuit.exists_eq_piKronecker,
        QuantumCircuit.Party.Word.exists_piKronecker_of_sources_eq_nil}
    \leanok
    \uses{def:party_local_circuit, def:party_allowed, def:party_word}
    The party family of a contraction consists of contractions. If a party circuit
    is allowed for some condition on its gates, then every party operation $c_p$ and
    every region operation $c_R$ is a contraction. If $P$ is finite and $c$ is moreover
    local, then $[c]=\bigotimes_pB_p$ for contractions $B_p$. In particular an
    allowed word $w$ without pair sources satisfies $[w]=\bigotimes_pA_p$ for
    contractions $A_p$ \cite[Theorem~5.2]{OpenAI2026PolynomialPEPS}.
\end{theorem}
\begin{proof}\leanok
    \uses{lem:matrix_l2_contractions, thm:party_local_op}
    The matrix of an identification of bases is a contraction. By induction on
    the circuit, $c_p$ is $1$, $0$, a member of the party family of a contraction,
    the contraction $A_p$, or a product $c'_pc_p$ of two contractions. A product of
    contractions on the parties of $R$ is a contraction. For a local circuit take
    $B_p=c_p$ in (\ref{eq:party_local_op}); a word without pair sources is local.
\end{proof}

\begin{theorem}[The parties of a word]
    \label{thm:party_parties}
    \lean{QuantumCircuit.Party.Word.usesOnly_iff_parties_subset,
        QuantumCircuit.Party.Word.usesOnly_parties}
    \leanok
    \uses{def:party_local_circuit, def:party_word}
    Let $P$ be finite. A word uses only the parties of a set $R$ if and only if its
    parties lie in $R$. In particular a word uses only its own parties.
\end{theorem}
\begin{proof}\leanok
    By induction on the word. A source uses only $R$ when its cell lies in $R$, a
    local operation at $p$ when $p\in R$, local operations at every party when they
    are identifications of equal bases at every party outside $R$, and a
    composition when both parts do; the parties of a composition are the union of
    those of its parts.
\end{proof}

\subsection{Maps of source owners}
\label{sec:party_owner_map}

\begin{definition}[Pair sources along a map of parties]
    \label{def:party_owner_map}
    \lean{QuantumCircuit.Party.PairSource.mapOwner,
        QuantumCircuit.Party.SourceInventory.mapOwner,
        QuantumCircuit.Party.SourceInventory.mapOwner_nil,
        QuantumCircuit.Party.SourceInventory.mapOwner_cons,
        QuantumCircuit.Party.SourceInventory.mapOwner_append,
        QuantumCircuit.Party.affectedOwner,
        QuantumCircuit.Party.affectedOwner_eq_none,
        QuantumCircuit.Party.isSome_affectedOwner}
    \leanok
    \uses{def:party_pair_gates}
    Let $f\colon P\to Q$ be a map of parties. A pair source $s$ with endpoints $l$
    and $r$ is moved along $f$ to the pair source $f_*s$ with endpoints $f(l)$ and
    $f(r)$, the same fresh bases and the same vector, when $f(l)\neq f(r)$; when
    $f(l)=f(r)$ it becomes local and is dropped. A list $S$ of pair sources is moved
    to the list $f_*S$ of the moved sources that are kept, in their order, so that
    $f_*$ of the empty list is empty and $f_*(S\,S')=f_*S\,f_*S'$ for a concatenation.
    For a set $A\subseteq P$ of affected parties, the owner map $\omega_A$ from $P$
    to $A\sqcup\{\ast\}$ sends each affected party to itself and every other party to
    the one exterior party $\ast$; thus $\omega_A(p)=\ast$ if and only if
    $p\notin A$. This is the grouping of parties in the proof of
    \cite[Theorem~5.2]{OpenAI2026PolynomialPEPS}.
\end{definition}

\begin{theorem}[Moved lists of pair sources]
    \label{thm:party_owner_map}
    \lean{QuantumCircuit.Party.SourceInventory.mem_mapOwner,
        QuantumCircuit.Party.SourceInventory.IsNormalized.mapOwner,
        QuantumCircuit.Party.SourceInventory.mapOwner_eq_nil_iff,
        QuantumCircuit.Party.SourceInventory.map_partyPair_mapOwner,
        QuantumCircuit.Party.SourceInventory.mem_partyPair_mapOwner}
    \leanok
    \uses{def:party_owner_map}
    Let $f\colon P\to Q$ and let $S$ be a list of pair sources on $P$. A pair source
    $t$ occurs in $f_*S$ if and only if $t=f_*s$ for some $s$ in $S$ whose endpoints
    have distinct images. If every vector of $S$ is normalized, so is every vector
    of $f_*S$. The list $f_*S$ is empty if and only if the endpoints of every source
    of $S$ have the same image. The list of unordered endpoint pairs of $f_*S$ is
    obtained from that of $S$ by applying $f$ to each pair and dropping the pairs
    with equal points, in order and with multiplicity; hence an unordered pair
    $\{a,b\}$ occurs in it if and only if $a\neq b$ and $\{a,b\}$ is the image of an
    endpoint pair of $S$ \cite[Theorem~5.2]{OpenAI2026PolynomialPEPS}.
\end{theorem}
\begin{proof}\leanok
    The list $f_*S$ is obtained from $S$ by keeping the moved sources whose image
    endpoints differ, which gives the membership and emptiness statements; a moved
    source has the vector of the source. The list identity follows by induction on
    $S$, splitting on whether the endpoints of the first source have the same image.
\end{proof}

\begin{theorem}[Affected owners]
    \label{thm:party_affected_owner}
    \lean{QuantumCircuit.Party.affectedOwner_eq_iff,
        QuantumCircuit.Party.PairSource.affectedOwner_ne_iff,
        QuantumCircuit.Party.SourceInventory.mem_mapOwner_affectedOwner,
        QuantumCircuit.Party.SourceInventory.map_partyPair_affectedOwner,
        QuantumCircuit.Party.SourceInventory.mapOwner_affectedOwner_eq_nil_iff,
        QuantumCircuit.Party.affectedOwner_true_ne_none,
        QuantumCircuit.Party.uniqueBlockSystemNone}
    \leanok
    \uses{def:party_owner_map, def:sites_block}
    Let $A\subseteq P$. Two parties $p,q$ have the same owner, $\omega_A(p)=\omega_A(q)$,
    if and only if $p=q$ or both are unaffected. The endpoints of a pair source get
    distinct owners if and only if at least one of them is affected. Hence the
    sources occurring in $(\omega_A)_*S$ are exactly the moved sources $(\omega_A)_*s$
    of the sources $s$ of $S$ with an affected endpoint; its list of unordered
    endpoint pairs is the list of images under $\omega_A$ of the endpoint pairs of
    the sources of $S$ with an affected endpoint, in order and with multiplicity; and
    $(\omega_A)_*S$ is empty if and only if both endpoints of every source of $S$
    are unaffected. When every party is affected, no party is sent to $\ast$, and, for
    $P$ finite, the block of $\ast$ in a blocked system along $\omega_P$ has exactly
    one configuration \cite[Theorem~5.2]{OpenAI2026PolynomialPEPS}.
\end{theorem}
\begin{proof}\leanok
    \uses{thm:party_owner_map}
    Split on whether $p$ and $q$ are affected: two affected parties have the same
    owner exactly when they are equal, an affected party and an unaffected one never
    do, and two unaffected parties always do. The endpoints of a source are distinct,
    so they share an owner exactly when both are unaffected. The statements on
    $(\omega_A)_*S$ follow from Theorem~\ref{thm:party_owner_map}, and the list
    identity by induction on $S$. When every party is affected, the preimage of
    $\ast$ is empty and its configurations form the empty product.
\end{proof}

\subsection{Words along a map of parties}
\label{sec:party_coarsen}

Let $P$ be finite and $f\colon P\to Q$. The blocked system $X^f$ has the basis
$X_{f^{-1}(q)}$ at $q\in Q$ (Definition~\ref{def:sites_block}), and its
configurations are identified with those of $X$ by a bijection
$\theta_X\colon X^f_Q\to X_P$. A pair source $g$ with fresh bases $\alpha_g$ and
$\beta_g$ from $X$ to $Y$ identifies the configurations of its output with pairs,
$\gamma_g\colon Y_P\to(\alpha_g\times\beta_g)\times X_P$.

\begin{definition}[A word read on blocks of parties]
    \label{def:party_coarsen}
    \lean{QuantumCircuit.Party.freshBlockEquiv,
        QuantumCircuit.Party.freshBlockEquiv_apply_fst,
        QuantumCircuit.Party.freshBlockEquiv_apply_snd_self,
        QuantumCircuit.Party.freshBlockEquiv_apply_snd_of_ne,
        QuantumCircuit.Party.SourceGate.image_eq,
        QuantumCircuit.Party.SourceGate.coarse,
        QuantumCircuit.Party.SourceGate.toPairSource_coarse,
        QuantumCircuit.Party.Word.coarsen, QuantumCircuit.Party.Word.coarsen_comp,
        QuantumCircuit.Party.Word.coarsen_localOp}
    \leanok
    \uses{def:party_word, def:party_pair_gates, def:sites_block,
        lem:sites_subregion_place_between}
    Let $B$ be a region, $l\in B$, let $Y_v=X_v$ at the sites $v\in B$ other than $l$,
    and let $e\colon Y_l\to\alpha\times X_l$ be a bijection. Then a configuration $z$
    of $Y$ on $B$ is a pair of a fresh part $e(z_l)_1$ and the configuration of $X$ on
    $B$ equal to $e(z_l)_2$ at $l$ and to $z_v$ at $v\neq l$. A pair source $g$ on its
    cell $D$ has the image cell $f(D)=\{f(l),f(r)\}$. If $f(l)\neq f(r)$, the coarse
    source $g^f$ is the pair source on $f(D)$ from $X^f$ to $Y^f$ with endpoints
    $f(l)$ and $f(r)$, the fresh bases and the vector of $g$, its fresh factors located
    in the bases of the blocks by the preceding bijection. The word $w^f$ read on the
    blocks of $f$ is defined by induction on the word $w$:
    \begin{enumerate}
        \item the empty word stays empty, and a composition becomes the
              composition of the coarse words;
        \item a pair source $g$ with $f(l)\neq f(r)$ becomes $g^f$;
        \item a pair source $g$ with $f(l)=f(r)=b$ becomes one local operation at
              $b$, the matrix of $g$ on its cell extended to the block $f^{-1}(b)$ by
              the identity;
        \item a local operation $A$ at $p$ becomes a local operation at $f(p)$, the
              matrix $A$ read on the region $\{p\}$ and extended to $f^{-1}(f(p))$
              by the identity;
        \item local operations $(A_v)_v$ at every party become the local operations
              $\bigotimes_{v\in f^{-1}(b)}A_v$ at the blocks $b$.
    \end{enumerate}
    This is the reading on the affected and exterior systems in the proof of
    \cite[Theorem~5.2]{OpenAI2026PolynomialPEPS}, where sources inside one block
    become local preparations.
\end{definition}

\begin{theorem}[The coarse source on configurations]
    \label{thm:party_coarse_source}
    \lean{QuantumCircuit.Party.SourceGate.configEquiv_coarse}
    \leanok
    \uses{def:party_coarsen}
    Let $g$ be a pair source whose endpoints lie in distinct blocks. For every
    configuration $z$ of $Y^f$, with $(a,x)=\gamma_g(\theta_Y(z))$,
    \begin{align}
        \gamma_{g^f}(z)=\bigl(a,\theta_X^{-1}(x)\bigr).
        \label{eq:party_coarse_source}
    \end{align}
\end{theorem}
\begin{proof}\leanok
    The fresh parts agree by construction. Compare the old parts at a block $b$ and a
    site $v$ of $f^{-1}(b)$. If $b$ is the block of an endpoint and $v$ is that
    endpoint, both are the old part at the endpoint; if $v$ is another site of that
    block, it lies outside the cell, since the other endpoint is in another block, and
    both are $z$ at $v$. If $b$ is neither block, $v$ lies outside the cell and both
    are again $z$ at $v$.
\end{proof}

\begin{theorem}[The operator of a word read on blocks]
    \label{thm:party_op_coarsen}
    \lean{QuantumCircuit.Party.Word.op_coarsen, QuantumCircuit.Party.Word.sum_op_coarsen}
    \leanok
    \uses{def:party_coarsen}
    Let $Q$ be finite. For every word $w$ from $X$ to $Y$ and all configurations $z$
    of $Y^f$ and $z'$ of $X^f$,
    \begin{align}
        [w^f]_{z,z'}=[w]_{\theta_Y(z),\theta_X(z')}.
        \label{eq:party_op_coarsen}
    \end{align}
    For finitely many words $w_\xi$ and coefficients $c_\xi\in\C$, the combination
    $\sum_\xi c_\xi[w_\xi^f]$ is likewise $\sum_\xi c_\xi[w_\xi]$ read along
    $\theta_Y$ and $\theta_X$ \cite[Theorem~5.2]{OpenAI2026PolynomialPEPS}.
\end{theorem}
\begin{proof}\leanok
    \uses{thm:party_coarse_source, lem:sites_place_between_blocks,
        lem:sites_subregion_place_between, thm:sites_block_prod, def:party_word}
    By induction on the word. A pair source inside the block $b$ becomes the placement
    at $b$ of its extended matrix; by Lemma~\ref{lem:sites_place_between_blocks} this
    is, on the configurations of the parties, the extended matrix placed on $f^{-1}(b)$,
    which is the matrix of the source placed on its cell by
    Lemma~\ref{lem:sites_subregion_place_between}. A local operation is treated in the
    same way. The operator of a pair source joining two blocks is the preparation of
    its vector with the output read through $\gamma_g$, and
    (\ref{eq:party_coarse_source}) compares $\gamma_{g^f}$ with $\gamma_g$. Local
    operations at every party give, by Theorem~\ref{thm:sites_block_prod}, the product
    over the blocks of the products over the parties of each block. Products of
    relabelled matrices are relabelled products. The statement on linear combinations
    follows from (\ref{eq:party_op_coarsen}) entrywise.
\end{proof}

\begin{theorem}[Sources and allowedness of a word read on blocks]
    \label{thm:party_coarsen_sources}
    \lean{QuantumCircuit.Party.Word.sources_coarsen,
        QuantumCircuit.Party.Word.isAllowed_coarsen}
    \leanok
    \uses{def:party_coarsen, def:party_owner_map, def:party_word}
    The sources of $w^f$ are the sources of $w$ moved along $f$, $f_*$ of the list of
    sources of $w$. If $w$ is allowed, so is $w^f$.
\end{theorem}
\begin{proof}\leanok
    \uses{lem:sites_subregion_place_between, def:party_pair_gates}
    By induction on the word. A pair source inside one block becomes a local operation
    and is dropped by $f_*$; one joining two blocks becomes $g^f$, which is $f_*g$. The
    lists of a composition are concatenated and $f_*$ preserves concatenation. The
    extension of a matrix to a block does not increase its norm, the matrix of a pair
    source has norm at most that of its vector, and a product of contractions is a
    contraction.
\end{proof}

\subsection{Words restricted to a set of parties}
\label{sec:party_restrict}

\begin{definition}[A word restricted to a set of parties]
    \label{def:party_restrict}
    \lean{QuantumCircuit.Party.SourceGate.restrict,
        QuantumCircuit.Party.SourceGate.regionMatrix_restrict,
        QuantumCircuit.Party.Word.restrict, QuantumCircuit.Party.Word.numGates_restrict}
    \leanok
    \uses{def:party_word, def:party_pair_gates, def:sites_subregion_place}
    Let $S\subseteq P$ be finite. A pair source on a cell $D\subseteq S$ is read on
    the parties of $S$ with the same endpoints, fresh bases and vector; its matrix is
    the matrix of the source, read along the identification of the configurations of
    $D$ as a region of $P$ and as a region of $S$. A word $w$ that uses only the
    parties of $S$ restricts to the word $w|_S$ on the parties of $S$, from
    $(X_s)_{s\in S}$ to $(Y_s)_{s\in S}$, defined by induction on $w$: each pair
    source is read on $S$, a local operation at $p\in S$ is the same local operation
    at $p$, local operations at every party become those at the parties of $S$, and
    the empty word and compositions are kept. The word $w|_S$ has as many gates as
    $w$.
\end{definition}

\begin{theorem}[Sources and allowedness of a restricted word]
    \label{thm:party_restrict_sources}
    \lean{QuantumCircuit.Party.Word.sources_restrict,
        QuantumCircuit.Party.Word.isAllowed_restrict}
    \leanok
    \uses{def:party_restrict, def:party_owner_map}
    Let $w$ use only the parties of $S$. The sources of $w|_S$, moved along the
    inclusion of $S$ in $P$, are the sources of $w$. If $w$ is allowed, so is $w|_S$.
\end{theorem}
\begin{proof}\leanok
    \uses{lem:matrix_l2_contractions}
    By induction on the word. The inclusion is injective, so it keeps every source and
    moves it back to its endpoints in $P$; $f_*$ preserves concatenation. The matrix of
    a restricted source is a relabelling of the matrix of the source, which does not
    increase its norm, and its vector is unchanged.
\end{proof}

\begin{theorem}[A word using a set of parties is placed there]
    \label{thm:party_op_restrict}
    \lean{QuantumCircuit.Party.Word.op_eq_placeBetween_restrict}
    \leanok
    \uses{def:party_restrict, def:sites_place_between}
    Let $P$ be finite and let a word $w$ from $X$ to $Y$ use only the parties of a
    finite set $S$. Then $Y$ agrees with $X$ outside $S$, and
    \begin{align}
        [w]=[w|_S]\otimes1,
        \label{eq:party_op_restrict}
    \end{align}
    the operator of the restricted word placed on $S$ between $X$ and $Y$: the
    registers outside $S$ are untouched \cite[Section~5.1]{OpenAI2026PolynomialPEPS}.
\end{theorem}
\begin{proof}\leanok
    \uses{def:party_word, lem:sites_place_between_products,
        lem:sites_subregion_place_between, thm:party_local_op,
        lem:sites_place_between_algebra}
    By induction on the word. The empty word gives $1=\bigotimes_p1$, and placing the
    product operator of identities on $S$ gives it back, by
    Lemma~\ref{lem:sites_place_between_products}. A pair source on $D\subseteq S$ is
    the placement of its matrix on $D$, and by
    Lemma~\ref{lem:sites_subregion_place_between} this is the placement on $S$ of its
    extension from $D$ to $S$, which is the operator of the restricted source. A local
    operation at $p\in S$ is the product operator of its party family by
    Theorem~\ref{thm:party_local_op}, and so is its restriction; placing the latter on
    $S$ gives the product operator equal to the family on $S$ and to the
    identifications outside $S$, which is the former. Local operations at every party
    are the identifications of equal bases outside $S$, and the same comparison
    applies. A composition follows from
    $([w'|_S]\otimes1)([w|_S]\otimes1)=[w'|_S][w|_S]\otimes1$.
\end{proof}

\subsection{Words on the parties of a cell}
\label{sec:party_on_cell}

\begin{definition}[A word run on the parties of a cell]
    \label{def:party_on_cell}
    \lean{QuantumCircuit.Party.cellEquiv,
        QuantumCircuit.Party.cellEquiv_symm_apply_of_notMem,
        QuantumCircuit.Party.Word.onCell}
    \leanok
    \uses{def:party_coarsen, def:party_gate_frame, def:party_relabel}
    Let $D\subseteq P$ be finite, let $\iota\colon D\to P$ be the inclusion, and let
    $Y$ agree with $X$ outside $D$. The block of a party $p$ along $\iota$ is $\{p\}$
    if $p\in D$ and empty otherwise, and the spectator factor of $p$ is $X_p$ if
    $p\notin D$ and a one-point set otherwise. The basis $Y_p$ is identified with the
    product of the spectator factor and the basis of $p$ in the blocked system
    $(Y_s)_{s\in D}^\iota$; at $p\notin D$ the inverse identification reads the
    spectator factor through $Y_p=X_p$. A word $w$ on the parties of $D$, from
    $(X_s)_{s\in D}$ to $(Y_s)_{s\in D}$, is run inside all parties as the word
    $\widehat w$ from $X$ to $Y$ that first relabels $X$ along these identifications,
    then runs the word $w^\iota$ read on the blocks of $\iota$ beside the spectator
    factors, and finally relabels back to $Y$.
\end{definition}

\begin{theorem}[The operator of a word run on its cell]
    \label{thm:party_op_on_cell}
    \lean{QuantumCircuit.Party.Word.op_onCell,
        QuantumCircuit.Party.Word.isAllowed_onCell,
        QuantumCircuit.Party.Word.sources_onCell}
    \leanok
    \uses{def:party_on_cell, def:sites_place_between, def:party_owner_map}
    Let $P$ be finite. For every word $w$ on the parties of $D$,
    \begin{align}
        [\widehat w]=[w]\otimes1,
        \label{eq:party_op_on_cell}
    \end{align}
    the operator of $w$ placed on $D$ between $X$ and $Y$: an operation on the parties
    of a cell acts as the identity on the other parties
    \cite[Section~5.1]{OpenAI2026PolynomialPEPS}. The word $\widehat w$ is allowed
    when $w$ is, and its sources are those of $w$ moved along $\iota$.
\end{theorem}
\begin{proof}\leanok
    \uses{thm:party_op_frame, thm:party_op_coarsen, def:party_relabel,
        thm:party_frame_allowed, thm:party_coarsen_sources, thm:party_frame_syntax}
    The operator of $\widehat w$ is the product of the two relabellings with
    $[(w^\iota)^W]$, which is $1_{W_P}\otimes[w^\iota]$ by
    Theorem~\ref{thm:party_op_frame} and, by Theorem~\ref{thm:party_op_coarsen},
    $1_{W_P}\otimes[w]$ read on the configurations of $D$. Under the relabellings, the
    block part of a configuration $y$ of $Y$ is its restriction to $D$ and its
    spectator part is its restriction outside $D$, read in $X$. The element at $(y,x)$
    is therefore $[w]_{y|_D,x|_D}$ if $y$ and $x$ agree outside $D$ and $0$ otherwise,
    which is the element of $[w]\otimes1$. Relabellings are allowed and have no
    sources, framing keeps words allowed and keeps their sources by
    Theorems~\ref{thm:party_frame_allowed} and~\ref{thm:party_frame_syntax}, and
    reading on blocks does so up to moving the sources along $\iota$, by
    Theorem~\ref{thm:party_coarsen_sources}.
\end{proof}

\section{Elimination of pair effects}
\label{sec:party_effect_elimination}

This section proves the part of \cite[Lemma~5.1]{OpenAI2026PolynomialPEPS}
that does not refer to parties. A monomial is written as a chain of
contractions and pair effects on finite basis sets, each pair effect is
replaced by the cyclic insertion into its own stack of $m$ pair registers,
and the replaced gate is compared with the original gate followed by a
fixed normalized vector. The pieces of a chain are arbitrary contractions
here; the statement about circuits on parties, whose pieces are local
contractions, is Theorem~\ref{thm:party_pair_effect_elimination}.

Throughout, $X,Y,S,\alpha,\beta,\iota$ denote finite sets. A matrix
$F\in\C^{Y\times X}$, with rows indexed by $Y$, is read as a map
$\C^X\to\C^Y$, and $\lVert F\rVert$ is its operator norm for the Euclidean
norms. We write $\mathbf 1_X$ for the identity matrix of $\C^X$,
$\C^{X\times Y}=\C^X\otimes\C^Y$ for products of basis sets, and
$\alpha^m$ for the set of maps $\{0,\dots,m-1\}\to\alpha$. For
$\eta\in\C^\alpha$, the tensor power $\eta^{\otimes m}\in\C^{\alpha^m}$
has coefficients $\eta^{\otimes m}_f=\prod_{j=0}^{m-1}\eta_{f(j)}$, and
$e_*$ is the unit vector of a one-point set.

\subsection{Preparations, effects and insertion matrices}
\label{sec:party_effect_matrices}

\begin{definition}[Preparation and effect matrices]
    \label{def:party_effect_matrices}
    \lean{Matrix.effectMatrix, Matrix.effectMatrix_apply,
        Matrix.prepMatrix_apply, Matrix.prepMatrix_mul_apply,
        Matrix.prepMatrix_prodComm_mul_apply, Matrix.prepMatrix_eq_submatrix,
        Matrix.reindex_symm_prepMatrix,
        Matrix.toEuclideanLin_prepMatrix_prodComm,
        Matrix.conjTranspose_mul_prepMatrix, Matrix.norm_prepMatrix_le,
        Matrix.norm_prepMatrix, Matrix.effectMatrix_eq_conjTranspose,
        Matrix.conjTranspose_prepMatrix,
        Matrix.toEuclideanLin_effectMatrix_vecKron,
        Matrix.norm_effectMatrix_le}
    \leanok
    \uses{def:construction_idle_gates, lem:kronecker_vectors}
    Let $e:M\to\alpha\times N$ be a bijection of finite sets and
    $v\in\C^\alpha$. The preparation of $v$ along $e$ is the matrix
    $J_e(v)\in\C^{M\times N}$ with $J_e(v)_{y,x}=v_a$ if $e(y)=(a,x)$ and
    $J_e(v)_{y,x}=0$ otherwise
    (Definition~\ref{def:construction_idle_gates}). It maps $|x\rangle$ to
    $v\otimes|x\rangle$, read in the basis $M$ through $e$. For $e$ the
    identity of $\alpha\times N$ it is $|v\rangle\otimes\mathbf 1_N$, and
    for the exchange $N\times\alpha\to\alpha\times N$ it is
    $\mathbf 1_N\otimes|v\rangle$. For $\eta\in\C^\iota$, the effect
    $\langle\eta|\otimes\mathbf 1_S\in\C^{S\times(\iota\times S)}$ has
    entries
    $(\langle\eta|\otimes\mathbf 1_S)_{s,(i,s')}=\overline{\eta_i}\,\delta_{ss'}$.

    These matrices have the following properties.
    \begin{enumerate}
        \item If $e(y)=(a,x)$, then $(J_e(v)A)_{y,z}=v_aA_{x,z}$ for every
              matrix $A$ with rows indexed by $N$. In particular
              $((\mathbf 1_N\otimes|v\rangle)A)_{(x,a),z}=v_aA_{x,z}$, and
              $(\mathbf 1_N\otimes|v\rangle)u=u\otimes v$ for
              $u\in\C^N$.
        \item $J_e(v)$ is $|v\rangle\otimes\mathbf 1_N$ with its rows read
              through $e$. For a bijection $g:L\to M$, the matrix $J_e(v)$
              with its rows read through $g$ is $J_{e\circ g}(v)$.
        \item $J_e(v)^*J_e(v)=\lVert v\rVert^2\,\mathbf 1_N$. Hence
              $\lVert J_e(v)\rVert\le\lVert v\rVert$, with equality when
              $N$ is nonempty.
        \item $\langle\eta|\otimes\mathbf 1_S
                  =(|\eta\rangle\otimes\mathbf 1_S)^*$, and for a bijection
              $e:M\to\iota\times S$ the adjoint $J_e(\eta)^*$ is
              $\langle\eta|\otimes\mathbf 1_S$ with its columns read through
              $e$. The effect sends $h\otimes v$ to
              $\langle\eta,h\rangle\,v$, and
              $\lVert\langle\eta|\otimes\mathbf 1_S\rVert\le1$ when
              $\lVert\eta\rVert=1$.
    \end{enumerate}
\end{definition}

\begin{theorem}[Operator norm of a finite tensor product]
    \label{thm:party_pi_kronecker_norm}
    \lean{Matrix.l2_opNorm_piKronecker_le, Matrix.l2_opNorm_piKronecker_le_one,
        Matrix.piKronecker_eq_reindex_kronecker}
    \leanok
    \uses{def:elementary_finite_tensor_product}
    Let $\nu$ be a finite set and let $K_k\in\C^{I_k\times J_k}$ for
    $k\in\nu$. For every $v\in\nu$, the tensor product $\bigotimes_kK_k$ is
    $K_v\otimes\bigotimes_{j\neq v}K_j$, the index products split at $v$.
    Moreover
    \begin{align}
        \Bigl\lVert\bigotimes_{k\in\nu}K_k\Bigr\rVert
        \le\prod_{k\in\nu}\lVert K_k\rVert,
        \notag
    \end{align}
    so a tensor product of contractions is a contraction.
\end{theorem}

\begin{proof}\leanok
    \uses{thm:detailed_balance_tensor_norm}
    The splitting compares entries: both sides have entry
    $\prod_k(K_k)_{x_k,y_k}$. The bound follows by induction on the
    number of elements of $\nu$. For $\nu$ empty the tensor product is the
    $1\times1$ identity, read along a bijection, and has norm at most $1$.
    Otherwise split off one $v\in\nu$; relabelling rows and columns does not
    increase the operator norm, and
    $\lVert A\otimes B\rVert\le\lVert A\rVert\,\lVert B\rVert$ with the
    induction hypothesis for $\nu\setminus\{v\}$ gives the bound.
\end{proof}

\begin{definition}[Insertion matrices]
    \label{def:party_insertion_matrices}
    \lean{CyclicInsertion.insertAtMatrix, CyclicInsertion.cyclicInsertionMatrix,
        CyclicInsertion.idealInsertionMatrix, CyclicInsertion.insertAtMatrix_apply,
        CyclicInsertion.idealInsertionMatrix_apply,
        CyclicInsertion.cyclicInsertionMatrix_eq_sum,
        CyclicInsertion.insertAtMatrix_eq_insertNth}
    \leanok
    \uses{def:peps_cyclic_insertion, def:party_effect_matrices}
    Let $\eta\in\C^\iota$ and let $\kappa$ be a finite set. For $k\in\kappa$,
    the insertion matrix $I_k\in\C^{\iota^\kappa\times\iota}$ is the matrix
    of the insertion of Definition~\ref{def:peps_cyclic_insertion}, which
    places its input at position $k$ and $\eta$ at every other position:
    \begin{align}
        (I_k)_{f,i}=\delta_{f(k),i}\prod_{j\neq k}\eta_{f(j)}.
        \notag
    \end{align}
    For $m\in\N$, the cyclic and ideal insertion matrices
    $T_m,P_m\in\C^{\iota^m\times\iota}$ are the matrices of the cyclic
    insertion and of $|\eta^{\otimes m}\rangle\langle\eta|$. Thus
    $(P_m)_{f,i}=\eta^{\otimes m}_f\,\overline{\eta_i}$, and for $m\ge1$
    \begin{align}
        T_m=\frac1m\sum_{k=0}^{m-1}I_k.
        \notag
    \end{align}
    For $\kappa=\{0,\dots,n\}$, the insertion $I_k$ is the preparation
    $\mathbf 1_\iota\otimes|\eta^{\otimes n}\rangle$ followed by the
    relabelling that reads $(i,f)\in\iota\times\iota^n$ as the
    configuration with $i$ at position $k$ and the values of $f$, in
    order, at the other positions \cite[proof of
        Lemma~5.1]{OpenAI2026PolynomialPEPS}.
\end{definition}

\begin{theorem}[Norms of the insertion matrices]
    \label{thm:party_insertion_matrix_norms}
    \lean{CyclicInsertion.norm_cyclicInsertionMatrix_sub_ideal_le,
        CyclicInsertion.norm_insertAtMatrix_le_one,
        CyclicInsertion.norm_cyclicInsertionMatrix_le_one}
    \leanok
    \uses{def:party_insertion_matrices}
    Let $\eta\in\C^\iota$ be a unit vector. Then $\lVert I_k\rVert\le1$
    for every position $k$, $\lVert T_m\rVert\le1$ for every $m\in\N$, and
    for $m\ge1$
    \begin{align}
        \lVert T_m-P_m\rVert\le\frac1{\sqrt m}.
        \notag
    \end{align}
    This is the one-effect estimate in the proof of
    \cite[Lemma~5.1]{OpenAI2026PolynomialPEPS}.
\end{theorem}

\begin{proof}\leanok
    \uses{thm:peps_cyclic_insertion_basic, thm:peps_one_effect_estimate}
    The operator norm of a matrix equals the operator norm of the linear map
    it represents. Each insertion of a unit vector is an isometry, and the
    bounds for $T_m$ and $T_m-P_m$ are
    Lemmas~\ref{thm:peps_cyclic_insertion_basic}
    and~\ref{thm:peps_one_effect_estimate}.
\end{proof}

\subsection{Effect chains and the telescoping bound}
\label{sec:party_effect_chains}

\begin{definition}[Effect chain]
    \label{def:party_effect_chain}
    \lean{QuantumCircuit.Party.EffectChain, QuantumCircuit.Party.EffectChain.eval,
        QuantumCircuit.Party.EffectChain.effectCount,
        QuantumCircuit.Party.EffectChain.IsAllowed,
        QuantumCircuit.Party.EffectChain.stackIndex,
        QuantumCircuit.Party.EffectChain.stackVector}
    \leanok
    \uses{def:party_effect_matrices}
    Effect chains from $X$ to $Y$ are defined inductively. A matrix
    $F\in\C^{Y\times X}$ is an effect chain with no effect. If $\alpha,\beta,S$
    are finite sets, $F\in\C^{((\alpha\times\beta)\times S)\times X}$,
    $\eta\in\C^{\alpha\times\beta}$ and $M'$ is an effect chain from $S$ to
    $Y$, then $(F,\eta,M')$ is an effect chain from $X$ to $Y$: it applies
    $F$, which ends with a pair register $\alpha\times\beta$ in front of a
    spectator register $S$, contracts the pair register with $\langle\eta|$,
    and continues with $M'$. An effect chain $M$ with $n$ effects is thus
    given by pieces $F_0,\dots,F_n$ and pair vectors
    $\eta_i\in\C^{\alpha_i\times\beta_i}$ with spectator registers $S_i$,
    and its matrix is
    \begin{align}
        \operatorname{ev}(M)
        =F_n(\langle\eta_n|\otimes\mathbf 1_{S_n})F_{n-1}\cdots
        F_1(\langle\eta_1|\otimes\mathbf 1_{S_1})F_0\in\C^{Y\times X},
        \label{eq:party_effect_chain}
    \end{align}
    defined by $\operatorname{ev}(F)=F$ and
    $\operatorname{ev}(F,\eta,M')
        =\operatorname{ev}(M')(\langle\eta|\otimes\mathbf 1_S)F$. We write
    $n(M)=n$ for its number of effects. The chain is allowed if
    $\lVert F_i\rVert\le1$ and $\lVert\eta_i\rVert=1$ for all $i$
    \cite[Section~5.1]{OpenAI2026PolynomialPEPS}.

    For $m\in\N$, the stack registers of $M$ form the set
    $\Sigma_m(M)=(\alpha_1\times\beta_1)^m\times\bigl(\cdots\times
        ((\alpha_n\times\beta_n)^m\times\{*\})\bigr)$, one stack of $m$ pair
    registers for each effect occurrence, and their ideal content is
    $\Gamma_m(M)=\eta_1^{\otimes m}\otimes\cdots\otimes\eta_n^{\otimes m}
        \otimes e_*$.
\end{definition}

\begin{definition}[Replaced effect chain]
    \label{def:party_effect_chain_replacement}
    \lean{QuantumCircuit.Party.EffectChain.replace,
        QuantumCircuit.Party.EffectChain.replaceTerm}
    \leanok
    \uses{def:party_effect_chain, def:party_insertion_matrices}
    Let $M$ be an effect chain from $X$ to $Y$ and $m\in\N$. The replaced
    chain $R_m(M)\in\C^{(Y\times\Sigma_m(M))\times X}$ is defined by
    $R_m(F)=(\mathbf 1_Y\otimes|e_*\rangle)F$ and
    \begin{align}
        R_m(F,\eta,M')=\Pi\bigl[(\mathbf 1_{(\alpha\times\beta)^m}\otimes
                              R_m(M'))(T_m\otimes\mathbf 1_S)F\bigr],
        \label{eq:party_replace}
    \end{align}
    where $T_m$ is the cyclic insertion matrix of $\eta$ and $\Pi$ reads
    $(\alpha\times\beta)^m\times(Y\times\Sigma_m(M'))$ as
    $Y\times\bigl((\alpha\times\beta)^m\times\Sigma_m(M')\bigr)
        =Y\times\Sigma_m(M)$. Each effect $\langle\eta_i|$ is thus replaced
    by the cyclic insertion into its own stack, and the stack stays idle
    afterwards \cite[proof of Lemma~5.1]{OpenAI2026PolynomialPEPS}.

    For $\kappa:\{0,\dots,n(M)-1\}\to\{0,\dots,m-1\}$, the term
    $R_{m,\kappa}(M)$ is defined in the same way, with the insertion matrix
    $I_{\kappa(0)}$ of $\eta$ in place of $T_m$ and the term of $M'$ at
    $(\kappa(1),\dots,\kappa(n(M)-1))$ in place of $R_m(M')$. It applies no
    pair bra: it composes the pieces $F_i$ with insertions, which are
    preparations of copies of $\eta_i$ followed by relabellings of the
    stack registers (Definition~\ref{def:party_insertion_matrices}).
\end{definition}

\begin{theorem}[Contractivity of effect chains]
    \label{thm:party_effect_chain_contraction}
    \lean{QuantumCircuit.Party.EffectChain.norm_eval_le_one,
        QuantumCircuit.Party.EffectChain.norm_stackVector,
        QuantumCircuit.Party.EffectChain.norm_replace_le_one,
        QuantumCircuit.Party.EffectChain.norm_replaceTerm_le_one}
    \leanok
    \uses{def:party_effect_chain_replacement}
    Let $M$ be an allowed effect chain and $m\in\N$. Then
    $\lVert\operatorname{ev}(M)\rVert\le1$,
    $\lVert\Gamma_m(M)\rVert=1$, $\lVert R_m(M)\rVert\le1$, and
    $\lVert R_{m,\kappa}(M)\rVert\le1$ for every $\kappa$.
\end{theorem}

\begin{proof}\leanok
    \uses{def:party_effect_matrices, thm:party_insertion_matrix_norms,
        lem:kronecker_vectors, thm:detailed_balance_tensor_norm,
        def:construction_idle_gates}
    By induction on the chain. The operator norm is submultiplicative,
    $\lVert\langle\eta|\otimes\mathbf 1_S\rVert\le1$ for a unit vector
    $\eta$, the vectors $\eta^{\otimes m}$ and $e_*$ are unit vectors and
    $\lVert x\otimes y\rVert=\lVert x\rVert\,\lVert y\rVert$, tensoring with
    an identity and relabelling rows do not increase the operator norm, and
    $\lVert T_m\rVert\le1$ and $\lVert I_k\rVert\le1$ by
    Theorem~\ref{thm:party_insertion_matrix_norms}.
\end{proof}

\begin{theorem}[Telescoping bound for one monomial]
    \label{thm:party_effect_chain_replace}
    \lean{QuantumCircuit.Party.EffectChain.norm_replace_sub_le,
        QuantumCircuit.Party.EffectChain.prepMatrix_mul_eval_effect}
    \leanok
    \uses{def:party_effect_chain_replacement}
    Let $m\ge1$ and let $M$ be an allowed effect chain from $X$ to $Y$ with
    $n$ effects. Then
    \begin{align}
        \bigl\lVert R_m(M)-(\mathbf 1_Y\otimes|\Gamma_m(M)\rangle)
        \operatorname{ev}(M)\bigr\rVert\le\frac n{\sqrt m}.
        \label{eq:party_chain_error}
    \end{align}
    For every effect chain $M=(F,\eta,M')$, allowed or not, the ideal
    augmentation replaces the first effect by $P_m$:
    \begin{align}
        (\mathbf 1_Y\otimes|\Gamma_m(M)\rangle)\operatorname{ev}(M)
        =\Pi\bigl[(\mathbf 1_{(\alpha\times\beta)^m}\otimes
                (\mathbf 1_Y\otimes|\Gamma_m(M')\rangle)\operatorname{ev}(M'))
                (P_m\otimes\mathbf 1_S)F\bigr],
        \label{eq:party_chain_ideal}
    \end{align}
    with $\Pi$ as in~(\ref{eq:party_replace}). This is the telescoping step
    in the proof of \cite[Lemma~5.1]{OpenAI2026PolynomialPEPS}.
\end{theorem}

\begin{proof}\leanok
    \uses{def:party_effect_matrices, thm:party_effect_chain_contraction,
        thm:party_insertion_matrix_norms, thm:detailed_balance_tensor_norm,
        def:construction_idle_gates}
    The identity~(\ref{eq:party_chain_ideal}) compares entries, using
    $(P_m)_{f,i}=\eta^{\otimes m}_f\overline{\eta_i}$. The bound is proved
    by induction on the chain; for a chain without effects both sides
    of~(\ref{eq:party_chain_error}) vanish. For $M=(F,\eta,M')$ write
    $A=R_m(M')$ and
    $B=(\mathbf 1_Y\otimes|\Gamma_m(M')\rangle)\operatorname{ev}(M')$. By
    (\ref{eq:party_replace}) and~(\ref{eq:party_chain_ideal}), the
    difference in~(\ref{eq:party_chain_error}) is
    \begin{align}
        \Pi\Bigl[\bigl((\mathbf 1\otimes(A-B))(T_m\otimes\mathbf 1_S)
               +(\mathbf 1\otimes B)((T_m-P_m)\otimes\mathbf 1_S)\bigr)F\Bigr].
        \notag
    \end{align}
    The induction hypothesis gives $\lVert A-B\rVert\le(n-1)/\sqrt m$,
    Theorem~\ref{thm:party_effect_chain_contraction} gives
    $\lVert B\rVert\le1$, and Theorem~\ref{thm:party_insertion_matrix_norms}
    gives $\lVert T_m\rVert\le1$ and
    $\lVert T_m-P_m\rVert\le1/\sqrt m$. With $\lVert F\rVert\le1$ the norm
    is at most $(n-1)/\sqrt m+1/\sqrt m$.
\end{proof}

\begin{theorem}[Expansion of a replaced chain]
    \label{thm:party_effect_chain_expansion}
    \lean{QuantumCircuit.Party.EffectChain.replace_eq_sum}
    \leanok
    \uses{def:party_effect_chain_replacement}
    Let $m\ge1$ and let $M$ be an effect chain with $n$ effects. Then
    \begin{align}
        R_m(M)=\frac1{m^n}\sum_{\kappa}R_{m,\kappa}(M),
        \notag
    \end{align}
    the sum over all $\kappa:\{0,\dots,n-1\}\to\{0,\dots,m-1\}$
    \cite[proof of Lemma~5.1]{OpenAI2026PolynomialPEPS}.
\end{theorem}

\begin{proof}\leanok
    \uses{def:party_insertion_matrices}
    By induction on the chain. For $M=(F,\eta,M')$, substitute
    $T_m=m^{-1}\sum_kI_k$ and the expansion of $R_m(M')$
    into~(\ref{eq:party_replace}), expand by bilinearity of the tensor
    product and linearity of the relabelling $\Pi$, and index the double
    sum over $(k,\kappa')$ by the single map $\kappa$ with $\kappa(0)=k$
    and $(\kappa(1),\dots,\kappa(n-1))=\kappa'$.
\end{proof}

\subsection{Gate expansions and the replaced gate}
\label{sec:party_gate_expansions}

\begin{definition}[Gate expansion, inventory and replaced gate]
    \label{def:party_gate_expansion}
    \lean{QuantumCircuit.Party.GateExpansion,
        QuantumCircuit.Party.GateExpansion.gate,
        QuantumCircuit.Party.GateExpansion.gate_eq_sum,
        QuantumCircuit.Party.GateExpansion.inventory,
        QuantumCircuit.Party.GateExpansion.inventoryVector,
        QuantumCircuit.Party.GateExpansion.replaceGate,
        QuantumCircuit.Party.GateExpansion.termList,
        QuantumCircuit.Party.prodPrependEquiv}
    \leanok
    \uses{def:party_effect_chain_replacement}
    A gate expansion from $X$ to $Y$ is a finite list
    $L=((c_1,M_1),\dots,(c_K,M_K))$ of coefficients $c_\xi\in\C$ and effect
    chains $M_\xi$ from $X$ to $Y$; its gate is
    $G=\sum_\xi c_\xi\operatorname{ev}(M_\xi)$
    \cite[Lemma~5.1]{OpenAI2026PolynomialPEPS}. For $m\in\N$, its inventory
    is $\mathrm{Inv}_m(L)=\Sigma_m(M_1)\times\bigl(\cdots\times
        (\Sigma_m(M_K)\times\{*\})\bigr)$, the stacks of every effect
    occurrence of every monomial, and its ideal content is the vector
    $\Gamma_m=\Gamma_m(M_1)\otimes\cdots\otimes\Gamma_m(M_K)\otimes e_*$.

    The replaced gate
    $G'_m\in\C^{(Y\times\mathrm{Inv}_m(L))\times X}$ is $0$ for the empty
    list, and for a list $(c,M)$ followed by $L$ it is
    \begin{align}
        G'_m=c\,\Psi\bigl[(\mathbf 1_{Y\times\Sigma_m(M)}\otimes
                        |\Gamma_m(L)\rangle)R_m(M)\bigr]
        +\Pi'\bigl[(\mathbf 1_{Y\times\mathrm{Inv}_m(L)}\otimes
                 |\Gamma_m(M)\rangle)G'_m(L)\bigr],
        \label{eq:party_replace_gate}
    \end{align}
    where $\Psi$ reads $(Y\times\Sigma_m(M))\times\mathrm{Inv}_m(L)$ and
    $\Pi'$ reads $(Y\times\mathrm{Inv}_m(L))\times\Sigma_m(M)$ as
    $Y\times(\Sigma_m(M)\times\mathrm{Inv}_m(L))$. Thus in the branch of
    each monomial its effects are replaced by cyclic insertions, the stacks
    of all other monomials are prepared in their ideal contents
    $\Gamma_m(M_{\xi'})$, and every branch has the output register
    $Y\times\mathrm{Inv}_m(L)$ \cite[proof of
        Lemma~5.1]{OpenAI2026PolynomialPEPS}.

    The term list of $L$ is the list of pairs (coefficient, matrix) that is
    empty for the empty list and, for $(c,M)$ followed by $L$, consists of
    the pairs
    $\bigl(c\,m^{-n(M)},\Psi[(\mathbf 1\otimes|\Gamma_m(L)\rangle)
            R_{m,\kappa}(M)]\bigr)$ for all $\kappa$, followed by the pairs
    $\bigl(c',\Pi'[(\mathbf 1\otimes|\Gamma_m(M)\rangle)Q]\bigr)$ for the
    pairs $(c',Q)$ of the term list of $L$.
\end{definition}

\begin{theorem}[Elimination of normalized pair effects without parties]
    \label{thm:party_replace_gate}
    \lean{QuantumCircuit.Party.pairEffectElimination,
        QuantumCircuit.Party.GateExpansion.norm_inventoryVector,
        QuantumCircuit.Party.GateExpansion.norm_replaceGate_sub_le,
        QuantumCircuit.Party.GateExpansion.replaceGate_eq_sum_termList,
        QuantumCircuit.Party.GateExpansion.length_termList_le,
        QuantumCircuit.Party.GateExpansion.sum_norm_coeff_termList,
        QuantumCircuit.Party.GateExpansion.norm_term_le_one}
    \leanok
    \uses{def:party_gate_expansion}
    Let $G=\sum_\xi c_\xi\operatorname{ev}(M_\xi)$ be a gate expansion into
    $K$ allowed effect chains, each with at most $r$ effects, and let
    $m\ge1$. Then $\lVert\Gamma_m\rVert=1$,
    \begin{align}
        \bigl\lVert G'_m-(\mathbf 1_Y\otimes|\Gamma_m\rangle)G\bigr\rVert
        \le\frac r{\sqrt m}\sum_\xi|c_\xi|,
        \label{eq:party_gate_error}
    \end{align}
    and $G'_m=\sum c\,Q$, the sum over the pairs $(c,Q)$ of the term list.
    The term list has at most $Km^r$ pairs, the sum of the absolute values
    of its coefficients is $\sum_\xi|c_\xi|$, and each of its matrices $Q$
    satisfies $\lVert Q\rVert\le1$. The expansion of $G'_m$ holds for every
    gate expansion and every $m$; the normalization of $\Gamma_m$ and the
    contractivity of the terms need only that the chains are allowed.

    These are the error, count and coefficient-sum clauses of
    \cite[Lemma~5.1]{OpenAI2026PolynomialPEPS}, the last with equality where
    the source states an upper bound. The pieces of the chains are arbitrary
    contractions and $G$ need not be a contraction. The clauses of the
    source about parties (the expansion of $G'_m$ uses only local
    contractions and normalized pair sources, and its additional registers
    are owned by the original participating parties) concern circuits on
    parties and are the content of
    Theorem~\ref{thm:party_pair_effect_elimination}; two normalized pair
    vectors combine into one by Lemma~\ref{thm:peps_pair_combine}.
\end{theorem}

\begin{proof}\leanok
    \uses{thm:party_effect_chain_replace, thm:party_effect_chain_expansion,
        thm:party_effect_chain_contraction, def:party_effect_matrices,
        lem:kronecker_vectors}
    Every clause is proved by induction on the list. The vector $\Gamma_m$
    is a tensor product of the unit vectors $\Gamma_m(M_\xi)$ of
    Theorem~\ref{thm:party_effect_chain_contraction}. For the error, the
    two terms of~(\ref{eq:party_replace_gate}) give
    \begin{align}
        G'_m-(\mathbf 1\otimes|\Gamma_m\rangle)G
         & =c\,\Psi\bigl[(\mathbf 1\otimes|\Gamma_m(L)\rangle)
                       \bigl(R_m(M)-(\mathbf 1\otimes|\Gamma_m(M)\rangle)
                       \operatorname{ev}(M)\bigr)\bigr]
        \notag                                                                        \\
         & \quad+\Pi'\bigl[(\mathbf 1\otimes|\Gamma_m(M)\rangle)
                         \bigl(G'_m(L)-(\mathbf 1\otimes|\Gamma_m(L)\rangle)G(L)\bigr)
                         \bigr].
        \notag
    \end{align}
    The preparations of unit vectors are contractions, so
    Theorem~\ref{thm:party_effect_chain_replace} bounds the first term by
    $|c|\,n(M)/\sqrt m\le|c|\,r/\sqrt m$ and the induction hypothesis
    bounds the second. The expansion follows from
    Theorem~\ref{thm:party_effect_chain_expansion} and linearity. A
    monomial with $n\le r$ effects contributes $m^n\le m^r$ pairs, each with
    coefficient of absolute value $|c|m^{-n}$, so their absolute values sum
    to $|c|$. Each matrix of the term list is a preparation of a unit vector
    applied to a term $R_{m,\kappa}(M)$ or to an earlier matrix, hence a
    contraction by Theorem~\ref{thm:party_effect_chain_contraction}.
\end{proof}

\begin{definition}[Stack length for a target error]
    \label{def:party_stack_length}
    \lean{QuantumCircuit.Party.stackLength,
        QuantumCircuit.Party.stackLength_ne_zero,
        QuantumCircuit.Party.stackLength_spec}
    \leanok
    For $r\in\N$ and real numbers $S$ and $\delta>0$, the stack length is
    \begin{align}
        m(r,S,\delta)=\bigl\lceil(rS/\delta)^2\bigr\rceil+1.
        \notag
    \end{align}
    It satisfies $m(r,S,\delta)\ge1$, and if $S\ge0$ then
    $rS/\sqrt{m(r,S,\delta)}\le\delta$, since
    $rS/\delta\le\sqrt{m(r,S,\delta)}$. It depends only on $r$, $S$ and
    $\delta$ \cite[proof of Theorem~5.2]{OpenAI2026PolynomialPEPS}.
\end{definition}

\begin{theorem}[Rescaling the replaced gate]
    \label{thm:party_replace_gate_rescaled}
    \lean{QuantumCircuit.Party.GateExpansion.replaceGate_rescaled,
        QuantumCircuit.Party.GateExpansion.replaceGate_rescaled_stackLength}
    \leanok
    \uses{def:party_gate_expansion, def:party_stack_length}
    Let $G=\sum_\xi c_\xi\operatorname{ev}(M_\xi)$ be a gate expansion into
    allowed effect chains, each with at most $r$ effects, such that
    $\lVert G\rVert\le1$, and write
    $B_m=(\mathbf 1_Y\otimes|\Gamma_m\rangle)G$. If $m\ge1$ and
    $\delta\ge(r/\sqrt m)\sum_\xi|c_\xi|$, then
    \begin{align}
        \lVert G'_m-B_m\rVert               & \le\delta,   &
        \lVert G'_m\rVert                   & \le1+\delta,
        \notag                                               \\
        \lVert(1+\delta)^{-1}G'_m\rVert     & \le1,        &
        \lVert(1+\delta)^{-1}G'_m-B_m\rVert & \le2\delta.
        \notag
    \end{align}
    If $\delta>0$ and $\sum_\xi|c_\xi|\le S$, the stack length
    $m=m(r,S,\delta)$ satisfies $\lVert G'_m-B_m\rVert\le\delta$,
    $\lVert(1+\delta)^{-1}G'_m\rVert\le1$ and
    $\lVert(1+\delta)^{-1}G'_m-B_m\rVert\le2\delta$. This is the reduction
    to source-only gates in the proof of
    \cite[Theorem~5.2]{OpenAI2026PolynomialPEPS}.
\end{theorem}

\begin{proof}\leanok
    \uses{thm:party_replace_gate, lem:rescaled_contraction,
        def:party_stack_length, def:party_effect_matrices}
    The matrix $B_m$ is a contraction, as the preparation of the unit
    vector $\Gamma_m$ applied to $G$. The bound~(\ref{eq:party_gate_error})
    gives $\lVert G'_m-B_m\rVert\le\delta$, and
    Lemma~\ref{lem:rescaled_contraction} gives the other three bounds. For
    $m=m(r,S,\delta)$,
    $(r/\sqrt m)\sum_\xi|c_\xi|\le rS/\sqrt m\le\delta$ by
    Definition~\ref{def:party_stack_length}.
\end{proof}

\section{Stacks of pair registers and the elimination of pair effects}
\label{sec:party_effect_stacks}

This section proves the statements about parties in the elimination of
normalized pair effects \cite[Lemma~5.1]{OpenAI2026PolynomialPEPS}. Every effect
occurrence of a monomial receives its own stack of $m$ pair registers, owned by
the two endpoint parties of the effect; a circuit of pair sources and pair
effects is rewritten, by one recursion on the circuit, as an effect chain whose
replaced terms are words of local contractions and normalized pair sources. For a
weighted gate this gives the elimination theorem on parties, all clauses of the
manuscript's lemma except the combination of sources on one pair of parties.

Throughout, $P$ is a set of parties, a system $X$ is a family of finite basis
sets $X_s$, $s\in P$, with configurations $X_P=\prod_sX_s$, and the operator of a
circuit $c$ from $X$ to $Y$ is the matrix $\mathrm{op}(c)\in\C^{Y_P\times X_P}$
of Definition~\ref{def:party_circuit}. The circuit $c$ followed by $c'$ is
written $c'\circ c$, so $\mathrm{op}(c'\circ c)=\mathrm{op}(c')\,\mathrm{op}(c)$.
For distinct parties $p\ne q$ and finite sets $A,B$, the system $T^{pq}_{A,B}$ has
the basis $A\times T_p$ at $p$, the basis $B\times T_q$ at $q$ and $T_s$ at every
other party (Definition~\ref{def:party_add_pair}); its configurations are
identified with $(A\times B)\times T_P$. For a finite set $\alpha$ and $n\in\N$
write $\alpha^n$ for the maps $\Fin{n}\to\alpha$.

\subsection{Stacks of pair registers}
\label{sec:party_stacks}

\begin{definition}[Circuits beside fresh factors at two parties]
    \label{def:party_lift_pair}
    \lean{QuantumCircuit.Party.PartyCircuit.liftPair,
        QuantumCircuit.Party.PartyCircuit.isAllowed_liftPair,
        QuantumCircuit.Party.PartyCircuit.gateCells_liftPair,
        QuantumCircuit.Party.PartyCircuit.gateCount_liftPair,
        QuantumCircuit.Party.Word.framePairs,
        QuantumCircuit.Party.Word.isAllowed_framePairs,
        QuantumCircuit.Party.Word.sources_framePairs,
        QuantumCircuit.Party.Word.gateCount_framePairs}
    \leanok
    \uses{def:party_circuit, def:party_add_pair, def:party_gate_frame,
        def:party_relabel, def:party_allowed, def:party_gate_counts, def:party_word}
    Let $p\ne q$ and let $A,B$ be finite sets. For a circuit $c$ from $S$ to $T$, the
    circuit $c$ beside the fresh factors $A$ at $p$ and $B$ at $q$ is the circuit
    from $S^{pq}_{A,B}$ to $T^{pq}_{A,B}$ that reads every basis of $S^{pq}_{A,B}$
    as the product of the factor $A$, $B$ or a one-point set with the basis of
    $S$, runs $c$ framed by these factors, and reads the bases back. It has the gate
    cells and the gate count of $c$. For an interpretation compatible with framing,
    with one outcome per gate, and a condition on gates that passes from a gate to
    its framed gates, it is allowed when $c$ is.

    For $n\in\N$ and a word $w$ from $S$ to $T$, the word beside a stack of $n$ pair
    registers on $p,q$ with pair basis $\alpha\times\beta$ is $w$ beside the fresh
    factors $\alpha^n$ at $p$ and $\beta^n$ at $q$. It is allowed when $w$ is, and
    it has the pair sources and the gate count of $w$.
\end{definition}

\begin{theorem}[Operators beside fresh factors]
    \label{thm:party_op_lift_pair}
    \lean{QuantumCircuit.Party.PartyCircuit.op_liftPair,
        QuantumCircuit.Party.Word.op_framePairs,
        QuantumCircuit.Party.Word.op_liftPair_mul_op_sourceAt}
    \leanok
    \uses{def:party_lift_pair, def:party_stack_system, def:party_pair_gates}
    For an interpretation compatible with framing, with one outcome per gate, the
    operator of a circuit $c$ beside the fresh factors $A$ at $p$ and $B$ at $q$ is
    $1_{A\times B}\otimes\mathrm{op}(c)$, read along the identifications of the
    configurations of $S^{pq}_{A,B}$ and $T^{pq}_{A,B}$ with $(A\times B)\times S_P$
    and $(A\times B)\times T_P$. In particular a word beside a stack of $n$ pair
    registers has the operator $1_{(\alpha\times\beta)^n}\otimes\mathrm{op}(w)$, the
    stack configurations read as in Definition~\ref{def:party_stack_system}. If
    $\sigma_\eta^S$ denotes the pair source preparing
    $\eta\in\C^{\alpha\times\beta}$ on fresh factors in front of the bases of $p$ and
    $q$ of $S$, then for every word $c$ from $S$ to $T$
    \begin{align}
        \mathrm{op}(\tilde c)\,\mathrm{op}(\sigma_\eta^S)
         & =\mathrm{op}(\sigma_\eta^T)\,\mathrm{op}(c),
        \notag
    \end{align}
    where $\tilde c$ is $c$ beside the fresh factors $\alpha$ at $p$ and $\beta$ at
    $q$ \cite[Section~5]{OpenAI2026PolynomialPEPS}.
\end{theorem}

\begin{proof}\leanok
    \uses{thm:party_op_frame, def:party_relabel, thm:party_word_op}
    The operator of a composition is the product of the operators; the two
    relabellings are the permutation matrices of the identifications of bases, and
    the framed circuit has the operator $1\otimes\mathrm{op}(c)$ by
    Theorem~\ref{thm:party_op_frame}. Comparing entries gives the first identity,
    and the second is its case $A=\alpha^n$, $B=\beta^n$. The operator of the pair
    source is the preparation $x\mapsto\eta\otimes x$; both sides of the last
    identity have the entry $\eta_{(a,b)}\,\mathrm{op}(c)_{yx}$ at the output
    configuration $((a,b),y)$ and the input $x$.
\end{proof}

\begin{definition}[Stacks of pair registers]
    \label{def:party_stack_system}
    \lean{QuantumCircuit.Party.stackSys, QuantumCircuit.Party.stackSys_of_ne,
        QuantumCircuit.Party.stackConfigEquiv, QuantumCircuit.Party.stackZeroEquiv,
        QuantumCircuit.Party.pairSourceOf, QuantumCircuit.Party.stackSource}
    \leanok
    \uses{def:party_add_pair, def:party_pair_gates, def:sites_config}
    Let $p\ne q$, let $\alpha,\beta$ be finite sets and $n\in\N$. The system $T$ with
    a stack of $n$ pair registers on $p,q$ is $T^{pq}_{\alpha^n,\beta^n}$; at every
    party other than $p$ and $q$ it has the basis of $T$. Its configurations are
    identified with $(\alpha\times\beta)^n\times T_P$, the stack configuration
    $j\mapsto(a_j,b_j)$ collecting the factor $a\in\alpha^n$ at $p$ and
    $b\in\beta^n$ at $q$. For $n=0$ every basis of the stacked system is identified
    with the corresponding basis of $T$. For $\eta\in\C^{\alpha\times\beta}$, the
    pair source with endpoints $p,q$, fresh factors $\alpha,\beta$ and vector $\eta$
    is the occurrence $(p,q,\eta)$; its gate from $T^{pq}_{\alpha^n,\beta^n}$ to
    $T^{pq}_{\alpha^{n+1},\beta^{n+1}}$ on the cell $\{p,q\}$ prepares $\eta$ at
    position $0$ of the stack, reading $a\in\alpha^{n+1}$ as
    $(a_0,(a_1,\dots,a_n))$ and likewise at $q$ \cite[Lemma~5.1]{OpenAI2026PolynomialPEPS}.
\end{definition}

\begin{definition}[Preparation of a stack]
    \label{def:party_prep_word}
    \lean{QuantumCircuit.Party.Word.prepWord,
        QuantumCircuit.Party.Word.isAllowed_prepWord,
        QuantumCircuit.Party.Word.sources_prepWord,
        QuantumCircuit.Party.Word.sourceCount_prepWord,
        QuantumCircuit.Party.Word.usesOnly_prepWord}
    \leanok
    \uses{def:party_stack_system, def:party_word, def:party_relabel}
    The preparation of $n$ copies of $\eta\in\C^{\alpha\times\beta}$ on $p,q$ is the
    word from $T$ to $T^{pq}_{\alpha^n,\beta^n}$ defined by induction on $n$: for
    $n=0$ it is the relabelling of the bases, and for $n+1$ it is the preparation of
    $n$ copies followed by the pair source preparing $\eta$ at position $0$. Its pair
    sources are $n$ copies of $(p,q,\eta)$, so it has $n$ sources; it uses only the
    parties $p$ and $q$, and it is allowed when $\lVert\eta\rVert=1$.
\end{definition}

\begin{theorem}[The prepared stack]
    \label{thm:party_op_prep_word}
    \lean{QuantumCircuit.Party.Word.op_prepWord}
    \leanok
    \uses{def:party_prep_word, def:party_effect_matrices}
    Read along the identification of the configurations of
    $T^{pq}_{\alpha^n,\beta^n}$ with $(\alpha\times\beta)^n\times T_P$, the
    preparation of $n$ copies of $\eta$ has the operator
    $x\mapsto\eta^{\otimes n}\otimes x$
    \cite[Lemma~5.1]{OpenAI2026PolynomialPEPS}.
\end{theorem}

\begin{proof}\leanok
    \uses{def:party_relabel, thm:party_word_op, def:party_effect_matrices}
    By induction on $n$. For $n=0$ the relabelling has the permutation matrix of
    the identification of bases. For $n+1$, the operator of the pair source is the
    preparation of $\eta$ in front of its input, and the entries of the product
    match since
    $\eta^{\otimes(n+1)}(f)=\eta(f_0)\,\eta^{\otimes n}(f_1,\dots,f_n)$.
\end{proof}

\begin{definition}[Insertion words]
    \label{def:party_insertion_word}
    \lean{QuantumCircuit.Party.Word.insEquiv, QuantumCircuit.Party.Word.insWord,
        QuantumCircuit.Party.Word.isAllowed_insWord,
        QuantumCircuit.Party.Word.sources_insWord,
        QuantumCircuit.Party.Word.usesOnly_insWord}
    \leanok
    \uses{def:party_prep_word, def:party_pair_gates, def:party_relabel}
    Let $g$ be a pair effect from $Z$ to $Y$ on the endpoints $p\ne q$, with pair
    basis $\alpha\times\beta$ and vector $\eta$, so that $Z_p=\alpha\times Y_p$ and
    $Z_q=\beta\times Y_q$ up to the identifications of $g$. For $m\ge1$ and
    $k\in\Fin{m}$, the insertion word of $g$ at position $k$ is the word from $Z$ to
    $Y^{pq}_{\alpha^m,\beta^m}$ that prepares $m-1$ copies of $\eta$ on $p,q$ and
    then relabels the bases: at $p$ the basis
    $\alpha^{m-1}\times(\alpha\times Y_p)$ becomes $\alpha^m\times Y_p$, the factor
    $\alpha$ removed by $g$ going to position $k$ of the stack and the prepared
    copies to the other positions in order, and likewise at $q$. It is allowed when
    $\lVert\eta\rVert=1$, its pair sources are $m-1$ copies of $(p,q,\eta)$, and it
    uses only $p$ and $q$ \cite[Lemma~5.1]{OpenAI2026PolynomialPEPS}.
\end{definition}

\begin{theorem}[The insertion word is the insertion map]
    \label{thm:party_op_insertion_word}
    \lean{QuantumCircuit.Party.Word.op_insWord}
    \leanok
    \uses{def:party_insertion_word, def:party_insertion_matrices,
        def:peps_cyclic_insertion}
    Read along the identification $Z_P\cong(\alpha\times\beta)\times Y_P$ of the pair
    effect $g$ and the identification of the configurations of
    $Y^{pq}_{\alpha^m,\beta^m}$ with $(\alpha\times\beta)^m\times Y_P$, the insertion
    word of $g$ at position $k$ has the operator $I_k\otimes1_{Y_P}$, where
    $I_k:\C^{\alpha\times\beta}\to(\C^{\alpha\times\beta})^{\otimes m}$ places its
    input in factor $k$ and $\eta$ in the other factors
    \cite[Lemma~5.1]{OpenAI2026PolynomialPEPS}.
\end{theorem}

\begin{proof}\leanok
    \uses{thm:party_op_prep_word, def:party_relabel, def:party_insertion_matrices}
    The operator is the permutation matrix of the relabelling times the operator
    $x\mapsto\eta^{\otimes(m-1)}\otimes x$ of Theorem~\ref{thm:party_op_prep_word}.
    The insertion $I_k$ is the preparation of $\eta^{\otimes(m-1)}$ behind its input
    followed by the relabelling $(i,f)\mapsto f'$, where $f'$ has $i$ at position $k$
    and $f$ at the other positions in order. Comparing entries at $p$, at $q$ and at
    the other parties proves the identity.
\end{proof}

\subsection{The stacks of a list of effects}
\label{sec:party_effect_stack_lists}

For a pair-source occurrence $e$ with endpoints $p_e\ne q_e$, fresh factors
$\alpha_e,\beta_e$ and vector $\eta_e$, write $T^m_e$ for $T$ with a stack of $m$ pair
registers on $p_e,q_e$ with pair basis $\alpha_e\times\beta_e$.

\begin{definition}[Output stacks of a list of effects]
    \label{def:party_effect_stack}
    \lean{QuantumCircuit.Party.stackApp, QuantumCircuit.Party.stackIdx,
        QuantumCircuit.Party.stackAppEquiv, QuantumCircuit.Party.stackIdxVector,
        QuantumCircuit.Party.norm_stackIdxVector, QuantumCircuit.Party.prepStack,
        QuantumCircuit.Party.frameStack, QuantumCircuit.Party.isAllowed_prepStack,
        QuantumCircuit.Party.sources_prepStack,
        QuantumCircuit.Party.isAllowed_frameStack,
        QuantumCircuit.Party.sources_frameStack}
    \leanok
    \uses{def:party_stack_system, def:party_prep_word, def:party_lift_pair,
        lem:kronecker_vectors}
    Let $E=(e_1,\dots,e_\ell)$ be a list of pair-source occurrences and $m\in\N$. The
    stacks of $E$ in front of $T$ form the system $T^m_E$, defined by $T^m_{()}=T$ and
    $T^m_{eE'}=(T^m_{E'})^m_e$, the first effect outermost. The stack
    configurations of $E$ are
    $J_m(E)=\prod_{i=1}^\ell(\alpha_{e_i}\times\beta_{e_i})^m$, in the order of the
    list, and the configurations of $T^m_E$ are identified with $T_P\times J_m(E)$.
    The ideal content of the stacks is
    \begin{align}
        \Gamma_m(E)
         & =\eta_{e_1}^{\otimes m}\otimes\cdots\otimes\eta_{e_\ell}^{\otimes m}
        \in\C^{J_m(E)},
        \label{eq:party_stack_content}
    \end{align}
    the unit vector of the one-point set for the empty list; it is a unit vector
    when every $\eta_{e_i}$ is. The preparation of the stacks is the word from $T$ to
    $T^m_E$ that prepares, for $E=eE'$, the stacks of $E'$ and then $m$ copies of
    $\eta_e$ on $p_e,q_e$. It is allowed when every $\eta_{e_i}$ is normalized, and
    its pair sources are $m$ copies of $(p_{e_1},q_{e_1},\eta_{e_1})$, then $m$
    copies of the next occurrence, and so on. A word $w$ from $S$ to $T$ beside the
    stacks of $E$ is the word from $S^m_E$ to $T^m_E$ obtained by placing $w$ beside
    the stack of every occurrence in turn; it is allowed when $w$ is and has the pair
    sources of $w$ \cite[Lemma~5.1]{OpenAI2026PolynomialPEPS}.
\end{definition}

\begin{theorem}[Operators of the stacks]
    \label{thm:party_op_effect_stack}
    \lean{QuantumCircuit.Party.op_prepStack, QuantumCircuit.Party.op_frameStack}
    \leanok
    \uses{def:party_effect_stack, def:party_effect_matrices}
    Read along the identification of the configurations of $T^m_E$ with
    $T_P\times J_m(E)$, the preparation of the stacks of $E$ has the operator
    $x\mapsto x\otimes\Gamma_m(E)$, and a word $w$ beside the stacks of $E$ has the
    operator $\mathrm{op}(w)\otimes1_{J_m(E)}$.
\end{theorem}

\begin{proof}\leanok
    \uses{thm:party_op_prep_word, thm:party_op_lift_pair}
    By induction on the list $E$. For $E=eE'$ the operator of the preparation is
    the operator $y\mapsto\eta_e^{\otimes m}\otimes y$ of
    Theorem~\ref{thm:party_op_prep_word} times $x\mapsto x\otimes\Gamma_m(E')$,
    and $\Gamma_m(eE')=\eta_e^{\otimes m}\otimes\Gamma_m(E')$. The operator of a
    word beside the stack of $e$ is $1\otimes\mathrm{op}(w')$ by
    Theorem~\ref{thm:party_op_lift_pair}, with $w'$ the word beside the stacks of
    $E'$; comparing entries gives both identities.
\end{proof}

\begin{theorem}[Ownership and dimensions of the stacks]
    \label{thm:party_stack_ownership}
    \lean{QuantumCircuit.Party.stackApp_of_forall_ne,
        QuantumCircuit.Party.card_stackApp}
    \leanok
    \uses{def:party_effect_stack}
    At a party $s$ that is no endpoint of an occurrence of $E$, the stacked system
    $T^m_E$ has the basis $T_s$. At every party $s$,
    \begin{align}
        \lvert(T^m_E)_s\rvert
         & =\lvert T_s\rvert\prod_{e\in E}
        \lvert\alpha_e\rvert^{m[s=p_e]}\,\lvert\beta_e\rvert^{m[s=q_e]},
        \notag
    \end{align}
    where $[s=p]$ is $1$ if $s=p$ and $0$ otherwise
    \cite[Lemma~5.1]{OpenAI2026PolynomialPEPS}.
\end{theorem}

\begin{proof}\leanok
    \uses{def:party_stack_system, def:party_add_pair}
    By induction on $E$: a stack changes no basis outside its endpoints, and it
    multiplies the dimension at $p_e$ by $\lvert\alpha_e\rvert^m$ and that at $q_e$
    by $\lvert\beta_e\rvert^m$.
\end{proof}

\subsection{The chain form of a circuit of pair sources and pair effects}
\label{sec:party_chain_form}

A circuit of pair sources and pair effects is a circuit whose gates are pair
sources and pair effects (Definition~\ref{def:party_pair_gates}); a monomial is
such a circuit on the parties of a cell. An effect chain
$F_nE_n\cdots E_1F_0$ is as in Definition~\ref{def:party_effect_chain}.

\begin{definition}[The chain form]
    \label{def:party_chain_form}
    \lean{QuantumCircuit.Party.ChainForm, QuantumCircuit.Party.ChainForm.ofWord,
        QuantumCircuit.Party.ChainCont, QuantumCircuit.Party.chainFormAux,
        QuantumCircuit.Party.chainFormAux_comp, QuantumCircuit.Party.chainFormAux_done,
        QuantumCircuit.Party.chainFormAux_localOp,
        QuantumCircuit.Party.chainFormAux_localOps,
        QuantumCircuit.Party.chainFormAux_source,
        QuantumCircuit.Party.chainFormAux_effect,
        QuantumCircuit.Party.PairEffectGate.toChain,
        QuantumCircuit.Party.Word.thenEffect,
        QuantumCircuit.Party.PairEffectGate.op_toChain,
        QuantumCircuit.Party.Monomial.chainForm,
        QuantumCircuit.Party.Monomial.toEffectChain,
        QuantumCircuit.Party.Monomial.chainEffects,
        QuantumCircuit.Party.Monomial.stackEquiv,
        QuantumCircuit.Party.Monomial.replaceWord}
    \leanok
    \uses{def:party_effect_chain, def:party_effect_stack, def:party_insertion_word,
        def:party_lift_pair, def:party_pair_gates, def:party_word}
    A chain form from $X$ to $W$ consists of an effect chain $C$ from $X_P$ to $W_P$,
    a list $E$ of pair-source occurrences (its effects), for every $m$ a bijection
    between the stack registers of $C$ and the stack configurations $J_m(E)$, and for
    every $m$ and every $\kappa:\Fin{n}\to\Fin{m}$, $n$ the number of effects of $C$,
    a word $R_{m,\kappa}$ from $X$ to $W^m_E$. The chain form of a word $w$ has the
    chain $\mathrm{op}(w)$ without effects and the word $w$ itself. A completion rule
    $k$ from $Y$ to $W$ assigns a chain form $k(v)$ from $X$ to $W$ to every word $v$
    from any system $X$ to $Y$.

    For a circuit $c$ of pair sources and pair effects from $Z$ to $Y$, a completion
    rule $k$ from $Y$ to $W$ and a word $w$ from $X$ to $Z$, the chain form
    $\Phi(c,k)(w)$ is defined by recursion on $c$
    \cite[Lemma~5.1]{OpenAI2026PolynomialPEPS}.
    \begin{itemize}
        \item The empty circuit gives $k(w)$; a local operation, local operations at
              every party, or a pair source $h$ give $k(h\circ w)$.
        \item A pair effect $g$ on $p\ne q$ with pair basis $\alpha\times\beta$ and
              vector $\eta$ closes the current piece: let $k(\emptyset)$ be the
              completion of the empty word on $Y$, with chain $C'$, effects $E'$ and
              words $R'$. The chain of $\Phi(c,k)(w)$ is $C'E_gF$, where $F$ is
              $\mathrm{op}(w)$ read along $Z_P\cong(\alpha\times\beta)\times Y_P$ and
              $E_g$ is the effect of $\eta$; its effects are $(p,q,\eta)$ followed by
              $E'$, its stack registers are identified factor by factor, and its
              words are $R_{m,\kappa}=\tilde R'_{m,\kappa'}\circ\iota_{\kappa_0}\circ
                  w$, where $\iota_{\kappa_0}$ is the insertion word of $g$ at
              position $\kappa_0$, $\kappa'$ is $\kappa$ without its first value, and
              $\tilde R'_{m,\kappa'}$ is $R'_{m,\kappa'}$ beside the new stack.
        \item A composition $c'\circ c$ gives $\Phi(c,\Phi(c',k))$.
    \end{itemize}
    One pair effect is a circuit of pair sources and pair effects with the operator
    of the pair effect, and for a word $w$, a pair effect $g$ and a circuit $c$ the
    circuit $c\circ g\circ w$ is the word, then the effect, then $c$.

    The chain form of a monomial $M$ from $X$ to $Y$ is $\Phi(M,k_0)(\emptyset)$,
    where $k_0$ assigns to a word its chain form. Its chain $C_M$, its effects $E_M$
    and its words $R^M_{m,\kappa}$, from $X$ to $Y^m_{E_M}$, are the effect chain,
    the effects and the replacement words of $M$. Composing with the bijection of
    stack registers identifies the configurations of $T^m_{E_M}$ with $T_P$ times
    the stack registers of $C_M$.
\end{definition}

\begin{theorem}[Invariants of the chain form]
    \label{thm:party_chain_form_invariants}
    \lean{QuantumCircuit.Party.eval_chainFormAux,
        QuantumCircuit.Party.isAllowed_chainFormAux,
        QuantumCircuit.Party.isAllowed_replaceWord_chainFormAux,
        QuantumCircuit.Party.stackVector_chainFormAux,
        QuantumCircuit.Party.op_replaceWord_chainFormAux,
        QuantumCircuit.Party.chainFormAux_toChain,
        QuantumCircuit.Party.Monomial.effects_chainFormAux,
        QuantumCircuit.Party.Monomial.effectCount_chainFormAux,
        QuantumCircuit.Party.Monomial.sourceCount_replaceWord_chainFormAux}
    \leanok
    \uses{def:party_chain_form, def:party_allowed, def:party_word,
        def:party_effect_chain_replacement}
    Let $c$ be a circuit of pair sources and pair effects from $Z$ to $Y$ and $k$ a
    completion rule from $Y$ to $W$. Each of the following holds for every word $w$
    into $Z$ when the corresponding hypothesis on $k$ holds for every word $v$ into
    $Y$.
    \begin{enumerate}
        \item If the chain of $k(v)$ evaluates to $R\,\mathrm{op}(v)$ for a fixed
              matrix $R$, the chain of $\Phi(c,k)(w)$ evaluates to
              $R\,\mathrm{op}(c)\,\mathrm{op}(w)$.
        \item Let $c$ be allowed. If the chain of $k(v)$ is allowed for allowed
              $v$, the chain of $\Phi(c,k)(w)$ is allowed for allowed $w$. If every
              word $R_{m,\kappa}$ of $k(v)$ is allowed for allowed $v$, every word
              $R_{m,\kappa}$ of $\Phi(c,k)(w)$ is allowed for allowed $w$.
        \item If the ideal content of the stacks of the chain of $k(v)$, read along
              its bijection of stack registers, is $\Gamma_m$ of its effects, the
              same holds for $\Phi(c,k)(w)$.
        \item If every word $R_{m,\kappa}$ of $k(v)$ has the operator of the term
              of the replaced chain of $k(v)$ with insertion positions $\kappa$, read
              along the identification of the configurations of $W^m_E$ with $W_P$
              times the stack registers, the same holds for $\Phi(c,k)(w)$.
        \item If $c$ is a word, then $\Phi(c,k)(w)=k(w')$ for a word $w'$ with
              $\mathrm{op}(w')=\mathrm{op}(c)\,\mathrm{op}(w)$, with the pair sources
              of $c$ followed by those of $w$, and allowed when $c$ and $w$ are.
        \item Let $c$ be a monomial. If the effects of $k(v)$ are a fixed list
              $E$, the effects of $\Phi(c,k)(w)$ are those of $c$ followed by $E$; if
              the chain of $k(v)$ has $n$ effects, the chain of $\Phi(c,k)(w)$ has
              those of $c$ plus $n$; if every $R_{m,\kappa}$ of $k(v)$ has as many
              sources as $v$ plus $N$, every $R_{m,\kappa}$ of $\Phi(c,k)(w)$ has
              as many sources as $w$, plus those of $c$, plus $m-1$ for every effect
              of $c$, plus $N$.
    \end{enumerate}
\end{theorem}

\begin{proof}\leanok
    \uses{def:party_word, def:party_pair_gates, thm:party_op_contraction,
        def:party_insertion_word, def:party_lift_pair, thm:party_op_insertion_word,
        thm:party_op_lift_pair, def:party_effect_chain,
        def:party_effect_chain_replacement, thm:party_word_op}
    Each statement is proved by induction on $c$, a composition $c'\circ c''$ using
    the statement for $c''$ completed by $\Phi(c',k)$. A local operation or a pair
    source joins the current word, whose operator is unchanged when it is read as a
    circuit of pair sources and pair effects. For a pair effect $g$, the effect
    matrix $E_g=\bra{\eta}\otimes1$ is the operator of $g$, so the chain evaluates
    to $R\,E_g\,\mathrm{op}(w)$; the piece $F$ is a contraction by
    Theorem~\ref{thm:party_op_contraction}; the insertion word and the word beside
    the stack are allowed, have $m-1$ and $N$ further sources, and have the
    operators $I_{\kappa_0}\otimes1$ and $1\otimes\mathrm{op}(R'_{m,\kappa'})$ by
    Theorems~\ref{thm:party_op_insertion_word} and~\ref{thm:party_op_lift_pair},
    which is the recursive formula of the replaced term after regrouping the stack
    configurations.
\end{proof}

\begin{theorem}[Monomials as effect chains]
    \label{thm:party_monomial_chain}
    \lean{QuantumCircuit.Party.Monomial.eval_toEffectChain,
        QuantumCircuit.Party.Monomial.isAllowed_toEffectChain,
        QuantumCircuit.Party.Monomial.chainEffects_eq,
        QuantumCircuit.Party.Monomial.effectCount_toEffectChain,
        QuantumCircuit.Party.Monomial.effectParties_eq,
        QuantumCircuit.Party.Monomial.norm_vector_of_mem_chainEffects,
        QuantumCircuit.Party.Monomial.toEffectChain_toChain,
        QuantumCircuit.Party.Monomial.toEffectChain_thenEffect,
        QuantumCircuit.Party.Monomial.chainEffects_thenEffect,
        QuantumCircuit.Party.Monomial.isAllowed_thenEffect}
    \leanok
    \uses{def:party_chain_form, def:party_allowed, def:party_pair_gates}
    For a monomial $M$ from $X$ to $Y$, the chain $C_M$ evaluates to $\mathrm{op}(M)$
    and has as many effects as $M$; the effects $E_M$ of the chain form are the pair
    effects of $M$ in chronological order, so their endpoints are those of the pair
    effects of $M$. If $M$ is allowed, then $C_M$ is allowed and every vector of
    $E_M$ is normalized. A word $w$ read as a monomial has the chain
    $\mathrm{op}(w)$. For a word $w$ from $X$ to $Z$, a pair effect $g$ from $Z$ to
    $Y$ with vector $\eta$ and a monomial $c$ from $Y$ to $W$, the chain of
    $c\circ g\circ w$ is $C_cE_gF$ with $F$ the operator of $w$ read along the
    configurations of $g$, its effects are $g$ followed by $E_c$, and the circuit
    $c\circ g\circ w$ is allowed when $w$ and $c$ are allowed and
    $\lVert\eta\rVert=1$
    \cite[Lemma~5.1]{OpenAI2026PolynomialPEPS}.
\end{theorem}

\begin{proof}\leanok
    \uses{thm:party_chain_form_invariants, thm:party_op_contraction,
        def:party_pair_gates}
    Apply Theorem~\ref{thm:party_chain_form_invariants} with the completion $k_0$,
    for which the chain of $k_0(v)$ is $\mathrm{op}(v)$, has no effect, and is
    allowed for allowed $v$ by Theorem~\ref{thm:party_op_contraction}: the first
    statement with $R=1$, the second, and the sixth with $E$ empty and $n=0$.
    Normalization of the effects of an allowed monomial is part of allowedness.
    The chains of a word and of $c\circ g\circ w$ follow from the fifth statement;
    the operator of a pair effect has norm at most $\lVert\eta\rVert$.
\end{proof}

\begin{theorem}[The replaced terms of a monomial are source-only]
    \label{thm:party_monomial_replaced_terms}
    \lean{QuantumCircuit.Party.Monomial.isAllowed_replaceWord,
        QuantumCircuit.Party.Monomial.op_replaceWord,
        QuantumCircuit.Party.Monomial.sourceCount_replaceWord,
        QuantumCircuit.Party.Monomial.isAllowed_prepStack,
        QuantumCircuit.Party.Monomial.op_prepStack,
        QuantumCircuit.Party.Monomial.stackApp_of_notMem_effectParties}
    \leanok
    \uses{def:party_chain_form, def:party_effect_chain_replacement,
        def:party_effect_stack}
    Let $M$ be a monomial from $X$ to $Y$ with $n$ pair effects and $s_M$ pair
    sources, let $m\in\N$ and $\kappa:\Fin{n}\to\Fin{m}$. Read along the
    identification of the configurations of $Y^m_{E_M}$ with $Y_P$ times the stack
    registers of $C_M$, the word $R^M_{m,\kappa}$ has as operator the term of the
    replaced chain $C_M$ with insertion positions $\kappa$. It has
    $s_M+n(m-1)$ pair sources, and it is allowed when $M$ is. The preparation of the
    stacks of $E_M$ in front of $T$ has the operator $x\mapsto x\otimes\Gamma$,
    where $\Gamma$ is the ideal content of the stacks of $C_M$; it is allowed when
    $M$ is. At a party that is no endpoint of a pair effect of $M$, the system
    $T^m_{E_M}$ has the basis of $T$ \cite[Lemma~5.1]{OpenAI2026PolynomialPEPS}.
\end{theorem}

\begin{proof}\leanok
    \uses{thm:party_chain_form_invariants, thm:party_monomial_chain,
        thm:party_op_effect_stack, thm:party_stack_ownership, def:party_effect_stack}
    For the completion $k_0$, the word of a chain without effects is $v$ itself, and
    the replaced term of the chain $\mathrm{op}(v)$ is $\mathrm{op}(v)$ followed by
    the preparation of the unit vector of the one-point set. The fourth, second and
    sixth statements of Theorem~\ref{thm:party_chain_form_invariants} give the
    operator, allowedness and the source count. The operator of the preparation
    follows from Theorem~\ref{thm:party_op_effect_stack} and the third statement,
    its allowedness from the normalization of $E_M$
    (Theorem~\ref{thm:party_monomial_chain}), and the ownership from
    Theorem~\ref{thm:party_stack_ownership}, since the endpoints of $E_M$ are those
    of the pair effects of $M$.
\end{proof}

\subsection{Elimination of pair effects for a weighted gate}
\label{sec:party_weighted_elimination}

Let $D$ be a cell and $L=((c_1,M_1),\dots,(c_K,M_K))$ a list of coefficients
$c_\xi\in\C$ and monomials $M_\xi$ from $A$ to $B$ on the parties of $D$, such as the
terms of a weighted gate (Definition~\ref{def:party_weighted_gate}). Gate
expansions, their inventory, their ideal vector, the replaced gate and its term
list are those of Definition~\ref{def:party_gate_expansion}.

\begin{definition}[The gate expansion of a weighted gate]
    \label{def:party_weighted_expansion}
    \lean{QuantumCircuit.Party.WeightedGate.toGate,
        QuantumCircuit.Party.WeightedGate.length_toGate,
        QuantumCircuit.Party.WeightedGate.map_norm_toGate,
        QuantumCircuit.Party.WeightedGate.mem_toGate,
        QuantumCircuit.Party.WeightedGate.gate_toGate,
        QuantumCircuit.Party.WeightedGate.ev_eq_gate_toGate,
        QuantumCircuit.Party.WeightedGate.isAllowed_toGate}
    \leanok
    \uses{def:party_gate_expansion, def:party_chain_form, def:party_weighted_gate}
    The gate expansion $\mathcal G_L$ of $L$ is the list of pairs $(c_\xi,C_{M_\xi})$.
    It has $K$ terms, with the absolute values $\lvert c_\xi\rvert$, every term is
    $(c,C_M)$ for a term $(c,M)$ of $L$, and its gate is
    $G=\sum_\xi c_\xi\,\mathrm{op}(M_\xi)$; in particular the matrix of a weighted
    gate is the gate of the expansion of its terms. When every $M_\xi$ is allowed,
    every chain of $\mathcal G_L$ is allowed.
\end{definition}

\begin{definition}[The output system of the replaced gate and its words]
    \label{def:party_replaced_gate_output}
    \lean{QuantumCircuit.Party.WeightedGate.gateOut,
        QuantumCircuit.Party.WeightedGate.gateEquiv,
        QuantumCircuit.Party.WeightedGate.prepGate,
        QuantumCircuit.Party.WeightedGate.isAllowed_prepGate,
        QuantumCircuit.Party.WeightedGate.wordList,
        QuantumCircuit.Party.WeightedGate.isAllowed_wordList}
    \leanok
    \uses{def:party_weighted_expansion, def:party_chain_form, def:party_effect_stack,
        def:party_gate_expansion}
    Let $m\in\N$. The output system $B^m_L$ of the replaced gate is $B$ for the empty
    list and $(B^m_{L'})^m_{E_M}$ for $L=(c,M)L'$: the stacks of all effect
    occurrences of all monomials in front of $B$. Its configurations are identified
    with $B_P$ times the inventory of $\mathcal G_L$, one monomial at a time by the
    identifications of Definition~\ref{def:party_chain_form}. The preparation word
    $\pi_L$ from $B$ to $B^m_L$ is empty for the empty list and, for $L=(c,M)L'$, is
    $\pi_{L'}$ followed by the preparation of the stacks of $E_M$; it is allowed when
    every monomial is. The word list of the empty list is empty. The word list of
    $L=(c,M)L'$ consists of the pairs of the first kind
    $(c\,m^{-n_M},\tilde\pi_{L'}\circ R^M_{m,\kappa})$, one for every
    $\kappa:\Fin{n_M}\to\Fin{m}$, with $n_M$ the number of effects of $C_M$ and
    $\tilde\pi_{L'}$ the word $\pi_{L'}$ beside the stacks of $E_M$, followed by the
    pairs of the second kind $(c',\rho\circ w')$, one for every term $(c',w')$ of the
    word list of $L'$, with $\rho$ the preparation of the stacks of $E_M$ in front of
    $B^m_{L'}$. When every monomial is allowed, every word of the list is allowed
    \cite[Lemma~5.1]{OpenAI2026PolynomialPEPS}.
\end{definition}

\begin{theorem}[The replaced gate on parties]
    \label{thm:party_replaced_gate_words}
    \lean{QuantumCircuit.Party.WeightedGate.op_prepGate,
        QuantumCircuit.Party.WeightedGate.termList_toGate,
        QuantumCircuit.Party.WeightedGate.isSourceOnly_termList,
        QuantumCircuit.Party.WeightedGate.gateOut_of_forall_ne,
        QuantumCircuit.Party.WeightedGate.card_gateOut}
    \leanok
    \uses{def:party_replaced_gate_output, def:party_gate_expansion, def:party_word}
    Read along the identification of the configurations of $B^m_L$ with $B_P$ times
    the inventory, the preparation word $\pi_L$ has the operator
    $x\mapsto x\otimes\Gamma_m$, with $\Gamma_m$ the ideal vector of $\mathcal G_L$,
    and the term list of the replaced gate of $\mathcal G_L$ is, term by term, the
    list of pairs $(c,\mathrm{op}(w))$ for the terms $(c,w)$ of the word list. When
    every monomial is allowed, every term of the term list, read on $B^m_L$, is
    source-only. At a party $s\in D$ that is no endpoint of a pair effect of any
    monomial of $L$, the basis of $B^m_L$ is $B_s$, and at every party $s$
    \begin{align}
        \lvert(B^m_L)_s\rvert
         & =\lvert B_s\rvert\prod_{\xi=1}^K\prod_{e}
        \lvert\alpha_e\rvert^{m[s=p_e]}\,\lvert\beta_e\rvert^{m[s=q_e]},
        \notag
    \end{align}
    the inner product over the pair effects $e$ of $M_\xi$
    \cite[Lemma~5.1]{OpenAI2026PolynomialPEPS}.
\end{theorem}

\begin{proof}\leanok
    \uses{thm:party_monomial_replaced_terms, thm:party_op_effect_stack,
        thm:party_stack_ownership, thm:party_monomial_chain,
        def:party_replaced_gate_output}
    By induction on $L$. For $L=(c,M)L'$ the operator of $\pi_L$ is the product of
    the operator $x\mapsto x\otimes\Gamma_{m}(C_M)$ of the preparation of the stacks
    of $E_M$ (Theorem~\ref{thm:party_monomial_replaced_terms}) and that of
    $\pi_{L'}$. A pair of the first kind has the
    operator $(\mathrm{op}(\pi_{L'})\otimes1)\,\mathrm{op}(R^M_{m,\kappa})$ by
    Theorem~\ref{thm:party_op_effect_stack}, which is the replaced term of $C_M$
    followed by the preparation of the ideal vector of $\mathcal G_{L'}$, the term of
    the term list; a pair of the second kind is a term of $L'$ followed by the
    preparation of the ideal content of the stacks of $C_M$. Allowed words give
    source-only terms. Ownership and dimensions follow from
    Theorems~\ref{thm:party_monomial_replaced_terms}
    and~\ref{thm:party_stack_ownership}, since the effects $E_M$ are the pair
    effects of $M$.
\end{proof}

\begin{theorem}[Elimination of normalized pair effects on parties]
    \label{thm:party_pair_effect_elimination}
    \lean{QuantumCircuit.Party.WeightedGate.partyPairEffectElimination}
    \leanok
    \uses{def:party_replaced_gate_output, def:party_weighted_expansion,
        def:party_gate_expansion, def:party_word}
    Let $m\ge1$ and $r\in\N$, and let every monomial $M_\xi$ of $L$ be allowed and have
    at most $r$ pair effects. Let $G=\sum_\xi c_\xi\,\mathrm{op}(M_\xi)$ be the gate
    of $\mathcal G_L$, let $G'_m$
    be the replaced gate of $\mathcal G_L$ and $\Gamma_m$ its ideal vector. Then
    $\lVert\Gamma_m\rVert=1$, the allowed word $\pi_L$ prepares $\Gamma_m$, read
    along the identification of the configurations of $B^m_L$ with $B_P$ times the
    inventory, and
    \begin{align}
        \lVert G'_m-G\otimes\ket{\Gamma_m}\rVert
         & \le\frac{r}{\sqrt m}\sum_{\xi=1}^K\lvert c_\xi\rvert.
        \label{eq:party_elimination_error}
    \end{align}
    Read along the same identification, $G'_m=\sum_{(c,w)}c\,\mathrm{op}(w)$ over
    the word list, every word of the list is allowed, and every term of the term
    list is source-only. The word list has at most $Km^r$ terms and the absolute sum
    $\sum_{(c,w)}\lvert c\rvert$ of its coefficients equals
    $\sum_\xi\lvert c_\xi\rvert$. At every party of $D$ that is no endpoint of a
    pair effect of any monomial, the output system $B^m_L$ of $G'_m$ has the basis of
    $B$.
\end{theorem}

This covers \cite[Lemma~5.1]{OpenAI2026PolynomialPEPS} on parties: the
error~(\ref{eq:party_elimination_error}), the expansion of the replaced gate into
local contractions and normalized pair sources, the bound $Km^r$ on the number of
monomials, the coefficient sum, here with equality, and the ownership of the
additional registers by the endpoint parties of the eliminated effects, which are
parties of the cell. The hypothesis that $G$ is a contraction is not used. The last
clause of the manuscript's lemma, that all sources on one pair of parties may be
combined into one normalized pair source, is not part of this theorem; for two pair
vectors the norm of the combined vector is Lemma~\ref{thm:peps_pair_combine}.

\begin{proof}\leanok
    \uses{thm:party_replace_gate, def:party_weighted_expansion,
        thm:party_monomial_chain, thm:party_replaced_gate_words,
        def:party_replaced_gate_output}
    Every chain of $\mathcal G_L$ is allowed and has as many effects as its monomial
    (Theorem~\ref{thm:party_monomial_chain}), so Theorem~\ref{thm:party_replace_gate}
    applies to $\mathcal G_L$: it gives the norm of $\Gamma_m$, the
    error~(\ref{eq:party_elimination_error}), the expansion of $G'_m$ into the term
    list, the bound on its length and the coefficient sum, since $\mathcal G_L$ has
    $K$ terms with the coefficients of $L$. Theorem~\ref{thm:party_replaced_gate_words}
    rewrites the term list as the word list and gives the preparation of
    $\Gamma_m$, source-only terms and ownership, and
    Definition~\ref{def:party_replaced_gate_output} gives that the words are
    allowed.
\end{proof}

\begin{definition}[Rescaling a weighted gate]
    \label{def:party_rescale}
    \lean{QuantumCircuit.Party.WeightedGate.rescale,
        QuantumCircuit.Party.WeightedGate.monomials_rescale,
        QuantumCircuit.Party.WeightedGate.sum_norm_rescale_le,
        QuantumCircuit.Party.WeightedGate.ev_rescale,
        QuantumCircuit.Party.WeightedGate.gate_toGate_rescale,
        QuantumCircuit.Party.WeightedGate.isAllowed_rescale}
    \leanok
    \uses{def:party_weighted_gate, def:party_weighted_expansion,
        lem:rescaled_contraction}
    For $\delta\in\R$, the rescaling of a weighted gate multiplies every coefficient
    by $(1+\delta)^{-1}$, with $0^{-1}=0$, and keeps the ordered list of monomials.
    Its matrix, and the gate of the expansion of its terms, are $(1+\delta)^{-1}$
    times those of the gate; it is allowed when the gate is, and for $\delta\ge0$
    its absolute coefficient sum is at most that of the gate. By
    Lemma~\ref{lem:rescaled_contraction}, a gate of norm at most $1+\delta$ thus
    becomes a contraction, as in the proof of
    \cite[Theorem~5.2]{OpenAI2026PolynomialPEPS}.
\end{definition}

\section{Ownership sheets, births, deaths and exchanges}
\label{sec:party_sheets}

An ownership sheet assigns an owner to every site of a finite set $\Lambda$, and the register
$\C^q$ of a site is held by its owner. This section implements the homogeneous births and
deaths of \cite[Lemma~6.5]{OpenAI2026PolynomialPEPS} and the two-sheet exchange of
\cite[Lemma~6.6]{OpenAI2026PolynomialPEPS} by words and monomials of the party calculus, for
given splitting data, and assembles the memory of several sheets used in
\cite[Proposition~7.1]{OpenAI2026PolynomialPEPS}: lifts of changes of one or two sheets to the
whole stack, bounded changes, and the private product initialization. Sheets carry no encoded
holes, so the hole parts of the sets of the two lemmas are empty.

\subsection{Framing, relabelling and sequencing of original circuits}

\begin{definition}[Framing, relabelling and sequencing of original circuits]
    \label{def:party_original_frame_sequence}
    \lean{QuantumCircuit.Party.prodVec, QuantumCircuit.Party.prodVec_apply,
        QuantumCircuit.Party.norm_prodVec, QuantumCircuit.Party.prodVec_sub_left,
        QuantumCircuit.Party.prodVec_sub_right, QuantumCircuit.Party.OriginalGate.frame,
        QuantumCircuit.Party.OriginalGate.toWeightedGate_frame,
        QuantumCircuit.Party.OriginalCircuit.frame,
        QuantumCircuit.Party.OriginalCircuit.toEffectCircuit_frame,
        QuantumCircuit.Party.OriginalCircuit.nonprivateCount_frame,
        QuantumCircuit.Party.OriginalCircuit.participation_frame,
        QuantumCircuit.Party.OriginalCircuit.frameLocationsEquiv,
        QuantumCircuit.Party.OriginalCircuit.participants_frameLocationsEquiv,
        QuantumCircuit.Party.OriginalCircuit.forall_participants_frame,
        QuantumCircuit.Party.OriginalCircuit.isExpansionBounded_frame,
        QuantumCircuit.Party.OriginalCircuit.isMonomialBounded_frame,
        QuantumCircuit.Party.OriginalCircuit.isSourceBounded_frame,
        QuantumCircuit.Party.EffectCircuit.isExpansionBounded_frame,
        QuantumCircuit.Party.EffectCircuit.isMonomialBounded_frame,
        QuantumCircuit.Party.EffectCircuit.isSourceBounded_frame,
        QuantumCircuit.Party.OriginalCircuit.relabel,
        QuantumCircuit.Party.OriginalCircuit.toPartyCircuit_relabel,
        QuantumCircuit.Party.OriginalCircuit.nonprivateCount_relabel,
        QuantumCircuit.Party.OriginalCircuit.participation_relabel,
        QuantumCircuit.Party.OriginalCircuit.sequence,
        QuantumCircuit.Party.OriginalCircuit.sequence_zero,
        QuantumCircuit.Party.OriginalCircuit.sequence_succ,
        QuantumCircuit.Party.OriginalCircuit.op_sequence_zero,
        QuantumCircuit.Party.OriginalCircuit.op_sequence_succ,
        QuantumCircuit.Party.OriginalCircuit.isExpansionBounded_sequence,
        QuantumCircuit.Party.OriginalCircuit.isMonomialBounded_sequence,
        QuantumCircuit.Party.OriginalCircuit.isSourceBounded_sequence,
        QuantumCircuit.Party.OriginalCircuit.nonprivateCount_sequence,
        QuantumCircuit.Party.OriginalCircuit.forall_participants_sequence,
        QuantumCircuit.Party.OriginalCircuit.participation_sequence}
    \leanok
    \uses{def:party_original_circuit, def:party_gate_frame, def:party_relabel,
        def:party_gate_counts, def:party_weighted_gate, lem:kronecker_vectors}
    Let $X$, $Y$ and $W$ be systems on the parties. In this section $\operatorname{op}(w)$
    denotes the operator of an original circuit, word or monomial $w$, the letter $T$ being kept
    for a region of sites. For an original circuit, its number of
    nonprivate gates is its number of gates, every gate costing $1$ and local contractions
    nothing, and the participation of a party is the number of gates whose cell contains it.
    It satisfies the expansion bound $(r,S)$ if every monomial of every gate expansion has at
    most $r$ pair effects and the absolute values of the coefficients of every gate expansion
    sum to at most $S$, the monomial bound $K$ if every gate expansion has at most $K$
    monomials, and the pair-source bound $r$ if every monomial of every gate expansion has at
    most $r$ pair sources \cite[Theorem~5.2]{OpenAI2026PolynomialPEPS}.

    \emph{Product vectors.} The product system $W\times X$ has the basis $W_p\times X_p$ at the
    party $p$, and its configurations are the pairs of configurations of $W$ and $X$. For
    $x\in\mathcal H_W$ and $y\in\mathcal H_X$ the product vector
    $x\otimes y\in\mathcal H_{W\times X}$ has the coefficient $x(\sigma)y(\tau)$ at the
    configuration paired with $(\sigma,\tau)$. Thus
    $\lVert x\otimes y\rVert=\lVert x\rVert\,\lVert y\rVert$,
    $x\otimes y-x'\otimes y=(x-x')\otimes y$ and $x\otimes y-x\otimes y'=x\otimes(y-y')$.

    \emph{Framing.} For an original circuit $w$ from $X$ to $Y$, the framed circuit $w^W$ from
    $W\times X$ to $W\times Y$ runs $w$ beside the spectator factor $W_p$ placed in front of the
    basis of every party $p$: every original gate is replaced by its framed gate expansion, and
    every local contraction $A$ at a party $p$ by $1_{W_p}\otimes A$. The effect circuit of
    $w^W$ is the framed effect circuit of $w$. Framing keeps the number of nonprivate gates,
    the participation of every party, the gate occurrences together with their cells, hence
    every condition on the participants of the gates, and the expansion, monomial and
    pair-source bounds, of original circuits and of effect circuits alike.

    \emph{Relabelling.} For bijections $e_p:Y_p\to X_p$, the relabelling along $e$ is the
    original circuit from $X$ to $Y$ made of one local operation at every party, the matrix of
    $e_p$; it is the relabelling of party circuits, has no nonprivate gate, and no party
    participates in it.

    \emph{Sequencing.} For systems $a_0,a_1,\dots$ and original circuits $c_t$ from $a_t$ to
    $a_{t+1}$, the sequence of length $t$ is the empty circuit for $t=0$ and $c_0$ followed by
    $c_1,\dots,c_{t-1}$ in general, so the sequence of length $t+1$ is that of length $t$
    followed by $c_t$. Its operator is $\operatorname{op}(c_{t-1})\cdots\operatorname{op}(c_0)$,
    the identity for $t=0$. Its number of nonprivate gates and the participation of every party
    are the sums of those of the stages $c_s$, $s<t$; an expansion, monomial or pair-source
    bound satisfied by every stage $c_s$ with $s<t$ is satisfied by the sequence, and so is
    every condition on the participants of the gates.
\end{definition}

\begin{theorem}[Framed circuits and relabellings on product vectors]
    \label{thm:party_frame_prod_vec}
    \lean{QuantumCircuit.Party.OriginalCircuit.op_frame,
        QuantumCircuit.Party.OriginalGate.isCompatible_frame,
        QuantumCircuit.Party.OriginalCircuit.toEuclideanLin_op_frame_prodVec,
        QuantumCircuit.Party.OriginalCircuit.norm_op_frame_prodVec_sub_le,
        QuantumCircuit.Party.OriginalCircuit.op_relabel,
        QuantumCircuit.Party.OriginalCircuit.toEuclideanLin_op_relabel,
        QuantumCircuit.Party.OriginalCircuit.toEuclideanLin_op_relabel_prodComm_prodVec}
    \leanok
    \uses{def:party_original_frame_sequence}
    The operator of a framed original gate is the identity on the spectator factors tensored
    with the operator of the gate. For an original circuit $w$ from $X$ to $Y$ and a system $W$,
    read along the pairing of configurations,
    $\operatorname{op}(w^W)=1_{\mathcal H_W}\otimes\operatorname{op}(w)$. Hence
    $\operatorname{op}(w^W)(x\otimes y)=x\otimes\operatorname{op}(w)y$ for
    $x\in\mathcal H_W$ and $y\in\mathcal H_X$, and if $\lVert x\rVert\leq1$ then
    \begin{align}
        \lVert\operatorname{op}(w^W)(x\otimes v)-x\otimes v'\rVert
        \leq\lVert\operatorname{op}(w)v-v'\rVert
    \end{align}
    for all $v\in\mathcal H_X$ and $v'\in\mathcal H_Y$. The operator of the relabelling along
    bijections $e_p:Y_p\to X_p$ is the matrix of the bijection $\tau\mapsto(e_p(\tau_p))_p$ of
    configurations; it maps a vector $v\in\mathcal H_X$ to the vector
    $\tau\mapsto v((e_p(\tau_p))_p)$ of $\mathcal H_Y$. In particular the relabelling exchanging
    the two factors of $X\times Y$ maps $x\otimes y$ to $y\otimes x$.
\end{theorem}
\begin{proof}\leanok
    \uses{thm:party_op_frame, def:party_relabel, def:party_weighted_gate}
    The compatibility of framed original gates is that of their framed gate expansions, and
    Theorem~\ref{thm:party_op_frame} for party circuits then gives
    $\operatorname{op}(w^W)=1\otimes\operatorname{op}(w)$. Applied to $x\otimes y$, the sum over
    the first component reduces to its diagonal term, which gives
    $x\otimes\operatorname{op}(w)y$. The bound follows from
    $\operatorname{op}(w^W)(x\otimes v)-x\otimes v'=x\otimes(\operatorname{op}(w)v-v')$ and
    $\lVert x\otimes u\rVert=\lVert x\rVert\,\lVert u\rVert$. The operator of a relabelling is that
    of the relabelling of party circuits, a permutation matrix, and its action on vectors and on
    product vectors is a direct computation.
\end{proof}

\subsection{Ownership sheets}

Throughout, $\Lambda$ is a finite set of sites, $q\in\N$, $[q]=\{0,\dots,q-1\}$, and
$[q]^R$ is the set of configurations of a region $R\subseteq\Lambda$. The parties form a set
$\mathcal P$.

\begin{definition}[Ownership sheets and their reference vectors]
    \label{def:party_sheet}
    \lean{QuantumCircuit.Party.Sheet, QuantumCircuit.Party.Sheet.system,
        QuantumCircuit.Party.Sheet.refVec, QuantumCircuit.Party.Sheet.refVec_apply,
        QuantumCircuit.Party.Sheet.norm_refVec, QuantumCircuit.Party.Sheet.norm_refVec_le,
        QuantumCircuit.Party.Sheet.twoConfigEquiv,
        QuantumCircuit.Party.Sheet.twoConfigEquiv_apply,
        QuantumCircuit.Party.Sheet.twoConfigEquiv_fst_eq,
        QuantumCircuit.Party.Sheet.twoConfigEquiv_snd_eq,
        QuantumCircuit.Party.Sheet.prodVec_refVec}
    \leanok
    \uses{def:sites_block, def:party_original_frame_sequence, lem:kronecker_vectors}
    A sheet $S$ on $\Lambda$ is an owner map $o_S:\Lambda\to\mathcal P$; the register $\C^q$ of
    a site $x$ is held by $o_S(x)$ \cite[Section~6]{OpenAI2026PolynomialPEPS}. Its system $X^S$
    has at the party $p$ the basis $X^S_p=[q]^{o_S^{-1}(p)}$ of configurations of the sites $p$
    owns, the blocked system of $[q]^\Lambda$ along $o_S$. Its canonical coordinates are the
    bijection $\gamma_S$ from the configurations of $X^S$ to $[q]^\Lambda$ with
    $\gamma_S(z)(x)=z_{o_S(x)}(x)$. A vector $\Omega\in\C^{[q]^\Lambda}$ read on the sheet is
    $\Omega_S=\Omega\circ\gamma_S$, and $\lVert\Omega_S\rVert=\lVert\Omega\rVert$, so
    $\lVert\Omega_S\rVert\leq1$ when $\lVert\Omega\rVert\leq1$.

    Two sheets $S_1,S_2$ held at once form the product system $X^{S_1}\times X^{S_2}$. Its
    canonical coordinates are the bijection $\gamma_{S_1,S_2}$ onto $[q]^\Lambda\times[q]^\Lambda$,
    $\gamma_{S_1,S_2}(z)=(\gamma_{S_1}(z^1),\gamma_{S_2}(z^2))$ for the pair $(z^1,z^2)$ of
    configurations paired with $z$: the coordinate of a site $x$ on the first sheet is read at
    $o_{S_1}(x)$, on the second at $o_{S_2}(x)$. The product of the reference vectors is
    $\Omega_{S_1}\otimes\Omega_{S_2}=(\Omega\otimes\Omega)\circ\gamma_{S_1,S_2}$.
\end{definition}

\begin{definition}[Change of owner and the sets of a birth]
    \label{def:party_change_owner}
    \lean{QuantumCircuit.Party.Sheet.changeOwner, QuantumCircuit.Party.Sheet.birthEnv,
        QuantumCircuit.Party.Sheet.mem_birthEnv, QuantumCircuit.Party.Sheet.disjoint_birthEnv,
        QuantumCircuit.Party.Sheet.changeOwner_eq_self,
        QuantumCircuit.Party.Sheet.changeOwner_birth_owners,
        QuantumCircuit.Party.Sheet.owner_eq_iff_of_birth,
        QuantumCircuit.Party.Sheet.changeOwner_owner_eq_self_iff,
        QuantumCircuit.Party.Sheet.changeOwner_owner_eq_iff,
        QuantumCircuit.Party.Sheet.changeOwner_owner_eq_iff_of_ne,
        QuantumCircuit.Party.Sheet.changeOwner_self_owner_eq_iff,
        QuantumCircuit.Party.Sheet.mem_sdiff_birthEnv_iff}
    \leanok
    \uses{def:party_sheet}
    For $T\subseteq\Lambda$ and a party $Q$, the sheet $S^{T\to Q}$ hands the sites of $T$ to
    $Q$: its owner is $Q$ on $T$ and $o_S$ elsewhere. For a party $P_\circ$, the set
    $E_S(P_\circ)=\{x:o_S(x)\neq P_\circ\}$ is the set $E$ of
    \cite[Lemma~6.5]{OpenAI2026PolynomialPEPS}, whose hole part is empty on a sheet. Let
    $E=E_S(P_\circ)$ and $U=\Lambda\setminus(T\cup E)$.

    For every $T$ and $Q$, the sites owned by $Q$ in $S^{T\to Q}$ are those of $T$ and those
    owned by $Q$ in $S$; if $Q\neq P_\circ$ the sites owned by $P_\circ$ in $S^{T\to Q}$ are
    those of $U$; and the sites owned by $P_\circ$ in $S^{T\to P_\circ}$ are those of
    $T\cup U$. If moreover every site of $T$ is owned by $P_\circ$ in $S$, then $T$ and $E$ are
    disjoint, $S^{T\to P_\circ}=S$, the sites owned by $P_\circ$ in $S$ are those of $T\cup U$,
    a site lies in $(T\cup E)\setminus T$ exactly when it is not owned by $P_\circ$, every
    party $r\notin\{P_\circ,Q\}$ owns the same sites in $S$ and in $S^{T\to Q}$, the sites of
    $U$ are owned by $P_\circ$ in both sheets, the sites of $T$ by $Q$ in $S^{T\to Q}$, and the
    owners on $E$ are unchanged.
\end{definition}

\begin{definition}[Bases of a party read on regions]
    \label{def:party_sheet_region_coords}
    \lean{QuantumCircuit.Party.Sheet.regionEquiv,
        QuantumCircuit.Party.Sheet.regionEquiv_apply,
        QuantumCircuit.Party.Sheet.regionEquiv_symm_apply,
        QuantumCircuit.Party.Sheet.splitEquiv,
        QuantumCircuit.Party.Sheet.splitEquiv_apply_fst,
        QuantumCircuit.Party.Sheet.splitEquiv_apply_snd,
        QuantumCircuit.Party.Sheet.complSplitEquiv,
        QuantumCircuit.Party.Sheet.system_eq_of_owner_iff,
        QuantumCircuit.Party.Sheet.cast_system_eq_regionEquiv,
        QuantumCircuit.Party.Sheet.birthEquivP, QuantumCircuit.Party.Sheet.birthEquivP',
        QuantumCircuit.Party.Sheet.birthEquivQ, QuantumCircuit.Party.Sheet.birthEquivSelf,
        QuantumCircuit.Party.Sheet.cast_changeOwner_system_iff}
    \leanok
    \uses{def:party_sheet, def:party_change_owner}
    If the sites owned by $p$ in $S$ form the region $R$, restriction identifies $X^S_p$ with
    $[q]^R$; if they form the disjoint union $R_1\sqcup R_2$, it identifies $X^S_p$ with
    $[q]^{R_1}\times[q]^{R_2}$; if every site outside a region $B$ is owned by $P_\circ$, it
    identifies $X^S_{P_\circ}$ with $[q]^{B\cap o_S^{-1}(P_\circ)}\times[q]^{\Lambda\setminus B}$.
    These are relabellings of sites \cite[Section~6]{OpenAI2026PolynomialPEPS}. If $p$ owns the
    same sites in two sheets $S'$ and $S$, then $X^{S'}_p=X^S_p$, and the identification of these
    equal bases is the restriction to the sites of $p$ in $S'$ followed by the inverse of the
    restriction in $S$, site by site.

    In the setting of a birth, with every site of $T$ owned by $P_\circ$, $E=E_S(P_\circ)$ and
    $U=\Lambda\setminus(T\cup E)$, these give the bijections
    \begin{align}
        \beta_P     & : X^S_{P_\circ}\to[q]^T\times[q]^U,                    \notag \\
        \beta_\circ & : X^{S^{T\to P_\circ}}_{P_\circ}\to[q]^T\times[q]^U, \notag   \\
        \beta'_P    & : X^{S^{T\to Q}}_{P_\circ}\to[q]^U,                  \notag   \\
        \beta_Q     & : X^{S^{T\to Q}}_{Q}\to[q]^T\times X^S_Q,
        \label{eq:party_birth_bases}
    \end{align}
    the last two for $Q\neq P_\circ$, where $\beta_Q$ splits a configuration of the sites of $Q$
    after the birth into its part on $T$ and its part on the sites $Q$ owned before. For a party
    $r\notin\{P_\circ,Q\}$, a configuration $y$ of $X^{S^{T\to Q}}$ and a configuration $x$ of
    $X^S$ have the same component at $r$, the two equal bases being identified, exactly when
    $\gamma_{S^{T\to Q}}(y)$ and $\gamma_S(x)$ agree on the sites owned by $r$.
\end{definition}

\subsection{Splitting data and the canonical map of a birth}

\begin{definition}[Splitting data and the canonical map of a birth]
    \label{def:party_splitting_data}
    \lean{ManyBody.SplittingData, ManyBody.SplittingData.error, ManyBody.birthKernel,
        ManyBody.birthEffect, ManyBody.birthOp, ManyBody.configBirthOp}
    \leanok
    \uses{def:sites_config, thm:three_way_cut_coordinates,
        def:canonical_product_purification, def:construction_idle_gates}
    Let $X$ be a system on a finite set of sites $V$, let $A\subseteq B$ be regions and
    $U=V\setminus B$. Write $\mathcal B_A=\C^{X_A}$ and
    $\mathcal B_M=\C^{X_{B\setminus A}\sqcup X_U}$. Splitting data $\sigma$ for $(A,B)$
    \cite[Lemma~6.4]{OpenAI2026PolynomialPEPS} are an isometry
    $V_\sigma:\C^{X_U}\to\mathcal B_A\otimes\mathcal B_M$, a unit vector
    $s\in\C^{X_A}\otimes\mathcal B_A$ and a unit vector
    $s'\in\C^{X_{B\setminus A}}\otimes\mathcal B_M$. Their error on $\Omega\in\C^{X_V}$ is
    \begin{align}
        \varepsilon_\sigma(\Omega)
        =\lVert(1_{X_A\times X_{B\setminus A}}\otimes V_\sigma)\Omega-s\otimes s'\rVert,
        \label{eq:party_splitting_error}
    \end{align}
    with $\Omega$ read in the three-way cut coordinates $((x_A,x_{B\setminus A}),x_U)$ of
    Theorem~\ref{thm:three_way_cut_coordinates} and
    $(s\otimes s')((a,m),(b,b'))=s(a,b)\,s'(m,b')$.

    The kernel of $\sigma$ is the matrix $K$ from $X_A\times X_U$ to the basis of
    $\mathcal B_M$ with
    $K_{b,(a,u)}=\sum_{\beta\in X_A}\overline{s(a,\beta)}\,(V_\sigma)_{(\beta,b),u}$, that is
    $K=\langle s|(1_{X_A}\otimes V_\sigma)$. With $\mathcal V=1\otimes V_\sigma$ on
    $(X_A\times X_{B\setminus A})\times X_U$ and $S_s$ the preparation matrix of $s$ on
    $X_A\times\mathcal B_A$, the effect of a birth is $S_s^\dagger\mathcal V$ and the canonical
    map of the birth is
    \begin{align}
        \mathbf B_\sigma=(S_s^\dagger\mathcal V)^\dagger(S_s^\dagger\mathcal V)
        =\mathcal V^\dagger\bigl(|s\rangle\langle s|\otimes1\bigr)\mathcal V,
        \label{eq:party_birth_map}
    \end{align}
    read on the configurations of $X_V$ along the three-way cut coordinates
    \cite[Lemma~6.5]{OpenAI2026PolynomialPEPS}. In the coordinates $(x_{B\setminus A},(x_A,x_U))$
    it is $\mathbf B_\sigma=1_{X_{B\setminus A}}\otimes K^\dagger K$
    \cite[proof of Lemma~6.5]{OpenAI2026PolynomialPEPS}.
\end{definition}

For a sheet $S$ and a set $T$ of sites owned by $P_\circ$, splitting data for a birth are
splitting data $\sigma$ on $[q]^\Lambda$ for $(T,T\cup E)$ with $E=E_S(P_\circ)$; then
$\mathcal B_T=\C^{[q]^T}$, $\mathcal B_E=\C^{[q]^E\sqcup[q]^U}$ and $(T\cup E)\setminus T=E$.

\begin{theorem}[The canonical map of a birth in the coordinates of the parties]
    \label{thm:party_config_birth_op}
    \lean{QuantumCircuit.Party.Sheet.configBirthOp_apply_changeOwner,
        QuantumCircuit.Party.Sheet.configBirthOp_apply_changeOwner_self}
    \leanok
    \uses{def:party_splitting_data, def:party_sheet_region_coords, def:party_change_owner}
    Let every site of $T$ be owned by $P_\circ$ in the sheet $S$, let $\sigma$ be splitting data
    for the birth with kernel $K$, and use the bijections~(\ref{eq:party_birth_bases}). For
    $Q\neq P_\circ$, a configuration $y$ of $X^{S^{T\to Q}}$ and a configuration $x$ of $X^S$,
    \begin{align}
        \mathbf B_\sigma(\gamma_{S^{T\to Q}}(y),\gamma_S(x))
        =\delta\,(K^\dagger K)\bigl((\beta_Q(y_Q)_1,\beta'_P(y_{P_\circ})),
        \beta_P(x_{P_\circ})\bigr),
        \label{eq:party_birth_entry}
    \end{align}
    where $\delta=1$ if $\beta_Q(y_Q)_2=x_Q$ and $\gamma_{S^{T\to Q}}(y)$, $\gamma_S(x)$ agree on
    the sites of every party $r\notin\{P_\circ,Q\}$, and $\delta=0$ otherwise. For a configuration
    $y$ of $X^{S^{T\to P_\circ}}$,
    $\mathbf B_\sigma(\gamma_{S^{T\to P_\circ}}(y),\gamma_S(x))
        =\delta_\circ\,(K^\dagger K)(\beta_\circ(y_{P_\circ}),\beta_P(x_{P_\circ}))$, where
    $\delta_\circ=1$ if the two configurations agree on the sites of every party $r\neq P_\circ$
    and $\delta_\circ=0$ otherwise.
\end{theorem}
\begin{proof}\leanok
    \uses{def:party_splitting_data, def:party_sheet_region_coords, def:party_change_owner}
    By Definition~\ref{def:party_splitting_data},
    $\mathbf B_\sigma=1_{[q]^E}\otimes K^\dagger K$ in the coordinates $(x_E,(x_T,x_U))$. The
    factor $1_{[q]^E}$ compares the parts of the two configurations on $E$, the sites not owned
    by $P_\circ$: for $Q$ these are the sites it owned before, whose part in $y$ is
    $\beta_Q(y_Q)_2$, and for $r\notin\{P_\circ,Q\}$ the sites $r$ owns in both sheets. The parts
    on $T$ and $U$ are read at $Q$ and $P_\circ$ after the birth and at $P_\circ$ before it,
    through $\beta_Q$, $\beta'_P$ and $\beta_P$. For $Q=P_\circ$ the parts on $T$ and $U$ after
    the birth are both read at $P_\circ$ through $\beta_\circ$.
\end{proof}

\subsection{Births and deaths}

In the rest of this subsection $S$ is a sheet, every site of $T$ is owned by $P_\circ$,
$E=E_S(P_\circ)$, $U=\Lambda\setminus(T\cup E)$, and $\sigma$ are splitting data for the birth
with isometry $V_\sigma:\C^{[q]^U}\to\mathcal B_T\otimes\mathcal B_E$, unit vector
$s\in\C^{[q]^T}\otimes\mathcal B_T$ and kernel $K$.

\begin{definition}[Birth words and death chains]
    \label{def:party_birth_word}
    \lean{QuantumCircuit.Party.Sheet.birthPartyMatrix, QuantumCircuit.Party.Sheet.birthFamily,
        QuantumCircuit.Party.Sheet.birthWord, QuantumCircuit.Party.Sheet.deathFamily,
        QuantumCircuit.Party.Sheet.deathChain}
    \leanok
    \uses{def:party_splitting_data, def:party_sheet_region_coords, def:party_word,
        def:party_add_pair, def:party_pair_gates, def:party_circuit}
    Let $Q\neq P_\circ$. The contraction of the birth at $P_\circ$ is
    $M_\sigma=V_\sigma^\dagger(1_{\mathcal B_T}\otimes K)$, from
    $\mathcal B_T\otimes\C^{[q]^T}\otimes\C^{[q]^U}$ to $\C^{[q]^U}$.

    The birth word from $X^S$ to $X^{S^{T\to Q}}$ is the normalized pair source $s$ on a fresh
    copy of $[q]^T$ added at $Q$ and a fresh factor $\mathcal B_T$ added at $P_\circ$, followed
    by one local operation at every party: at $Q$ the matrix of the bijection $\beta_Q^{-1}$,
    so that the new copy of $T$ joins the sites of $Q$; at $P_\circ$ the matrix $M_\sigma$, read
    on the bases through $\beta_P$ and $\beta'_P$, so that the copy of $T$ held by $P_\circ$
    before the birth leaves $P_\circ$; and the identification of equal bases at every other
    party. In \cite[proof of Lemma~6.5]{OpenAI2026PolynomialPEPS} the contraction $K$ is applied
    before the pair source; it acts on factors other than those of the source, and here it is
    applied after the source, together with $V_\sigma^\dagger$, in one local operation at
    $P_\circ$.

    The death chain from $X^{S^{T\to Q}}$ to $X^S$ is the monomial made of one local operation
    at every party, the matrix of $\beta_Q$ at $Q$, the adjoint
    $M_\sigma^\dagger=(1_{\mathcal B_T}\otimes K^\dagger)V_\sigma$ at $P_\circ$ read through
    the same bijections, and the identification of equal bases elsewhere, followed by the
    normalized pair effect $\langle s|$ on the copy of $T$ at $Q$ and on $\mathcal B_T$ at
    $P_\circ$. Since $M_\sigma^\dagger$ acts on factors other than those of the pair effect, it is
    applied before the effect.
\end{definition}

\begin{theorem}[The contraction of a birth]
    \label{thm:party_birth_contraction}
    \lean{QuantumCircuit.Party.Sheet.norm_birthPartyMatrix_le_one,
        QuantumCircuit.Party.Sheet.sum_birthPartyMatrix_mul}
    \leanok
    \uses{def:party_birth_word}
    Let $T$, $U$, $B_T$ and $B_E$ be finite sets, $V$ a matrix from $\C^U$ to
    $\C^{B_T}\otimes\C^{B_E}$, $s$ a vector on $T\times B_T$, $K=\langle s|(1_T\otimes V)$ and
    $M=V^\dagger(1_{B_T}\otimes K)$. Then for all $t'\in T$, $u\in U$ and $(t,u')\in T\times U$,
    \begin{align}
        \sum_{b\in B_T}M_{u,(b,(t,u'))}\,s(t',b)=(K^\dagger K)_{(t',u),(t,u')},
        \label{eq:party_birth_contraction}
    \end{align}
    and if $V$ is an isometry and $s$ a unit vector, then $\lVert M\rVert\leq1$.
\end{theorem}
\begin{proof}\leanok
    \uses{def:party_splitting_data}
    Expanding the identity factor,
    $M_{u,(b,\tau)}=\sum_{b'}\overline{V_{(b,b'),u}}\,K_{b',\tau}$, and
    $\sum_b\overline{V_{(b,b'),u}}\,s(t',b)=\overline{K_{b',(t',u)}}$, which
    gives~(\ref{eq:party_birth_contraction}) after exchanging the two sums. For the norm,
    $\lVert V^\dagger\rVert\leq1$, and $\lVert1\otimes K\rVert\leq\lVert K\rVert\leq1$ because $K$
    is the effect $S_s^\dagger\mathcal V$ of Definition~\ref{def:party_splitting_data} with a
    one-point middle region, a product of two contractions.
\end{proof}

\begin{theorem}[Birth words and death chains implement the canonical map]
    \label{thm:party_birth_word}
    \lean{QuantumCircuit.Party.Sheet.op_birthWord,
        QuantumCircuit.Party.Sheet.isAllowed_birthWord,
        QuantumCircuit.Party.Sheet.usesOnly_birthWord,
        QuantumCircuit.Party.Sheet.sourceCount_birthWord,
        QuantumCircuit.Party.Sheet.op_deathChain,
        QuantumCircuit.Party.Sheet.isAllowed_deathChain,
        QuantumCircuit.Party.Sheet.usesOnly_deathChain,
        QuantumCircuit.Party.Sheet.forallGate_isRight_deathChain,
        QuantumCircuit.Party.Sheet.numGates_deathChain}
    \leanok
    \uses{def:party_birth_word, def:party_allowed, def:party_word, def:party_gate_counts}
    Let $Q\neq P_\circ$ and $S'=S^{T\to Q}$. The birth word is allowed, uses only $P_\circ$
    and $Q$, has exactly one pair source, and its operator is the canonical map in canonical
    coordinates: $\operatorname{op}(y,x)=\mathbf B_\sigma(\gamma_{S'}(y),\gamma_S(x))$ for
    configurations $y$ of $X^{S'}$ and $x$ of $X^S$. The death chain is allowed, uses only
    $P_\circ$ and $Q$, has exactly one gate, which is a pair effect, and its operator is
    $\operatorname{op}(x,y)=\mathbf B_\sigma(\gamma_S(x),\gamma_{S'}(y))$.
\end{theorem}
\begin{proof}\leanok
    \uses{thm:party_birth_contraction, thm:party_config_birth_op, def:party_birth_word,
        def:party_splitting_data}
    The operator of the birth word is the product over the parties of the local operations
    applied after the preparation of $s$. At the configurations $y$, $x$ and the fresh pair
    $(a,b)$, the factor at $Q$ is $1$ exactly when $\beta_Q(y_Q)=(a,x_Q)$, the factor at $P_\circ$
    is $M_\sigma(\beta'_P(y_{P_\circ}),(b,\beta_P(x_{P_\circ})))$, and the factor at every other
    party is $1$ exactly when the configurations agree on its sites. The sum over $a$ leaves
    $a=\beta_Q(y_Q)_1$, and the sum over $b$ against $s(a,b)$ gives $K^\dagger K$ by
    Theorem~\ref{thm:party_birth_contraction}; the result is the right side
    of~(\ref{eq:party_birth_entry}) in Theorem~\ref{thm:party_config_birth_op}. The death chain
    gives the same computation with the adjoint factors and the effect $\langle s|$, hence the
    complex conjugate of $\mathbf B_\sigma(\gamma_{S'}(y),\gamma_S(x))$, which is
    $\mathbf B_\sigma(\gamma_S(x),\gamma_{S'}(y))$ because $\mathbf B_\sigma$ is Hermitian
    by~(\ref{eq:party_birth_map}). Both are allowed: $s$ is a unit vector, the matrices of
    bijections are contractions and $\lVert M_\sigma\rVert\leq1$ by
    Theorem~\ref{thm:party_birth_contraction}. Outside $\{P_\circ,Q\}$ the local operations are
    identifications of equal bases, and the counts hold by definition.
\end{proof}

\begin{definition}[Births and deaths at one party]
    \label{def:party_private_birth}
    \lean{QuantumCircuit.Party.Sheet.privateBirthFamily,
        QuantumCircuit.Party.Sheet.privateBirthWord,
        QuantumCircuit.Party.Sheet.privateDeathChain}
    \leanok
    \uses{def:party_splitting_data, def:party_sheet_region_coords, def:party_word,
        def:party_pair_gates, def:party_circuit}
    When the party receiving $T$ is $P_\circ$ itself, no pair source is needed
    \cite[proof of Lemma~6.5]{OpenAI2026PolynomialPEPS}. The private birth word from $X^S$ to
    $X^{S^{T\to P_\circ}}$ is the word without gates made of one local operation at every party:
    the matrix $K^\dagger K$ at $P_\circ$, read on the bases through $\beta_P$ and $\beta_\circ$,
    and the identification of equal bases at every other party. The private death chain from
    $X^{S^{T\to P_\circ}}$ to $X^S$ is the monomial without gates made of the adjoints of these
    local operations.
\end{definition}

\begin{theorem}[Births and deaths at one party implement the canonical map]
    \label{thm:party_private_birth}
    \lean{QuantumCircuit.Party.Sheet.op_privateBirthWord,
        QuantumCircuit.Party.Sheet.isAllowed_privateBirthWord,
        QuantumCircuit.Party.Sheet.usesOnly_privateBirthWord,
        QuantumCircuit.Party.Sheet.sourceCount_privateBirthWord,
        QuantumCircuit.Party.Sheet.op_privateDeathChain,
        QuantumCircuit.Party.Sheet.isAllowed_privateDeathChain,
        QuantumCircuit.Party.Sheet.usesOnly_privateDeathChain,
        QuantumCircuit.Party.Sheet.numGates_privateDeathChain}
    \leanok
    \uses{def:party_private_birth, def:party_allowed, def:party_word, def:party_gate_counts}
    Let $S'=S^{T\to P_\circ}$. The private birth word is allowed, uses only $P_\circ$, has no
    pair source, and its operator is
    $\operatorname{op}(y,x)=\mathbf B_\sigma(\gamma_{S'}(y),\gamma_S(x))$. The private death chain
    is allowed, uses only $P_\circ$, has no gate, and its operator is
    $\operatorname{op}(x,y)=\mathbf B_\sigma(\gamma_S(x),\gamma_{S'}(y))$.
\end{theorem}
\begin{proof}\leanok
    \uses{thm:party_config_birth_op, def:party_private_birth, def:party_splitting_data}
    The operator of the private birth word is the product over the parties of its local
    operations: $(K^\dagger K)(\beta_\circ(y_{P_\circ}),\beta_P(x_{P_\circ}))$ at $P_\circ$ times
    the agreement of the configurations on the sites of every other party, which is the second
    formula of Theorem~\ref{thm:party_config_birth_op}. The operator of the death chain is the
    adjoint of this product, and $\mathbf B_\sigma$ is Hermitian by~(\ref{eq:party_birth_map}).
    Both are allowed because $\lVert K\rVert\leq1$ (see the proof of
    Theorem~\ref{thm:party_birth_contraction}), so $\lVert K^\dagger K\rVert\leq1$, and the
    identifications of equal bases are contractions; the remaining statements hold by
    definition.
\end{proof}

\begin{theorem}[Homogeneous birth and death for given splitting data]
    \label{thm:party_birth}
    \lean{QuantumCircuit.Party.Sheet.exists_birthWord_of_environment,
        QuantumCircuit.Party.Sheet.exists_deathChain_of_environment,
        QuantumCircuit.Party.Sheet.exists_birthWord_of_splittingData,
        QuantumCircuit.Party.Sheet.exists_deathChain_of_splittingData}
    \leanok
    \uses{def:party_change_owner, def:party_splitting_data, def:party_word,
        def:party_pair_gates, def:party_allowed, def:party_gate_counts}
    Let $S$ be a sheet, $P_\circ$ and $Q$ parties, not necessarily distinct, $T\subseteq\Lambda$
    a set of sites owned by $P_\circ$, and $E\subseteq\Lambda$ a set disjoint from $T$ such that
    every site outside $T\cup E$ is owned by $P_\circ$; sites of $E$ may be owned by $P_\circ$.
    Let $\sigma$ be splitting data on $[q]^\Lambda$ for $(T,T\cup E)$, let
    $\Omega\in\C^{[q]^\Lambda}$ and $S'=S^{T\to Q}$.
    \begin{enumerate}
        \item There is an allowed word $w$ from $X^S$ to $X^{S'}$ that uses only $P_\circ$ and $Q$,
              has at most one pair source, whose operator is
              $\operatorname{op}(w)(y,x)=\mathbf B_\sigma(\gamma_{S'}(y),\gamma_S(x))$, and with
              $\lVert\operatorname{op}(w)\Omega_S-\Omega_{S'}\rVert\leq\varepsilon_\sigma(\Omega)$.
        \item There is an allowed monomial $M$ from $X^{S'}$ to $X^S$ that uses only $P_\circ$
              and $Q$, whose gates are pair effects and at most one, whose operator is
              $\operatorname{op}(M)(x,y)=\mathbf B_\sigma(\gamma_S(x),\gamma_{S'}(y))$, and with
              $\lVert\operatorname{op}(M)\Omega_{S'}-\Omega_S\rVert\leq\varepsilon_\sigma(\Omega)$.
    \end{enumerate}
    In particular both hold for $E=E_S(P_\circ)$. This is the part of
    \cite[Lemma~6.5]{OpenAI2026PolynomialPEPS} that holds for any splitting data: the birth and
    the death involve only $P_\circ$ and $Q$, use at most one normalized pair source or pair
    effect, and move the reference vector by at most the error of the splitting data. The
    hypothesis of the lemma on the mutual information of $T$ and $E$ enters only through the
    existence of splitting data of small error \cite[Lemma~6.4]{OpenAI2026PolynomialPEPS},
    which is not part of this statement.
\end{theorem}
\begin{proof}\leanok
    \uses{def:party_change_owner, def:party_splitting_data, def:party_sheet}
    The canonical coordinates are bijections, so a word or monomial whose operator is
    $\mathbf B_\sigma$ in canonical coordinates moves $\Omega_S$ to within
    $\lVert\mathbf B_\sigma\Omega-\Omega\rVert$ of $\Omega_{S'}$. With the projection
    $\Pi=S_sS_s^\dagger$, which fixes $s\otimes s'$,
    $(\mathbf B_\sigma-1)\Omega=\mathcal V^\dagger(\Pi-1)(\mathcal V\Omega-s\otimes s')$ because
    $\mathcal V$ is an isometry, and $\lVert\mathcal V^\dagger\rVert\leq1$,
    $\lVert\Pi-1\rVert\leq1$ give
    $\lVert\mathbf B_\sigma\Omega-\Omega\rVert\leq\varepsilon_\sigma(\Omega)$
    by~(\ref{eq:party_splitting_error}). For $Q=P_\circ$ the sheet is unchanged,
    $S^{T\to P_\circ}=S$, and the private birth word and private death chain of
    Definition~\ref{def:party_private_birth} serve, modified so that the local operation at
    $P_\circ$ acts as the identity on the sites of $E$ owned by $P_\circ$. For $Q\neq P_\circ$ the
    birth word and death chain of Definition~\ref{def:party_birth_word} serve, with the same
    modification of the local operation at $P_\circ$. The computation of
    Theorems~\ref{thm:party_birth_word} and~\ref{thm:party_private_birth}, with the sites of
    $E$ owned by $P_\circ$ compared by the identity at $P_\circ$, shows that their operators are
    $\mathbf B_\sigma$ in canonical coordinates, and the counts and the norms of the local
    operations are as there. For $E=E_S(P_\circ)$, the set $T$ is disjoint from $E$ and the sites outside
    $T\cup E$ are owned by $P_\circ$ by Definition~\ref{def:party_change_owner}.
\end{proof}

\subsection{Two-sheet exchange}

\begin{definition}[Exchange of two sheets inside a region]
    \label{def:party_exchange_sets}
    \lean{QuantumCircuit.Party.Sheet.swapIn, QuantumCircuit.Party.Sheet.exchangeEnv,
        QuantumCircuit.Party.Sheet.mem_exchangeEnv,
        QuantumCircuit.Party.Sheet.owner_eq_of_notMem_exchangeEnv,
        QuantumCircuit.Party.Sheet.sdiff_subset_compl_inter_union_sdiff,
        QuantumCircuit.Party.Sheet.swapIn_owner_eq_of_notMem,
        QuantumCircuit.Party.Sheet.swapIn_owner_eq_of_notMem_union,
        QuantumCircuit.Party.Sheet.swapIn_owner_eq_of_notMem_union'}
    \leanok
    \uses{def:party_sheet}
    For sheets $S_1,S_2$ and a region $Y$, the sheet $S_1\ast_YS_2$ has the owner $o_{S_2}$ on
    $Y$ and $o_{S_1}$ outside $Y$. The exchange inside $Y$ replaces $(S_1,S_2)$ by
    $(S_1',S_2')=(S_1\ast_YS_2,S_2\ast_YS_1)$ \cite[Lemma~6.6]{OpenAI2026PolynomialPEPS}. For a
    party $P_\circ$, the set $Z_\circ$ of sites whose two owners are not both $P_\circ$ is the set
    $Z$ of the lemma without its hole parts, which are empty on sheets; so let $Z$ be any set of
    sites containing $Z_\circ$, that is, outside of which both owners are $P_\circ$. Put
    $T=Z\cap Y$, $E=Z\setminus Y$, $A=Y\setminus Z$ and $U=\Lambda\setminus Z$. Then
    $T\cup E=Z$, $A\subseteq\Lambda\setminus(T\cup E)$, and every site outside $T\cup E$ is owned
    by $P_\circ$ on both sheets after the exchange.
\end{definition}

\begin{definition}[The map of an exchange]
    \label{def:party_exchange_map}
    \lean{QuantumCircuit.Party.Sheet.exchangeOp}
    \leanok
    \uses{def:party_exchange_sets, def:party_splitting_data, def:sites_pair_swap}
    Let $\sigma$ be splitting data on $[q]^\Lambda$ for $(T,T\cup E)$, with isometry
    $V_\sigma:\C^{[q]^U}\to\mathcal B_T\otimes\mathcal B_E$. On two copies
    $\C^{[q]^\Lambda\times[q]^\Lambda}$, let $F_R$ be the exchange of the two copies on a region
    $R$, and let the buffer correction $D_U$ act, in the three-way cut coordinates of both
    copies, as the identity on the two copies of $[q]^T\times[q]^E$ and as
    $(V_\sigma\otimes V_\sigma)^\dagger F_{\mathcal B_T}(V_\sigma\otimes V_\sigma)$ on the two
    copies of $[q]^U$, where $F_{\mathcal B_T}$ exchanges the $\mathcal B_T$ factors of
    $(\mathcal B_T\otimes\mathcal B_E)^{\otimes2}$. The map of the exchange is
    $C_\sigma=D_UF_AF_Y$ \cite[proof of Lemma~6.6]{OpenAI2026PolynomialPEPS}.
\end{definition}

\begin{definition}[Renaming and corrections of an exchange]
    \label{def:party_exchange_words}
    \lean{QuantumCircuit.Party.Sheet.swapInEquiv, QuantumCircuit.Party.Sheet.renameWord,
        QuantumCircuit.Party.Sheet.exchangeEquivP, QuantumCircuit.Party.Sheet.correctionAt,
        QuantumCircuit.Party.Sheet.correctionWord}
    \leanok
    \uses{def:party_exchange_map, def:party_sheet_region_coords, def:party_relabel,
        def:party_word, def:party_circuit}
    At every party $p$, the bijection $\rho_p$ from $X^{S_1'}_p\times X^{S_2'}_p$ to
    $X^{S_1}_p\times X^{S_2}_p$ moves the coordinate of a site of $Y$ to the other sheet and keeps
    the coordinates of the other sites. The renaming $\mathcal R$ is the word relabelling the
    bases along $\rho$; every register keeps its party. Restriction identifies
    $X^{S_1'}_{P_\circ}\times X^{S_2'}_{P_\circ}$ with the product of the configurations of the
    sites of $Z$ owned by $P_\circ$ on each sheet and of two copies of $[q]^U$. The correction at
    $P_\circ$ is the identity on the first factor tensored with
    \begin{align}
        W_\sigma=(V_\sigma\otimes V_\sigma)^\dagger F_{\mathcal B_T}(V_\sigma\otimes V_\sigma)F_A
        \label{eq:party_correction}
    \end{align}
    on the two copies of $[q]^U$, with $F_A$ exchanging the parts on $A\subseteq U$ of the two
    copies. The correction word is the word without gates whose local operation is the
    correction at $P_\circ$ and the identity at every other party.
\end{definition}

\begin{theorem}[Renaming and corrections in canonical coordinates]
    \label{thm:party_exchange_words}
    \lean{QuantumCircuit.Party.Sheet.twoConfigEquiv_symm_permCongr_pairSwap,
        QuantumCircuit.Party.Sheet.op_renameWord,
        QuantumCircuit.Party.Sheet.isAllowed_renameWord,
        QuantumCircuit.Party.Sheet.sourceCount_renameWord,
        QuantumCircuit.Party.Sheet.norm_correctionAt_le_one,
        QuantumCircuit.Party.Sheet.op_correctionWord,
        QuantumCircuit.Party.Sheet.isAllowed_correctionWord,
        QuantumCircuit.Party.Sheet.usesOnly_correctionWord,
        QuantumCircuit.Party.Sheet.sourceCount_correctionWord,
        QuantumCircuit.Party.Sheet.op_comp_renameWord_correctionWord}
    \leanok
    \uses{def:party_exchange_words, def:party_allowed}
    In the setting of Definition~\ref{def:party_exchange_words}:
    \begin{enumerate}
        \item The exchange of the two copies on $Y$, read from the canonical coordinates of
              $(S_1',S_2')$ to those of $(S_1,S_2)$, is the bijection $(\rho_p)_p$, that is
              $\gamma_{S_1,S_2}^{-1}\circ F_Y\circ\gamma_{S_1',S_2'}=(\rho_p)_p$. Hence the operator
              of $\mathcal R$ is $F_Y$ in canonical coordinates, and $\mathcal R$ is allowed and
              has no pair source.
        \item The correction at $P_\circ$ is a contraction. The correction word is allowed, uses
              only $P_\circ$, has no pair source, and its operator in the canonical coordinates of
              $(S_1',S_2')$ is $D_UF_A$.
        \item The renaming followed by the correction word has the operator $C_\sigma$, read from
              the canonical coordinates of $(S_1,S_2)$ to those of $(S_1',S_2')$.
    \end{enumerate}
\end{theorem}
\begin{proof}\leanok
    \uses{def:party_exchange_words, def:party_exchange_sets, def:party_relabel,
        def:sites_pair_swap}
    The first statement is checked site by site: a site of $Y$ in the first sheet after the
    exchange is read at its owner in the second sheet before it, and conversely, while the other
    sites stay on their sheets. The operator of a relabelling is the matrix of the bijection of
    configurations, which gives $F_Y$. The norm of the correction is at most
    $\lVert W_\sigma\rVert\leq1$, since $V_\sigma\otimes V_\sigma$ is an isometry and the
    exchanges are permutation matrices. Since $A\subseteq U$, the map $D_UF_A$ acts as the
    identity on the two copies of the configurations of $Z$ and as $W_\sigma$ on the two copies
    of $[q]^U$. Every site outside $Z$ is owned by $P_\circ$ on both sheets after the exchange,
    so the parts of two configurations on the sites of the other parties are compared by the
    identity local operations, and the parts on the sites of $Z$ owned by $P_\circ$ by the
    identity factor of the correction; this gives the operator of the correction word. The last
    statement is the product of the first two.
\end{proof}

\begin{theorem}[Two-sheet exchange for given splitting data]
    \label{thm:party_exchange}
    \lean{QuantumCircuit.Party.Sheet.exists_exchangeWord_of_splittingData,
        QuantumCircuit.Party.Sheet.exchangeOp_eq}
    \leanok
    \uses{def:party_exchange_map, def:party_exchange_words, def:party_allowed, def:party_word}
    Let $S_1,S_2$ be sheets, $Y$ a region, $P_\circ$ a party, $Z$ a set of sites outside of
    which both owners are $P_\circ$, with $T$, $E$, $A$, $U$ as in
    Definition~\ref{def:party_exchange_sets}, and let $\sigma$ be splitting data on
    $[q]^\Lambda$ for $(T,T\cup E)$. Then
    \begin{align}
        C_\sigma=D_UF_AF_Y=D_UF_T,
        \label{eq:party_exchange_net}
    \end{align}
    and there is an allowed word $c$ from $X^{S_1'}\times X^{S_2'}$ to itself that uses only
    $P_\circ$ and has no pair source, such that the renaming $\mathcal R$ followed by $c$ has the
    operator $C_\sigma$, read from the canonical coordinates of $(S_1,S_2)$ to those of
    $(S_1',S_2')$, and, for every unit vector $\Omega\in\C^{[q]^\Lambda}$,
    \begin{align}
        \bigl\lVert\operatorname{op}(\mathcal R;c)(\Omega_{S_1}\otimes\Omega_{S_2})
        -\Omega_{S_1'}\otimes\Omega_{S_2'}\bigr\rVert\leq4\,\varepsilon_\sigma(\Omega).
        \label{eq:party_exchange_error}
    \end{align}
    This is \cite[Lemma~6.6]{OpenAI2026PolynomialPEPS} for given splitting data: the exchange is
    a relabelling of the bases of the parties whose sites change sheet, no register changing
    party, followed by corrections private to $P_\circ$. It has no gate, so it is a bounded
    change for every set of participants (Theorem~\ref{thm:party_resources_bounded}).
\end{theorem}
\begin{proof}\leanok
    \uses{thm:party_exchange_words, def:party_sheet, def:sites_pair_swap}
    Since $A\subseteq Y$, the exchanges compose to $F_AF_Y=F_{Y\setminus A}$, and
    $Y\setminus(Y\setminus Z)=Z\cap Y=T$, which is~(\ref{eq:party_exchange_net}). The word $c$ is
    the correction word, and Theorem~\ref{thm:party_exchange_words} gives its properties and the
    operator of $\mathcal R$ followed by $c$. By Definition~\ref{def:party_sheet},
    $\Omega_{S_1}\otimes\Omega_{S_2}$ and $\Omega_{S_1'}\otimes\Omega_{S_2'}$ are $\Omega\otimes\Omega$
    in the canonical coordinates, so the left side of~(\ref{eq:party_exchange_error}) is
    $\lVert D_UF_T(\Omega\otimes\Omega)-\Omega\otimes\Omega\rVert$. With
    $\mathcal V=1\otimes V_\sigma$ and $\zeta=s\otimes s'$,
    $D_UF_T=(\mathcal V^{\otimes2})^\dagger F_{T\mathcal B_T}\mathcal V^{\otimes2}$, where
    $F_{T\mathcal B_T}$ exchanges the parts on $[q]^T\times\mathcal B_T$ of the two copies, and
    $F_{T\mathcal B_T}$ fixes $\zeta\otimes\zeta$, so
    $D_UF_T(\Omega\otimes\Omega)-\Omega\otimes\Omega=(\mathcal V^{\otimes2})^\dagger
        (F_{T\mathcal B_T}-1)\bigl((\mathcal V\Omega)^{\otimes2}-\zeta\otimes\zeta\bigr)$. Since
    $\lVert F_{T\mathcal B_T}-1\rVert\leq2$ and
    $\lVert(\mathcal V\Omega)^{\otimes2}-\zeta\otimes\zeta\rVert\leq
        2\lVert\mathcal V\Omega-\zeta\rVert$ for unit vectors, the bound
    is~$4\,\varepsilon_\sigma(\Omega)$ \cite[proof of Lemma~6.6]{OpenAI2026PolynomialPEPS}.
\end{proof}

\subsection{Stacks of sheets}

In this subsection the sheets form an arbitrary set $\kappa$; every $F\in\kappa$ carries a system
$X^F$ on the parties and, where needed, a reference vector $r_F\in\mathcal H_{X^F}$. Ownership
sheets with $X^S$ and $r_S=\Omega_S$ are one instance.

\begin{definition}[Stacks of sheets and lifted changes]
    \label{def:party_sheet_stack}
    \lean{QuantumCircuit.Party.stackSystem, QuantumCircuit.Party.stackVec,
        QuantumCircuit.Party.norm_stackVec, QuantumCircuit.Party.norm_stackVec_le,
        QuantumCircuit.Party.framesVec, QuantumCircuit.Party.norm_framesVec,
        QuantumCircuit.Party.liftOne, QuantumCircuit.Party.liftTwo,
        QuantumCircuit.Party.nonprivateCount_liftOne,
        QuantumCircuit.Party.forall_participants_liftOne,
        QuantumCircuit.Party.isExpansionBounded_liftOne,
        QuantumCircuit.Party.isMonomialBounded_liftOne,
        QuantumCircuit.Party.isSourceBounded_liftOne,
        QuantumCircuit.Party.nonprivateCount_liftTwo,
        QuantumCircuit.Party.forall_participants_liftTwo,
        QuantumCircuit.Party.isExpansionBounded_liftTwo,
        QuantumCircuit.Party.isMonomialBounded_liftTwo,
        QuantumCircuit.Party.isSourceBounded_liftTwo}
    \leanok
    \uses{def:party_original_frame_sequence}
    For a finite list $\mathbf F=(F_1,\dots,F_n)$ of sheets and an idle system $W$, the stack
    system $\mathrm{St}(\mathbf F;W)$ has at the party $p$ the basis
    $X^{F_1}_p\times(X^{F_2}_p\times\cdots\times(X^{F_n}_p\times W_p))$, and $W_p$ for the empty
    list. For $\xi\in\mathcal H_W$ the reference vector of the stack is
    $\omega_{\mathbf F}(\xi)=r_{F_1}\otimes(r_{F_2}\otimes\cdots\otimes(r_{F_n}\otimes\xi))$
    \cite[Proposition~7.1]{OpenAI2026PolynomialPEPS}. Its norm is
    $\lVert\omega_{\mathbf F}(\xi)\rVert=\prod_i\lVert r_{F_i}\rVert\,\lVert\xi\rVert$, at most
    $\lVert\xi\rVert$ when every $\lVert r_F\rVert\leq1$. For $W$ the system with one basis vector
    at every party and $\xi$ that vector, $\omega_{\mathbf F}=\omega_{\mathbf F}(\xi)$ has norm
    $\prod_i\lVert r_{F_i}\rVert$.

    Write $(\mathbf G,F,\mathbf H)$ for the list $\mathbf G$ followed by $F$ and then $\mathbf H$.
    For an original circuit $c$ from $X^F$ to $X^{F'}$, the lift $L_1(c)$ from
    $\mathrm{St}((\mathbf G,F,\mathbf H);W)$ to $\mathrm{St}((\mathbf G,F',\mathbf H);W)$ is, for
    $\mathbf G$ empty, the relabelling exchanging the two factors of
    $X^F\times\mathrm{St}(\mathbf H;W)$, then $c$ framed by $\mathrm{St}(\mathbf H;W)$, then the
    relabelling exchanging the two factors back; for $\mathbf G=(G,\mathbf G')$ it is the lift
    for $\mathbf G'$ framed by $X^G$. The lift $L_2(c)$ of an original circuit $c$ from
    $X^{F_1}\times X^{F_2}$ to $X^{F_1'}\times X^{F_2'}$, from
    $\mathrm{St}((\mathbf G,F_1,F_2,\mathbf H);W)$ to $\mathrm{St}((\mathbf G,F_1',F_2',\mathbf H);W)$,
    is defined in the same way, with relabellings moving $\mathrm{St}(\mathbf H;W)$ in front of
    the product of the two sheets and back. Lifting keeps the number of nonprivate gates, every
    condition on the participants of the gates, and the expansion, monomial and pair-source
    bounds.
\end{definition}

\begin{theorem}[The reference error of a lifted change]
    \label{thm:party_stack_lift_error}
    \lean{QuantumCircuit.Party.norm_op_liftOne_stackVec_sub_le,
        QuantumCircuit.Party.norm_op_liftTwo_stackVec_sub_le}
    \leanok
    \uses{def:party_sheet_stack}
    Let $\lVert r_G\rVert\leq1$ for every sheet $G$ and $\lVert\xi\rVert\leq1$. For an original
    circuit $c$ from $X^F$ to $X^{F'}$,
    \begin{align}
        \bigl\lVert\operatorname{op}(L_1(c))\,\omega_{(\mathbf G,F,\mathbf H)}(\xi)
        -\omega_{(\mathbf G,F',\mathbf H)}(\xi)\bigr\rVert
        \leq\lVert\operatorname{op}(c)r_F-r_{F'}\rVert,
    \end{align}
    and for an original circuit $c$ from $X^{F_1}\times X^{F_2}$ to $X^{F_1'}\times X^{F_2'}$ the
    lift $L_2(c)$ moves $\omega_{(\mathbf G,F_1,F_2,\mathbf H)}(\xi)$ to within
    $\lVert\operatorname{op}(c)(r_{F_1}\otimes r_{F_2})-r_{F_1'}\otimes r_{F_2'}\rVert$ of
    $\omega_{(\mathbf G,F_1',F_2',\mathbf H)}(\xi)$ \cite[Proposition~7.1]{OpenAI2026PolynomialPEPS}.
\end{theorem}
\begin{proof}\leanok
    \uses{thm:party_frame_prod_vec, def:party_sheet_stack}
    By induction on $\mathbf G$. For $\mathbf G$ empty, write $\zeta=\omega_{\mathbf H}(\xi)$, so
    $\lVert\zeta\rVert\leq\lVert\xi\rVert\leq1$. By Theorem~\ref{thm:party_frame_prod_vec} the first
    relabelling maps $r_F\otimes\zeta$ to $\zeta\otimes r_F$, the framed circuit maps it to
    $\zeta\otimes\operatorname{op}(c)r_F$, and the last relabelling to
    $\operatorname{op}(c)r_F\otimes\zeta$; the difference with $r_{F'}\otimes\zeta$ is
    $(\operatorname{op}(c)r_F-r_{F'})\otimes\zeta$, of norm at most
    $\lVert\operatorname{op}(c)r_F-r_{F'}\rVert$. For two sheets the relabellings map
    $r_{F_1}\otimes(r_{F_2}\otimes\zeta)$ to $\zeta\otimes(r_{F_1}\otimes r_{F_2})$ and back, up to
    the reassociation of the factors, which preserves norms. For $\mathbf G=(G,\mathbf G')$, the
    bound of Theorem~\ref{thm:party_frame_prod_vec} for the circuit framed by $X^G$, with
    $\lVert r_G\rVert\leq1$, reduces the statement to that for $\mathbf G'$.
\end{proof}

\begin{theorem}[Reordering the sheets is a relabelling]
    \label{thm:party_stack_perm}
    \lean{QuantumCircuit.Party.exists_relabel_stackVec_of_perm,
        QuantumCircuit.Party.perm_map_update_cons, QuantumCircuit.Party.perm_map_cons}
    \leanok
    \uses{def:party_sheet_stack}
    If the list $\mathbf F'$ is a permutation of the list $\mathbf F$, there are bijections
    $e_p:\mathrm{St}(\mathbf F';W)_p\to\mathrm{St}(\mathbf F;W)_p$ such that the relabelling
    along $e$ maps $\omega_{\mathbf F}(\xi)$ to $\omega_{\mathbf F'}(\xi)$ for every
    $\xi\in\mathcal H_W$. Moreover, for a map $a$ on a set of indices, an element $x$ of a list
    $l$ of indices, and values $m$ and $f$, the list $(m,a(l))$ is a permutation of
    $(a(x),m,a(l\setminus x))$, where $l\setminus x$ is $l$ with one occurrence of $x$ removed;
    if $l$ has no repetition and $a'$ is $a$ with the value at $x$ replaced by $f$, then
    $(m,a'(l))$ is a permutation of $(f,m,a(l\setminus x))$.
\end{theorem}
\begin{proof}\leanok
    \uses{thm:party_frame_prod_vec}
    By induction on the permutation. The empty list needs the identity. Adding the same sheet
    $F$ in front of both lists keeps the identity on $X^F$ beside the given bijections. Exchanging
    the first two sheets is the reassociation and commutation of the factors
    $X^{F}_p\times(X^{G}_p\times R)\to X^{G}_p\times(X^{F}_p\times R)$, which exchanges $r_F$ and
    $r_G$ in the product vector. A composite of permutations is handled by composing the
    bijections. The list statements follow from the permutation of $l$ to $(x,l\setminus x)$,
    mapped by $a$ or $a'$, followed by the exchange of the first two entries; when $l$ has no
    repetition, $a'$ and $a$ agree on $l\setminus x$.
\end{proof}

\subsection{Bounded changes and the initialization}

\begin{definition}[Bounded changes]
    \label{def:party_bounded_change}
    \lean{QuantumCircuit.Party.OriginalCircuit.IsBoundedChange,
        QuantumCircuit.Party.OriginalCircuit.IsBoundedChange.mono,
        QuantumCircuit.Party.HasBoundedChange, QuantumCircuit.Party.HasSourceBoundedChange,
        QuantumCircuit.Party.HasSourceBoundedChange.hasBoundedChange,
        QuantumCircuit.Party.HasBoundedChange.mono,
        QuantumCircuit.Party.HasSourceBoundedChange.mono}
    \leanok
    \uses{def:party_original_frame_sequence, def:party_sheet_stack}
    Let $c$ be an original circuit from $X$ to $Y$, $S$ a set of parties, $r,K\in\N$,
    $\varepsilon\in\R$, $x\in\mathcal H_X$ and $y\in\mathcal H_Y$. The circuit $c$ is a bounded
    change from $x$ to $y$ with participants in $S$, bounds $r$, $K$ and error $\varepsilon$ if
    \begin{enumerate}
        \item $c$ has at most one nonprivate gate;
        \item the cell of every gate of $c$ lies in $S$;
        \item every monomial of every gate expansion has at most $r$ pair effects and the
              coefficient sum of every gate expansion is at most $K$;
        \item every gate expansion has at most $K$ monomials;
        \item $\lVert\operatorname{op}(c)x-y\rVert\leq\varepsilon$.
    \end{enumerate}
    Two lists $\mathbf F$, $\mathbf F'$ of sheets with an idle system $W$ and
    $\xi\in\mathcal H_W$ have a bounded change if some original circuit from
    $\mathrm{St}(\mathbf F;W)$ to $\mathrm{St}(\mathbf F';W)$ is a bounded change from
    $\omega_{\mathbf F}(\xi)$ to $\omega_{\mathbf F'}(\xi)$, and a source-bounded change with
    constant $r_s$ if moreover that circuit satisfies the pair-source bound $r_s$
    \cite[Proposition~7.1]{OpenAI2026PolynomialPEPS}, \cite[Theorem~5.2]{OpenAI2026PolynomialPEPS}.
    A source-bounded change is a bounded change, and all three notions persist when $S$ is
    enlarged and $r$, $K$, $\varepsilon$ and $r_s$ are increased.
\end{definition}

\begin{theorem}[Bounded changes and circuits without nonprivate gates]
    \label{thm:party_bounded_change_basic}
    \lean{QuantumCircuit.Party.OriginalCircuit.norm_toEuclideanLin_op_le,
        QuantumCircuit.Party.OriginalCircuit.IsBoundedChange.comp_right,
        QuantumCircuit.Party.OriginalCircuit.IsBoundedChange.comp_left,
        QuantumCircuit.Party.OriginalCircuit.isBoundedChange_of_nonprivateCount_eq_zero,
        QuantumCircuit.Party.exists_isBoundedChange_of_eq,
        QuantumCircuit.Party.OriginalCircuit.isEmpty_nonprivateLocations_of_nonprivateCount_eq_zero,
        QuantumCircuit.Party.OriginalCircuit.forall_participants_of_nonprivateCount_eq_zero}
    \leanok
    \uses{def:party_bounded_change}
    An original circuit does not increase the norm of a vector. An original circuit without
    nonprivate gates has no gate occurrence, so every condition on the participants of its gates
    holds, and it is a bounded change from $x$ to $\operatorname{op}(c)x$ for every set $S$, all
    $r,K\in\N$ and every $\varepsilon\geq0$; in particular, for equal sheets $F=F'$ the empty
    circuit is a bounded change from $r_F$ to $r_{F'}$. If $c$ is a bounded change from $x$ to $y$
    and $e$ has no nonprivate gate, then $c$ followed by $e$ is a bounded change from $x$ to
    $\operatorname{op}(e)y$, and $e$ followed by $c$ is a bounded change from every $x'$ with
    $\operatorname{op}(e)x'=x$ to $y$, with the same $S$, $r$, $K$ and $\varepsilon$.
\end{theorem}
\begin{proof}\leanok
    \uses{thm:party_op_contraction, def:party_original_frame_sequence}
    The operator of an original circuit is a contraction by
    Theorem~\ref{thm:party_op_contraction}. A circuit without nonprivate gates has as many gate
    occurrences as nonprivate gates, none, and satisfies every bound vacuously, with error $0$.
    Under composition the numbers of nonprivate gates add and the bounds hold factor by factor;
    the errors satisfy
    $\lVert\operatorname{op}(e)(\operatorname{op}(c)x-y)\rVert\leq\lVert\operatorname{op}(c)x-y\rVert$
    and $\operatorname{op}(c)\operatorname{op}(e)x'=\operatorname{op}(c)x$.
\end{proof}

\begin{theorem}[Words and chains with one pair resource are bounded changes]
    \label{thm:party_resources_bounded}
    \lean{QuantumCircuit.Party.Word.isBoundedChange_toOriginalCircuit,
        QuantumCircuit.Party.Word.isBoundedChange_toOriginalCircuit_of_sourceCount_eq_zero,
        QuantumCircuit.Party.Monomial.isBoundedChange_toOriginalCircuit}
    \leanok
    \uses{def:party_bounded_change, def:party_word, def:party_pair_gates, def:party_allowed}
    Read an allowed word, or an allowed monomial of local contractions, normalized pair sources
    and normalized pair effects, as an original circuit whose local contractions are private and
    whose pair sources and pair effects are gates on the cells of their two parties, each with the
    single monomial ``this gate'' of coefficient one. Let $x$ and $y$ be vectors with
    $\lVert\operatorname{op}x-y\rVert\leq\varepsilon$ for the operator of the word or
    monomial.
    \begin{enumerate}
        \item An allowed word that uses only the parties of $S$ and has at most one pair source is
              a bounded change from $x$ to $y$ with participants in $S$, every $r$, every
              $K\geq1$ and error $\varepsilon$; these are the words of the births of
              Theorem~\ref{thm:party_birth}.
        \item An allowed word without pair sources is a bounded change from $x$ to $y$ for every
              $S$, $r$ and $K$, with error $\varepsilon$; these are the words of the exchanges of
              Theorem~\ref{thm:party_exchange}.
        \item An allowed monomial that uses only the parties of $S$ and has at most one gate is a
              bounded change from $x$ to $y$ with participants in $S$, every $r\geq1$, every
              $K\geq1$ and error $\varepsilon$; these are the chains of the deaths of
              Theorem~\ref{thm:party_birth}.
    \end{enumerate}
\end{theorem}
\begin{proof}\leanok
    \uses{def:party_original_circuit, def:party_word}
    The original circuit has one nonprivate gate for every pair source and pair effect, whose
    cell is the pair of its two parties, inside $S$ when the word or monomial uses only $S$. Its
    gate expansions have one monomial, of coefficient one, with at most one pair effect, so the
    expansion and monomial bounds hold for $r\geq1$ and $K\geq1$, and for all $r$ and $K$ when
    there is no gate. Its operator is that of the word or monomial.
\end{proof}

\begin{theorem}[Bounded changes of the stack of sheets]
    \label{thm:party_stack_bounded_change}
    \lean{QuantumCircuit.Party.hasSourceBoundedChange_of_perm,
        QuantumCircuit.Party.hasBoundedChange_of_perm,
        QuantumCircuit.Party.hasSourceBoundedChange_refl,
        QuantumCircuit.Party.hasBoundedChange_refl,
        QuantumCircuit.Party.HasSourceBoundedChange.perm_left,
        QuantumCircuit.Party.HasBoundedChange.perm_left,
        QuantumCircuit.Party.HasSourceBoundedChange.perm_right,
        QuantumCircuit.Party.HasBoundedChange.perm_right,
        QuantumCircuit.Party.isBoundedChange_liftOne,
        QuantumCircuit.Party.hasBoundedChange_liftOne,
        QuantumCircuit.Party.hasSourceBoundedChange_liftOne,
        QuantumCircuit.Party.isBoundedChange_liftTwo,
        QuantumCircuit.Party.hasBoundedChange_liftTwo,
        QuantumCircuit.Party.hasSourceBoundedChange_liftTwo,
        QuantumCircuit.Party.hasBoundedChange_activate,
        QuantumCircuit.Party.hasSourceBoundedChange_activate,
        QuantumCircuit.Party.hasBoundedChange_perm_liftOne,
        QuantumCircuit.Party.hasSourceBoundedChange_perm_liftOne}
    \leanok
    \uses{def:party_bounded_change, def:party_sheet_stack}
    Fix an idle system $W$ and $\xi\in\mathcal H_W$.
    \begin{enumerate}
        \item If $\mathbf F'$ is a permutation of $\mathbf F$, then $\mathbf F$, $\mathbf F'$ have
              a source-bounded change, hence a bounded change, for every $S$, $r$, $K$, $r_s$ and
              $\varepsilon\geq0$; in particular every list has the identity change. A
              (source-)bounded change from $\mathbf G$ to $\mathbf F'$ is one from every
              permutation $\mathbf F$ of $\mathbf G$, and one from $\mathbf F$ to $\mathbf G$ is one
              from $\mathbf F$ to every permutation $\mathbf F'$ of $\mathbf G$, with the same
              constants.
        \item Let $\lVert r_G\rVert\leq1$ for every sheet $G$ and $\lVert\xi\rVert\leq1$. If $c$ is
              a bounded change of $F$ into $F'$ from $r_F$ to $r_{F'}$, then $L_1(c)$ is a bounded
              change from $\omega_{(\mathbf G,F,\mathbf H)}(\xi)$ to
              $\omega_{(\mathbf G,F',\mathbf H)}(\xi)$ with the same participants, bounds and
              error, and a source-bounded one with constant $r_s$ if $c$ satisfies the pair-source
              bound $r_s$. The same holds for $L_2(c)$ when $c$ is a bounded change of $F_1,F_2$
              into $F_1',F_2'$ from $r_{F_1}\otimes r_{F_2}$ to $r_{F_1'}\otimes r_{F_2'}$.
        \item Under the norm conditions of the second item, if $\mathbf F$ is a permutation of
              $(\mathbf G,F_1,F_2,\mathbf H)$, for instance with an idle auxiliary sheet $F_2$ moved
              next to the main sheet $F_1$, a bounded change of the two adjacent sheets as in the
              second item gives a (source-)bounded change from $\mathbf F$ to
              $(\mathbf G,F_1',F_2',\mathbf H)$; likewise, if $\mathbf F$ is a permutation of
              $(\mathbf G,F,\mathbf H)$, a bounded change of $F$ gives one from $\mathbf F$ to
              $(\mathbf G,F',\mathbf H)$.
    \end{enumerate}
\end{theorem}
\begin{proof}\leanok
    \uses{thm:party_stack_perm, thm:party_bounded_change_basic, thm:party_stack_lift_error,
        def:party_sheet_stack, def:party_original_frame_sequence}
    A permutation of the sheets is the relabelling of Theorem~\ref{thm:party_stack_perm}, an
    original circuit without nonprivate gates, hence without pair sources, that maps one
    reference vector exactly to the other; Theorem~\ref{thm:party_bounded_change_basic} gives
    the first item. For the second, lifting keeps the number of nonprivate gates, the
    participants and the bounds by Definition~\ref{def:party_sheet_stack}, and
    Theorem~\ref{thm:party_stack_lift_error} bounds the error. The third item combines the second
    with the reordering of the first.
\end{proof}

\begin{theorem}[Private product initialization]
    \label{thm:party_stack_initialization}
    \lean{QuantumCircuit.Party.partyFactor, QuantumCircuit.Party.eq_of_apply_holder_eq,
        QuantumCircuit.Party.prod_partyFactor, QuantumCircuit.Party.norm_partyFactor,
        QuantumCircuit.Party.exists_productVector_eq_framesVec,
        QuantumCircuit.Party.Sheet.exists_productVector_eq_framesVec,
        QuantumCircuit.Party.exists_originalCircuit_framesVec,
        QuantumCircuit.Party.Sheet.exists_originalCircuit_framesVec}
    \leanok
    \uses{def:party_sheet_stack, def:party_sheet, def:party_original_circuit,
        def:siteconfig_product_vector}
    Let $h:\kappa\to\mathcal P$ assign a holder to every sheet, and let $\mathbf F$ be a list of
    sheets such that $X^F_p$ is a single point for every $F$ in $\mathbf F$ and every $p\neq h(F)$.
    \begin{enumerate}
        \item A configuration $z$ of $X^F$ is determined by $z_{h(F)}$. The factor of $F$ at the
              party $p$ is the vector $\pi_{F,p}$ on $X^F_p$ with
              $\pi_{F,p}(a)=\sum_{z:z_p=a}r_F(z)$ if $p=h(F)$ and $\pi_{F,p}=1$ otherwise; then
              $\prod_p\pi_{F,p}(z_p)=r_F(z)$, and $\lVert\pi_{F,p}\rVert$ is $\lVert r_F\rVert$ for
              $p=h(F)$ and $1$ otherwise.
        \item There are vectors $\varphi_p$ on $\mathrm{St}(\mathbf F;W)_p$, with $W$ the system of
              one basis vector at every party, such that
              $\lVert\varphi_p\rVert=\prod_{F\in\mathbf F,\,h(F)=p}\lVert r_F\rVert$ and the product
              vector $\bigotimes_p\varphi_p$ is $\omega_{\mathbf F}$.
        \item If moreover $\lVert r_F\rVert=1$ for every $F$ in $\mathbf F$, there is an original
              circuit without nonprivate gate from the system of one basis vector at every party
              to $\mathrm{St}(\mathbf F;W)$ whose operator maps that basis vector to
              $\omega_{\mathbf F}$.
        \item For a list of ownership sheets $S$ each owned entirely by $h(S)$, a unit vector
              $\Omega\in\C^{[q]^\Lambda}$ and $r_S=\Omega_S$, the second item holds with
              $\lVert\varphi_p\rVert=1$ for every party, and the third item holds.
    \end{enumerate}
    So the initial memory of \cite[Proposition~7.1]{OpenAI2026PolynomialPEPS}, every copy of
    $\Omega$ initialized privately at its future owner, is a pure product across the parties,
    prepared without nonprivate gates.
\end{theorem}
\begin{proof}\leanok
    \uses{def:party_sheet_stack, def:party_sheet, def:party_word,
        thm:party_resources_bounded}
    Two configurations of $X^F$ that agree at $h(F)$ agree everywhere, since the other bases are
    single points; so the sum defining $\pi_{F,h(F)}(a)$ has the single term $r_F(z)$ with
    $z_{h(F)}=a$, which gives the product formula and, the map $z\mapsto z_{h(F)}$ being a
    bijection, the norm. The second item follows by induction on the list: for $(F,\mathbf F')$
    take $\varphi_p=\pi_{F,p}\otimes\varphi'_p$, whose norm is
    $\lVert\pi_{F,p}\rVert\,\lVert\varphi'_p\rVert$ and whose product over the parties is
    $r_F\otimes\omega_{\mathbf F'}$. For the third, every $\varphi_p$ is a unit vector, and the
    word preparing $\varphi_p$ at every party by one local preparation is allowed and has no
    pair source, so as an original circuit it has no nonprivate gate, and it maps the basis
    vector to $\bigotimes_p\varphi_p$. For a sheet owned entirely by $h(S)$, every other party owns
    no site, so its basis $[q]^\emptyset$ is a single point, and
    $\lVert\Omega_S\rVert=\lVert\Omega\rVert=1$ by Definition~\ref{def:party_sheet}.
\end{proof}

